%!TEX root = /Users/kbj/Dropbox/*classes/601book/601master.tex
\chapter{Complement Structure}
\label{ch:compstructure}
\setcounter{exx}{0}
%\setcounter{page}{131}

\noindent
In this chapter, we'll take a closer look at the principles that organize verbs and their complements. We'll see that there is evidence that complements are arranged in a more structured way than our present system claims. Let's start by reviewing how our present grammar works.

\section{Accusatives Move} 
\label{sec:accmove}

We have a semantics that derives the fact that verbs must combine with the arguments they select. If we restrict our attention to the special case where the arguments of verbs are DPs (and maybe PPs), this is done with the assistance of $\theta$-roles. In chapter~\ref{ch:features}, we explored the idea that these $\theta$-roles are always separate morphological parts that are associated with a verb. For instance, the verb \emph{meet} has three morphological parts to it:%
%
\marginnote[0pt]{I have suppressed the c-selection information that comes with the meanings of these morphemes. Both Agent and Patient c-select~DPs.}%
%
\begin{exe}
	\raggedright
	\ex \label{ex5:meet1} 
	\small
	\begin{forest}
	[\obar{V}
		[$\surd$meet\\ $\theta$-role: \emph{agent}\\ $\theta$-role: \emph{patient} ]
	]
\end{forest}
\hspace*{8pt} $\Rightarrow$ \hspace*{8pt} \emph{meet}
	
\medskip
	
	\begin{xlist} 
		\ex \denotes{$\surd$meet} = $\lambda \epsilon$ \textsc{meet}($\epsilon$)
		\ex \denotes{agent} = $\lambda P$ $\lambda x$ $\lambda \epsilon$ $P(\epsilon)$ $\wedge$ \textsc{agent}($x$)($\epsilon$)
		\ex \denotes{patient} = $\lambda P$ $\lambda x$ $\lambda \epsilon$ $P(\epsilon)$ $\wedge$ \textsc{patient}($x$)($\epsilon$)
	\end{xlist}
\end{exe}
A D-structure that can house these morphemes and create a semantically well-formed phrase is in~(\ref{ex5:meet1t}).
\begin{exe}
	\raggedright
	\ex \label{ex5:meet1t} 
	\small
	\begin{forest}
	[vP
		[DP [Jesse, roof]]
		[\ibar{v} 
			[\obar{v} [agent] ]
			[VP [\ibar{V}
				[\obar{V}
					[\obar{V} [$\surd$meet] ]
					[\obar{v} [patient] ]
				 ]
				[DP [Robert, roof] ]
			] ]
		]
	]
	\end{forest}
\end{exe}
We derived the Theta Criterion from the requirement that vP must have a certain denotation---it must be a function whose domain is eventualities.
\begin{exe}
	\raggedright
	\ex \label{ex5:tco} Theta Criterion
	\begin{xlist}
		\ex Every XP that must be interpreted as an argument must be ``assigned'' a $\theta$-role.
		\ex Every $\theta$-role must be ``assigned'' to an argument.
	\end{xlist}
	\ex \label{ex5:vPden} A vP must have a denotation that makes it a function whose domain is eventualities and whose range is truth values.
\end{exe}
If the rules of semantic composition are just those in (\ref{ex5:semrules}), then the denotations we've assigned to the morphemes in (\ref{ex5:meet1t}) can be composed in a way that satisfies~(\ref{ex5:vPden}).
\begin{exe}
	\raggedright
	\ex \label{ex5:semrules} 
	\begin{xlist}
		\ex Function Application\\
		If $\gamma$ has daughters $\alpha$ and $\beta$, then \denotes{$\gamma$} is the result of applying the function \denotes{$\alpha$} to~\denotes{$\beta$}.
		\ex Trivial\\
		If $\gamma$ has only daughter $\alpha$, then \denotes{$\gamma$} = \denotes{$\alpha$}
	\end{xlist}
\end{exe}
As shown in (\ref{ex5:meettden}), for vP to have the denotation required of it by (\ref{ex5:vPden}) it is necessary for the vP to contain the arguments that the two $\theta$-roles require. If neither of those DP arguments were present, vP would be a function whose domain is things, not eventualities. Moreover, because the two DPs in (\ref{ex5:meettden}) can only semantically combine if they are arguments of some function---this is what distinguishes them from DPs like \emph{tomorrow}---they cannot be in a vP unless it contains $\theta$-roles that apply to them. In this way, then, the Theta Criterion is derived.
\begin{exe}
	\ex \label{ex5:meettden}
	\small
	\begin{forest}
	[$\lambda \epsilon$ \textsc{meet}($\epsilon$) $\wedge$ \textsc{patient}(\denotes{Robert})($\epsilon$) $\wedge$ \textsc{agent}(\denotes{Jesse})($\epsilon$)\\ vP
		[DP [Jesse, roof] ]
		[$\lambda x \lambda \epsilon$ \textsc{meet}($\epsilon$) $\wedge$ \textsc{patient}(\denotes{Robert})($\epsilon$) $\wedge$ \textsc{agent}($x$)($\epsilon$) \\\ibar{v} 
			[$\lambda P \lambda x \lambda \epsilon\ P(\epsilon)$ $\wedge$ \textsc{agent}($x$)($\epsilon$) \\\obar{v} [agent] ]
			[$\lambda \epsilon$ \textsc{meet}($\epsilon$) $\wedge$ \textsc{patient}(\denotes{Robert})($\epsilon$)\\VP [$\lambda \epsilon$ \textsc{meet}($\epsilon$) $\wedge$ \textsc{patient}(\denotes{Robert})($\epsilon$)\\ \ibar{V}
				[$\lambda y \lambda \epsilon$ \textsc{meet}($\epsilon$) $\wedge$ \textsc{patient}($y$)($\epsilon$) \\\obar{V} [meet] ]
				[DP [Robert, roof] ]
			]]
		]
	]
	\end{forest}
\end{exe}

To ensure that when non-arguments are present inside VP, they appear higher than the arguments, we have the Projection Principle.%
%
\marginnote[0pt]{This is stated as a condition on Lexical Insertion in chapter~\ref{ch:positargs}. Recall that an \underline{adopted daughter} of $\alpha$ is the daughter of $\alpha$ that $\alpha$ is not a projection of.}%
%
\begin{exe}
	\raggedright
	\ex \label{ex5:pp} Projection Principle\\
	The c-selection requirements of the morphemes in \obar{X} must be satisfied by every projection of \obar{X}. A phrase, $\pi$, satisfies a c-selection requirement of \obar{X} if that requirement hasn't already been satisfied and the adopted daughter of $\pi$ isn't different from what the c-selection requires. 
\end{exe}
This will correctly cause an argument to be closer to the verb than a non-argument linearly:
\begin{exe}
	\raggedright
	\ex \label{ex5:plin} 
	\begin{xlist}
		\ex[]{She talked to me yesterday.}
		\ex[??]{She talked yesterday to me.}
	\end{xlist}
\end{exe}
We also have an account of why when a verb has two internal arguments and they are PP and CP, there is a preference for the CP being last. CPs, we have discovered, move to the right edge of the clause.
\begin{exe}
	\raggedright
	\ex \label{ex5:ppcpo} 
	\begin{xlist}
		\ex[]{She explained to me that nattoo stinks.}
		\ex[??]{She explained that nattoo stinks to me.}
	\end{xlist}
\end{exe}

What we lost in chapter~\ref{ch:positargs}, however, was an explanation for the fact that when the two internal arguments of a verb are DP and PP, it's the DP that comes first.
\begin{exe}
	\raggedright
	\ex \label{ex5:dpppv} 
	\begin{xlist}
		\ex[]{She explained it to me.}
		\ex[*]{She explained to me it.}
	\end{xlist}
\end{exe}
Stowell's proposal was that the property of DPs relevant for this fact was their susceptibility to the Case filter. Their positions are determined by where Case can be licensed. If we could make the position that accusative Case is licensed by the position adjacent to the verb, then we will be able to understand the contrast in examples like (\ref{ex5:dpppv}). But we've now seen evidence that accusative Case is assigned by \obar{v}, not verbs, and this undoes Stowell's solution to the contrast in~(\ref{ex5:dpppv}). 

Or does it? If \obar{v} is the source of accusative Case, could we maintain Stowell's adjacency-based account of the placement of object DPs? Consider the D-structure our syntax assigns to~(\ref{ex5:dpppv}).
\begin{exe}
	\raggedright
	\ex \label{ex5:dpppvt} 
	\small
	\begin{forest}
	[TP[\ibar{T}
		[\obar{T} [tense: \emph{past}\\ case: \emph{nom}\\ $\bullet$$\phi$: ] ]
		[vP
			[DP\\ case:\\ $\phi$: \emph{\nth{3},sing} [she, roof] ]
			[\ibar{v}
				[\obar{v} [$\theta$-role: \emph{agent}\\ case: \emph{acc}] ]
				[VP [\ibar{V}
					[\ibar{V} 
						[\obar{V} [$\surd$explain\\ $\theta$-role:\\ $\phi$:\\ tense:] ]
						[DP\\ case: [it, roof] ]
					 ]
					[PP [to me, roof] ]
				] ]
			 ]
		 ]
	]]
	\end{forest}%
	%
	\marginnote[0pt]{I adopt here the convention of bundling the person and number features together into one feature: $\phi$.}%
	%
\end{exe}
There are unvalued Case features on the two DPs, there are unvalued $\phi$ features on \obar{T} and \obar{V}, and there are unvalued $\theta$-role and tense features on the verb. These can be valued by putting these features into the \textsc{agree} relation with the matching features in \obar{T} and \obar{v}. The $\bullet$ on $\phi$ will cause Argument Movement to bring the nominative Case marked DP into Specifier of TP. The derivation at this stage will look like~(\ref{ex5:dpppvt2}).
\begin{exe}
	\raggedright
	\ex \label{ex5:dpppvt2} 
	\small
	\begin{forest}
	[TP
		[DP\\ case: \emph{nom}\\ $\phi$: \emph{\nth{3},sing} [she, roof, name=binder] ]
		[\ibar{T}
			[\obar{T} [tense: \emph{past}\\ case: \emph{nom}\\  $\bullet$$\phi$: \emph{\nth{3},sing} ] ]
			[vP 
				[{\ }, name=trace ]
				[\ibar{v}
						[\obar{v} [$\theta$-role: \emph{agent}\\ case: \emph{acc}] ]
						[VP [\ibar{V}
							[\ibar{V}
								[\obar{V} [$\surd$explain\\ $\theta$-role:\\ $\phi$:\\ tense:] ]
								[DP\\ case: [it, roof] ]
							 ]
							[PP [to me, roof] ]
						]]
					 ]
			]
		 ]
	]
	\draw[->,>=stealth', ultra thin, gray] (trace) to[out=south west,in=south,looseness=2] node[pos=0.75,xshift=-10mm]{} (binder);
	\end{forest}
\end{exe}
Imagine that the $\theta$-role feature in \obar{v} has a $\bullet$ as well. That is, imagine that movement of the verb to \obar{v} occurs. This will bring \obar{V} within range of \obar{T}, so that its $\phi$ and tense features can be valued. In the resulting structure, the accusative Case marked DP will be string adjacent to the accusative Case assigner:~\obar{v}.
\begin{exe}
	\raggedright
	\ex \label{ex5:dpppvt3} 
	\small
	\begin{forest}
	[TP
		[DP\\ case: \emph{nom} [she, roof] ]
		[\ibar{T}
			[\obar{T} [tense: \emph{past}\\ $\bullet$case: \emph{nom}] ]
			[vP
				[{\ } ]
				[\ibar{V}
					[\obar{v}
						[\obar{V} [$\surd$explain\\ $\theta$-role: \emph{agent}, name=binder] ]
						[\obar{v} [$\bullet$$\theta$-role: \emph{agent}] ]
					 ]
					[VP [\ibar{V}
						[\ibar{V}
						[{\ }, name=trace ]
						[DP\\ case: \emph{acc} [it, roof]]
						]
						[PP [to me, roof]]
					] ]
				 ]
			 ]
		 ]
	]
	\draw[->,>=stealth', ultra thin, gray] (trace) to[out=south west,in=south,looseness=1] node[pos=0.75,xshift=-10mm]{} (binder);
	\end{forest}	
\end{exe}

There are some reasons to think that the verb in English moves in the way indicated in (\ref{ex5:dpppvt3}). Recall that Pollock's argument for movement of verbs in French comes in part from the observation that the verb can be separated from its internal argument by an adverb. This is possible in English as well; it's only blocked for accusative Case-marked objects.
\pagebreak
\begin{exe}
	\raggedright
	\ex \label{ex5:advasy} 
	\begin{xlist}
		\ex[]{She spoke carefully to me.}
		\ex[*]{She described carefully it.}
	\end{xlist}
\end{exe}
One way we have seen that internal arguments can be separated from their verbs is by way of a rightward movement, Complex NP Shift. But it seems unlikely that this is the method by which the PP in (\ref{ex5:advasy}a) manages to get separated from its verb, as \emph{to me} is neither necessarily focused nor ``heavy'' in the way that we have seen necessary for Complex NP Shift to arise. If verbs move, we can give (\ref{ex5:advasy}a) the syntax in~(\ref{ex5:advasyt}).
\begin{exe}
	\raggedright
	\ex \label{ex5:advasyt} 
	\small
	\begin{forest}
	[TP
		[DP\\ case: \emph{nom} [she, roof]]
		[\ibar{T} 
			[\obar{T} [tense: \emph{past}\\ case: \emph{nom}\\ $\bullet$$\phi$: \emph{\nth{3},sing}] ]
			[vP
				[{\ } ]
				[\ibar{v}
					[\obar{v}
						[\obar{V} [$\surd$speak\\ $\theta$-role: \emph{agent}, name=binder] ]
						[\obar{v} [$\bullet$$\theta$-role: \emph{agent}] ]
					 ]
					[VP [\ibar{V}
						[AdvP [carefully, roof] ]
						[\ibar{V}
							[{\ }, name=trace ]
							[PP [to me, roof] ]
						 ]
					] ]
				 ]
			 ]
		]
	]
	\draw[->,>=stealth', ultra thin, gray] (trace) to[out=south west,in=south,looseness=1] node[pos=0.75,xshift=-10mm]{} (binder);
	\end{forest}
\end{exe}
A parallel derivation is prevented when the internal argument of the verb is an accusative Case marked DP, as in (\ref{ex5:advasy}b), because the result would not cause that DP to be adjacent to the accusative Case assigner,~\obar{v}.

This account of (\ref{ex5:advasy}) carries over to the ungrammaticality of (\ref{ex5:shift1}a), since adjacency between the accusative Case marked DP and \obar{v} doesn't arise here either, but it doesn't explain the grammaticality of~(\ref{ex5:shift1}b).
\begin{exe}
	\raggedright
	\ex \label{ex5:shift1} 
	\begin{xlist}
		\ex[*]{She explained carefully it to me.}
		\ex[]{She explained it carefully to me.}
	\end{xlist}
\end{exe}
I\cite{Johnson:1991} argued that this word order involves Argument Movement of the accusative Case-marked DP into Specifier of VP, as in~(\ref{ex5:shift2}).
\begin{exe}
	\raggedright
	\ex \label{ex5:shift2} 
	\small
	\begin{forest}
	[TP
		[DP\\ case: \emph{nom} [she, roof]]
		[\ibar{T} 
			[\obar{T} [tense: \emph{past}\\ case: \emph{nom}\\ $\bullet$$\phi$: \emph{\nth{3}, sing}]]
			[vP
				[{\ } ]
				[\ibar{v}
					[\obar{v}
						[\obar{V} [$\surd$speak\\ $\theta$-role: \emph{agent}] ]
						[\obar{v} [$\bullet$$\theta$-role: \emph{agent}] ]
					 ]
					[VP
						[DP\\ case: \emph{acc} [it, name=binder, roof] ]
						[\ibar{V}
							[AdvP [carefully, roof] ]
							[\ibar{V}
								[\ibar{V}
									[{\ } ]
									[{\ }, name=trace]
								]
								[PP [to me, roof] ]
							 ]
						 ]
					 ]
				 ]
			 ]
		]
	]
	\draw[->,>=stealth', ultra thin, gray] (trace) to[out=south west,in=south,looseness=1] node[pos=0.75,xshift=-10mm]{} (binder);
	\end{forest}
\end{exe}

My argument hinged on the behavior of a class of verbs which seem to have two parts. These are called ``phrasal verbs'' by Dwight Bolinger\cite{Bolinger:1971} and sometimes go by the name ``particle verbs.'' Some of these verbs have a meaning that is not compositional; that is, their two parts do not have identifiable denotations whose meanings combine by a productive rule of semantic composition. Two examples are \emph{look up} and \emph{throw up}.
\begin{exe}
	\raggedright
	\ex \label{ex5:lookup} 
	She looked it up.\\ She threw it up.
\end{exe}
The eventuality described by \emph{look up} does not include the denotation associated with \emph{up}, nor are the eventualities described by \emph{throw up} a special kind of those eventualities described by \emph{throw}. \emph{look up} and \emph{throw up} have the idiosyncratic meanings that individual words typically have. These expressions also undergo morphological processes which target words, rather than phrases. For instance, there is a morphological process that makes nouns from verbs by suffixing \emph{ing}. The nouns in (\ref{ex5:ing1a}) are produced by this rule.
\begin{exe}
	\raggedright
	\ex \label{ex5:ing1a} 
	\begin{xlist}
		\ex the eating of cake
		\ex the drying of the towels
		\ex Shai's talking about Hebrew paradigms
	\end{xlist}
\end{exe}
This process is capable of making nouns out of particle verbs: \enlargethispage*{14pt}
\begin{exe}
	\raggedright
	\ex \label{ex5:ing2a} 
	\begin{xlist}
		\ex the throwing up of cake
		\ex the looking up of telephone numbers
		\ex Shai's shouting out at Kiyomi
	\end{xlist}
\end{exe}
It's clear that to whatever extent these expressions are single verbs, they occupy different syntactic positions in (\ref{ex5:lookup}). The cognate expressions in the verb final Germanic languages are separated by Head Movement. The verb and particle are adjacent, and usually written as one word, when Head Movement has not relocated \obar{V} to second position. But when \obar{V} is moved into \obar{C} position in verb second contexts, the verb and its particle are separated. This is illustrated by~(\ref{ex5:seppref}).
\begin{exe}
	\raggedright
	\ex \label{ex5:seppref} 
	\begin{xlist}
		\ex[]{\gll Sie sagte dass Ich den Brief ausstellte. \\ She said that I the letter out.stood  \\ `She said that I wrote the leter out.'}
		\ex[]{\gll Ich stellte den Brief aus. \\ I stood the letter out \\ `I wrote the letter out.'}
		\ex[*]{\gll Ich ausstellte den Brief. \\ I out.stood the letter \\ `I wrote the letter out.'}
	\end{xlist}
\end{exe}
I suggested that Head Movement is also responsible for separating the verbal part from its particle in English as well. I suggested a D-structure representation that allowed Head Movement to target the verbal portion of the particle verb, as in~(\ref{ex5:part1}).
\begin{exe}
	\raggedright
	\ex \label{ex5:part1} 
	\small
	\begin{forest}
	[vP
		[DP [she, roof]]
		[\ibar{v} 
			[\obar{v}
				[\obar{V} [$\surd$throw\\ $\theta$-role: \emph{agent}, name=binder] ]
				[\obar{v} [$\bullet$$\theta$-role: \emph{agent}] ]
			 ]
			[VP
				[DP [it, roof] ]
				[\ibar{V}
					[\obar{V}
						[{\ }, name=trace ]
						[up ]
					 ]
					[{\ } ]
				 ]
			 ]
		]
	]
	\draw[->,>=stealth', ultra thin, gray] (trace) to[out=south west,in=south,looseness=1] node[pos=0.75,xshift=-10mm]{} (binder);
	\end{forest}
\end{exe}

An interesting feature of constructions with particle verbs is that though they c-select internal arguments of a variety of kinds, it is only DPs that can be positioned between the verbal part and its particle.
\begin{exe}
	\raggedright
	\ex \label{ex5:dppart} 
	\begin{xlist}
		\ex[*]{Claire shouted [\PP to me] out.}
		\ex[*]{Aleksei pointed [\CP that natto stinks] out.}
		\ex[*]{Nick turned [\AP good] out.}
		\sn \emph{compare:}
		\ex[]{Claire shouted out [\PP to me].}
		\ex[]{Aleksei pointed out [\CP that natto stinks].}
		\ex[]{Nick tuned out [\AP good].}
	\end{xlist}
\end{exe}
Richard Kayne\cite{Kayne:1984b} argued that this fact can be explained if the particle heads a small clause that the verbal part of the particle verb selects. Small clauses can only have DPs in their Specifier positions, as (\ref{ex5:smdp}) indicates.
\begin{exe}
	\raggedright
	\ex \label{ex5:smdp} 
	\begin{xlist}
		\ex[]{She considers [[\DP the woman] happy].}
		\ex[*]{She considers [[\PP to me] obvious that natto stinks].}
		\ex[*]{She considers [[\CP that natto stinks] obvious to me].}
	\end{xlist}
\end{exe}
So whatever ensures this could be relied on to derive the same asymmetry with particle verbs. On his account, then, the representation for particle verbs would be (\ref{ex5:kaynep}).
\begin{exe}
	\raggedright
	\ex \label{ex5:kaynep} 
	\small
	\begin{forest}
	[vP
		[DP [Claire, roof]]
		[\ibar{v}
			[\obar{v} 
				[\obar{V} [$\surd$point\\ $\theta$-role: \emph{agent}, name=binder] ]
				[\obar{v} [$\bullet$$\theta$-role: \emph{agent}\\ case: \emph{acc}] ]
			 ]
			[VP [\ibar{V}
				[{\ }, name=trace ]
				[PP
					[DP [it, roof] ]
					[\ibar{P} [\obar{P} [out]] ]
				 ]
			] ]
		 ]
	]
	\draw[->,>=stealth', ultra thin, gray] (trace) to[out=-110,in=-80,looseness=1] node[pos=0.75,xshift=-10mm]{} (binder);
	\end{forest}
\end{exe}

Unlike (\ref{ex5:part1}), Kayne's representation in (\ref{ex5:kaynep}) correctly explains why the verb and particle cannot Head Move together. If that were possible, we would expect examples like \emph{She pointed out it} to be grammatical. On the view in \citet{Johnson:1991}, there would have to be another reason why the verb+particle unit cannot move by Head Movement.

On the view of particle verbs that involves movement of the accusative Case-marked DP into Specifier of VP, the fact that only DPs can come between the verbal part and its particle can be credited to Earliness, repeated here.%
%
\marginnote[0pt]{Recall that a sentence is the set of grammatical derivations which have the same D-structure.}%
%
\begin{exe}
	\raggedright
	\ex \label{ex5:earliness} Earliness\\
	Let the set of derivations, $\mathscr{D} = \{ D, D_1,\ldots,D_n \}$, be a sentence. There can be no $D_i \in \mathscr{D}$ if there is a $D_j$ in $\mathscr{D}$ such that $|D_i|\ >\ |D_j|$, where ``$|D|$'' is the number of elements in~$D$.
\end{exe}
If argument DPs are uniquely subject to the Case filter, and if we could manufacture a theory that let Specifier of VP be the only position that is assigned accusative Case, then movement of an argument DP into Specifier of VP will be required to meet the Case filter. Derivations that involve movement of an argument DP into Specifier of VP will therefore be permitted by Earliness---a derivation that didn't include that movement wouldn't be in the set of grammatical derivations that make a sentence -- but movement of other phrases would be blocked. Any derivation that includes the step in (\ref{ex5:needless1}), for instance, will be longer than an equivalent, and equally grammatical derivation, that didn't include that step.
\begin{exe}
	\raggedright
	\ex \label{ex5:needless1} 
	\small
	\begin{forest}
	[\ibar{v}
		[\obar{v} 
			[\obar{V} [$\surd$shout] ]
			[\obar{v} ]
		]
		[VP 
			[PP [to me, name=binder, roof]]
			[\ibar{V} 
				[\obar{V} 
					[{\ } ]
					[out ]
				]
				[{\ }, name=trace ]
			 ]
		]
	]
	\draw[->,>=stealth', ultra thin, gray] (trace) to[out=south west,in=south,looseness=2] node[pos=0.75,xshift=-10mm]{} (binder);
	\end{forest}
\end{exe}
But if accusative Case is only assigned to Specifier of VP, then we don't have an understanding of how the word orders in~(\ref{ex5:prtDP1}) arise.
\begin{exe}
	\raggedright
	\ex \label{ex5:prtDP1} 
	\begin{xlist}
		\ex She pointed out the mistake to us.
		\ex She threw up the burrito.
	\end{xlist}
\end{exe}

I'll sketch an account of these facts that closely follows \citet{Johnson:1991}, but makes use of the framework we've developed in the last chapter for Case assignment. That system allows Case assignment---that is, \textsc{agree}---to hold between a head and a phrase that head c-commands across long distances. The locality condition we posited for \textsc{agree} merely prevented two terms from being in the \textsc{agree} relation if there was a third term between them that could be in that same \textsc{agree} relation. If \obar{v} is the source of accusative Case, then this theory should allow an accusative Case feature in \obar{v} to value an unvalued Case feature on an object DP in~(\ref{ex5:long}).
\begin{exe}
	\raggedright
	\ex \label{ex5:long} 
	\small
	\begin{forest}
	[\ibar{v}
		[\obar{v} 
			[$\theta$-role: \emph{agent} ]
			[case: \emph{acc} ]
		]
		[VP [\ibar{V}
			[\ibar{V}
				[\obar{V} 
					[\obar{V} [$\surd$point\\ $\theta$-role:] ]
					[\obar{P} [out] ]
				 ]
				[DP\\ case: [the problem, roof] ]
			 ]
			[PP [to us, roof] ]
		]]
	]
	\end{forest}
\end{exe}
To force the accusative Case-marked DP to move, our present system equips one of the terms in the \textsc{agree} relation with a $\bullet$. If we associated a $\bullet$ with the Case feature in \obar{v}, this would produce a parse like~(\ref{ex5:long2}).
\begin{exe}
	\raggedright
	\ex \label{ex5:long2} 
	\small
	\begin{forest}
	[vP
		[DP\\ case: \emph{acc} [the problem, name=binder, roof] ]
		[\ibar{v} 
			[\obar{v} [$\theta$-role: \emph{agent}\\ $\bullet$case: \emph{acc}] ]
			[VP [\ibar{V}
				[\obar{V}
					[\obar{V} [$\surd$point\\ $\theta$-role:] ]
					[\obar{P} [out] ]
				 ]
				[{\ }, name=trace ]
			] ]
		]
	]
	\draw[->,>=stealth', ultra thin, gray] (trace) to[out=south west,in=south,looseness=1] node[pos=0.75,xshift=-10mm]{} (binder);
	\end{forest}
\end{exe}
This would correctly put the object DP first among internal arguments, but it will require that the verb in English move even farther than \obar{v}, for otherwise accusative Case-marked objects would erroneously be forced to precede the verb. Let's call the mysterious position to which the verb would have to move on this account~\obar{X}.
\begin{exe}
	\raggedright
	\ex \label{ex5:long3} 
	\small
	\begin{forest}
	[\ibar{X}
		[\obar{X}
			[\obar{v}, name=binder 
				[\obar{V} [$\surd$point] ]
				[\obar{v} [$\theta$-role: agent\\ $\bullet$case: \emph{acc}] ]
			 ]
			[\obar{X} ]
		 ]
		[vP
			[DP [the problem, roof] ]
			[vP
				[DP [she, roof] ]
				[\ibar{v}
					[{\ }, name=tracev ]
					[VP [\ibar{V}
							[\ibar{V} 
								[\obar{V}
									[{\ }, name=trace]
									[\obar{P} [out]]
								 ]
								[PP [to us, roof] ]
							 ]
							[{\ } ]
					 ]]
				 ]
			 ]
		 ]
	]
	\draw[->,>=stealth', ultra thin, gray] (trace) to[out=west,in=-110,looseness=1] node[pos=0.75,xshift=-10mm]{} (tracev);
	\draw[->,>=stealth', ultra thin, gray] (tracev) to[out=south west,in=-80,looseness=2] node[pos=0.75,xshift=-10mm]{} (binder);
	\end{forest}
\end{exe}
This is a popular approach, but I will present problems for its semantics later in this chapter.

A slight variant of this treatment adopts the view that the two pieces of the \obar{v} position that we have adopted to account for Burzio's Generalization can be syntactically separated in much the same way that French separates tense and agreement. We'll continue to adopt the solution to Burzio's Generalization that credits the external $\theta$-role with being the type of word that comes with an accusative Case feature. But rather than let the D-structure assemble these two pieces into a single \obar{X} position, we will assume that the derivation brings them together. At D-structure, let's put the accusative feature in a lower \obar{X} position; \citet{Johnson:1991} calls this position ``$\mu$.''%
%
\marginnote[0pt]{I took the name from a paper that was circulating at the time by David Pesetsky. Pesetsky's paper had many of the same ideas as mine---we had independently arrived at very similar proposals. My paper got published first, robbing Professor Pesetsky of the credit he deserves. More common names for this head position found in the literature are \obar{Acc} for ``accusative,'' and \obar{Agr_o}, for ``object agreement.''}%
%
\begin{exe}
	\raggedright
	\ex \label{ex5:long4} 
	\small
	\begin{forest}
	[vP
		[DP [she, roof]]
		[\ibar{v} 
			[\obar{v} [$\bullet$$\theta$-role: \emph{agent}\\ case: \emph{acc}] ]
			[$\mu$P [\ibar{$\mu$}
				[\obar{$\mu$} [ $\bullet$case:] ]
				[VP [\ibar{V}
					[\ibar{V}
						[\obar{V}
							[\obar{V} [$\surd$point\\ $\theta$-role:] ]
							[\obar{P} [out] ]
						 ]
						[PP [to us, roof] ]
					 ]
					[DP\\ case: [the problem, roof] ]
				] ]
			]]
		]
	]
	\end{forest}
\end{exe}
This view is sometimes called the ``Split VP Hypothesis,'' and it is given an extensive defense in \citet{Koizumi:1995}. From (\ref{ex5:long4}), repeated applications of \textsc{agree} can value the unvalued Case features on \obar{$\mu$} and the DP, and because of the $\bullet$s, this will trigger a movement of \obar{$\mu$} and DP. The result is~(\ref{ex5:long4b}).
\begin{exe}
	\raggedright
	\ex \label{ex5:long4b} 
	\small
	\begin{forest}
	[vP
		[DP [she, roof]]
		[\ibar{v} 
			[\obar{v}
				[\obar{$\mu$} [$\bullet$case: \emph{acc}, name=binderm] ]
				[\obar{v} [$\bullet$$\theta$-role: \emph{agent}] ]
			 ]
			[$\mu$P
				[DP\\ case: \emph{acc} [the problems, roof, name=binder]]
				[\ibar{$\mu$}
					[{\ }, name=tracem ]
					[VP [\ibar{V}
						[\ibar{V}
							[\obar{V}
								[\obar{V} [$\surd$point\\ $\theta$-role:] ]
								[\obar{P} [out] ]
							 ]
							[PP [to us, roof]]
						 ]
						[{\ }, name=trace ]
					] ]
				]
			 ]
		]
	]
	\draw[->,>=stealth', ultra thin, gray] (trace) to[out=-110,in=-80,looseness=1.3] node[pos=0.75,xshift=-10mm]{} (binder);
	\draw[->,>=stealth', ultra thin, gray] (tracem) to[out=south west,in=-80,looseness=1.5] node[pos=0.75,xshift=-10mm]{} (binderm);
	\end{forest}
	\marginnote[0pt]{The proposal that accusative Case-marked DPs move in English interacts badly with the model of quantifier float that we reviewed in section~\ref{sec:qfloat}. We would expect ungrammatical sentences like \emph{Meg pointed the problems out all to me} to be well-formed. See \citet{Bobaljik:2003} for some discussion of this problem and ideas about its solution.}
\end{exe}
This representation allows \textsc{agree} to value the $\theta$-role feature on the verb, and the $\bullet$ requires the movement shown in~(\ref{ex5:long4c}).
\begin{exe}
	\raggedright
	\ex \label{ex5:long4c} 
	\small
	\begin{forest}
	[vP
		[DP [she, roof]]
		[\ibar{v} 
			[\obar{v}
				[\obar{V} [$\surd$point\\ $\theta$-role: \emph{agent}, name=binder] ]
				[\obar{v}
					[\obar{$\mu$} [$\bullet$case: \emph{acc}] ]
					[\obar{v} [$\theta$-role: \emph{agent}\\ case: \emph{acc}] ]
				 ]
			 ]
			[$\mu$P
				[DP\\ case: \emph{acc} [the problems, roof] ]
				[\ibar{$\mu$} 
					[{\ } ]
					[VP [\ibar{V}
						[\ibar{V} 
							[\obar{V}
								[{\ }, name=trace ]
								[\obar{P} [out] ]
							 ]
							[PP [to us, roof] ]
						]
						[{\ } ]
					] ]
				 ]
			 ]
		]
	]
	\draw[->,>=stealth', ultra thin, gray] (trace) to[out=south west,in=south,looseness=1] node[pos=0.75,xshift=-10mm]{} (binder);
	\end{forest}
\end{exe}
The $\bullet$ on the $\theta$-role feature can also be satisfied with the movement in~(\ref{ex5:long4d}).
\begin{exe}
	\raggedright
	\ex \label{ex5:long4d} 
	\small
	\begin{forest}
	[vP
		[DP [she, roof]]
		[\ibar{v} 
			[\obar{v}
				[\obar{V} [$\surd$point out, name=binder, roof] ]
				[\obar{v} 
					[\obar{$\mu$} [$\bullet$case: \emph{acc}] ]
					[\obar{v} [$\bullet$$\theta$-role: \emph{agent}\\ case: \emph{acc}] ]
				 ]
			 ]
			[ $\mu$P
				[DP\\ case: \emph{acc} [the problems, roof] ]
				[\ibar{$\mu$}
					[{\ } ]
					[VP [\ibar{V}
						[\ibar{V}
							[{\ }, name=trace ]
							[PP [to us, roof] ]
						 ]
						[{\ } ]
					] ]
				 ]
			]
		]
	]
	\draw[->,>=stealth', ultra thin, gray] (trace) to[out=south west,in=south,looseness=1] node[pos=0.75,xshift=-10mm]{} (binder);
	\end{forest}
\end{exe}
Note that because both of these instances of Head Movement satisfy the same requirement of the $\bullet$, Earliness does not distinguish between them. Both outcomes contribute to the derivations that contain them an equal number of elements, and therefore both word orders are available. 

The key to meeting the requirements of Earliness is the view that the verb and its particle make an \obar{X} that can move. In the cognate constructions in the other Germanic languages, it is clear that Head movement of the verb and particle together is sometimes blocked. If the verb and particle were able to Head Move together to \obar{C}, for instance, sentences like (\ref{ex5:seppref}c), where the verb+particle occupy second position, would be grammatical in German. 

I suggested in \citet{Johnson:1991} that this might be connected to the fact that morphological processes generally resist taking the verb+particle as their host. For instance, the process that makes nouns from verbs by suffixing \emph{ing} (illustrated by (\ref{ex5:ing1})) targets only the verbal part of the particle verb, as (\ref{ex5:ing2}) indicates.%
%
\marginnote[0pt]{The possible exception to this trend is \emph{able}, where English speakers consistently judge \emph{pointoutable} to be at least as grammatical as \emph{pointableout}.}%
%
\begin{exe}
	\raggedright
	\ex \label{ex5:ing1} \small
	{}[_{\text{\obar{V}}} read] (e.g., \emph{read the book}) $\Rightarrow$ [_{\text{\obar{N}}} reading] (e.g. \emph{the reading of the book})
	\ex \label{ex5:ing2}
	\begin{xlist}
		\sn {}[_{\text{\obar{V}}}point out] $\Rightarrow$ [_{\text{\obar{N}}} point outing]
		\sn {}[_{\text{\obar{V}}}point out] $\Rightarrow$ [_{\text{\obar{N}}} pointing out]
	\end{xlist}	
\end{exe}
Imagine that this force influences how the morphemes in an \obar{X} that expresses verbal morphology can be arranged. For instance, consider how a particle verb that is embedded under the inflectional head that makes a perfect participle might deploy itself. There will be a parse in the derivation for the intransitive particle verb \emph{turn out} that will look like~(\ref{ex5:inflpv1}).
\begin{exe}
	\raggedright
	\ex \label{ex5:inflpv1} 
	\small
	\begin{forest}
	[\ibar{AsP}
		[\obar{Asp} [$\bullet$aspect: \emph{perf}]]
		[vP 
			[DP [they, roof] ]
			[\ibar{v}
				[\obar{v}
					[\obar{V}, name=binder
						[\obar{V} [$\surd$turn\\ $\theta$-role:\\ aspect: \emph{perf}] ]
						[\obar{P} [out]]
					 ]
					[\obar{v} [$\bullet$$\theta$-role: \emph{agent}] ]
				 ]
				[VP [\ibar{V} [{\ }, name=trace]]]
			 ]
		]
	]
	\draw[->,>=stealth', ultra thin, gray] (trace) to[out=south west,in=south,looseness=1] node[pos=0.75,xshift=-10mm]{} (binder);
	\end{forest}
\end{exe}
If Head Movement were to move the verb+particle pair to \obar{Asp} in (\ref{ex5:inflpv1}), the result would put the aspect morphology outside the particle syntactically.
\begin{exe}
	\raggedright
	\ex \label{ex5:inflpv2} 
	\small
	\begin{forest}
	[\ibar{Asp}
		[\obar{Asp} 
			[\obar{V}, name=binder
				[\obar{V} [$\surd$turn\\ $\theta$-role: \emph{agent}\\ aspect: \emph{perf}] ]
				[\obar{P} [out] ]
			 ]
			[\obar{Asp} [$\bullet$aspect: \emph{perfect}] ]
		]
		[vP 
			[DP [they, roof] ]
			[\ibar{v}
				[\obar{v}
					[{\ }, name=trace ]
					[\obar{v} [$\bullet$$\theta$-role: \emph{agent}] ]
				 ]
				[VP [\ibar{V} [{\ }]] ]
			 ]
		]
	]
	\draw[->,>=stealth', ultra thin, gray] (trace) to[out=south west,in=-80,looseness=1] node[pos=0.75,xshift=-10mm]{} (binder);
	\end{forest}
\end{exe}
If the syntactic structure mirrors the morphological structure, this plausibly violates whatever is responsible for the contrast in (\ref{ex5:ing2}). If, by contrast, Head Movement combines the verb and the aspectual morphology by moving just the verbal portion, the resulting syntax aligns with what (\ref{ex5:ing2}) teaches is the right morphological structure.\marginnote[0pt]{If this representation is grammatical, there will have to be a way of matching a vocabulary item with the \obar{v} containing the particle and the agent $\theta$-role. And, relatedly, the $\bullet$ will have to be satisfied by a representation in which the verbal part of the particle verb has been separated. Perhaps the agent $\theta$-role feature should be associated with the \obar{V} that contains both the verbal part and its particle, rather than, as indicated in the representation here, with just the verbal part.}
\begin{exe}
	\raggedright
	\ex \label{ex5:inflpv3} 
	\small
	\begin{forest}
	[\ibar{Asp}
		[\obar{Asp} 
			[\obar{V} [$\surd$turn\\ $\theta$-role: \emph{agent}\\ aspect: \emph{perf}, name=binder] ]
			[\obar{Asp} [$\bullet$aspect: \emph{perf}] ]
		]
		[vP 
			[DP [they, roof]]
			[\ibar{v}
				[\obar{v}
					[\obar{V}
						[{\ }, name=trace ]
						[\obar{P} [out] ]
					 ]
					[\obar{v} [$\bullet$$\theta$-role: \emph{agent}] ]
				 ]
				[VP [\ibar{V} []] ]
			 ]
		]
	]
	\draw[->,>=stealth', ultra thin, gray] (trace) to[out=-110,in=-80,looseness=1] node[pos=0.75,xshift=-10mm]{} (binder);
	\end{forest}
\end{exe}

It is difficult to establish that this same kind of restriction on the morphology of particle verbs holds in German because the particle in German is a ``prefix.'' But if there is a similar restriction, then this could be credited with ensuring that the verb and particle separate when the verb gains its finite morphology. This will account for why the verb+particle does not move as a unit to \obar{C} position in verb second contexts.

\section{Split VP} 
\label{sec:split}

In chapter \ref{ch:positargs} we saw reasons for thinking that the meaning one might na\"ively have associated with one position---the position the verb occupies---is actually spread across two syntactic positions. In chapter \ref{ch:features} we developed a theory of how the syntax can assemble the morphological pieces of a word, and imported that theory to the task of assembling the two pieces of a verb's meaning. In this chapter, we've see that accusative Case assigning function of verbs is part of the morpheme that expresses the external $\theta$-role, but it too is syntactically separated from that $\theta$-role. Over the course of these chapters, what we've discovered is that a variety of attributes we associate with verbs---their external $\theta$-role and their structural Case---are spread across different syntactic positions. What appeared to be a single projection of a verb---its VP---has turned into a variety of phrases, one embedded in the other. 

There are other aspects of the meaning associated with a verb that have been argued to involve splitting the VP up further. We will look at a couple of these here, using the view of constituency that we get from the meaning of~\emph{again}.

\subsection{Hidden states and ``again''} 
\label{sub:again}

An interesting feature of \emph{again} is that its surface syntactic distribution is arguably determined entirely by its semantics. The position of many other modifiers appears to be at least partly determined by the morpho-syntactic category of the phrases they modify (adjectives modify projections of nouns and not projections of verbs, for instance) and partly by way of movement, allowing them to be spoken in positions distant from where they are semantically interpreted, as in: \emph{Quickly, Trump should leave the news-cycle.}

To see this, consider the meanings of the sentences in~(\ref{ex5:againbase1}), which are fashioned from a paper by James McCawley.\cite{McCawley:1973a}
\begin{exe}
	\raggedright
	\ex \label{ex5:againbase1} 
	\begin{xlist}
		\ex \label{ex5:againbase1a} Again Della didn't kick it.
		\sn $\approx$ Della didn't kick it and she didn't kick it previously.
		\ex Della didn't kick it again.
		\sn $\approx$ Della didn't kick it and it hadn't been kicked previously.
	\end{xlist}
\end{exe}
In (\ref{ex5:againbase1}a), \emph{not} is part of what \emph{again} semantically combines with, and in (\ref{ex5:againbase1}b), by contrast, \emph{again} semantically combines with just \emph{kick it}. If \emph{again} cannot move, this ambiguity follows straightforwardly from the thesis of semantic compositionality, and the plausible parses for (\ref{ex5:againbase1}) in~(\ref{ex5:againbase2}).
\begin{exe}
	\raggedright
	\ex \label{ex5:againbase2} 
	\small
\begin{forest}
[TP
	[AdvP [again, roof]]
	[TP 
		[DP [Della, roof]]
		[\ibar{T}
			[\obar{T} [didn't] ]
			[vP [kick it, roof] ]
		 ]
	]
]
\end{forest}
\hspace*{10pt}
\begin{forest}
[TP
	[DP [Della, roof, name=binder]]
	[\ibar{T} 
		[\obar{T} [didn't] ]
		[vP
			[{\ }, name=trace ]
			[\ibar{v}
				[\obar{v} ]
				[VP [\ibar{V}
					[\ibar{V} 
						[\obar{V} [kick] ]
						[DP [it, roof] ]
					 ]
					[AdvP [again, roof] ]
				] ]
			 ]
		 ]
	]
]
\draw[->,>=stealth', ultra thin, gray] (trace) to[out=south west,in=south,looseness=1] node[pos=0.75,xshift=-10mm]{} (binder);
\end{forest}
	\marginnote[0pt]{Note that the first of these parses requires that we allow D-structures to adjoin a phrase to a maximal projection. That is a form that, so far, we have only seen the output of Complex NP Shift to have.}
\end{exe}

I'll provide a sketch of a semantics for \emph{again} here that is based on proposals by Arnim von Stechow.\cite{vonStechow:1996} Roughly, the idea is that \emph{again} combines with phrases whose denotations are descriptions of eventualities, and adds that those descriptions held of a previously occurring eventuality.
\begin{exe}
	\raggedright
	\ex \label{ex5:againden} 
	\denotes{again} = $\lambda P \lambda \epsilon\ P(\epsilon)$ $\wedge$  $\exists \epsilon '$ [$\epsilon ' < \epsilon \wedge\ P(\epsilon ') ]$\\
	``$<$'' means ``occurs previously to''
\end{exe} 
This will give to the \ibar{V}s in the second parse of (\ref{ex5:againbase2}) the meanings indicated in~(\ref{ex5:base3}).
\begin{fullwidth}
\begin{exe}
	\medskip
	\raggedright
	\ex \label{ex5:base3} 
	\hspace*{-20pt}
	\small
	\begin{forest}
	[$\lambda \epsilon$ \textsc{kick}($\epsilon$) $\wedge$ \textsc{patient}(it)($\epsilon$) $\wedge$ $\exists \epsilon'$ {[}$\epsilon' < \epsilon$ $\wedge$ \textsc{kick}($\epsilon'$) $\wedge$ \textsc{patient}(it)($\epsilon'$){]}\\ \ibar{V}
		[$\lambda \epsilon$ \textsc{kick}($\epsilon$) $\wedge$ \textsc{patient}(it)($\epsilon$)\\ \ibar{v} 
			[$\lambda x \lambda \epsilon$ \textsc{kick})($\epsilon$) $\wedge$ \textsc{patient}($x$)($\epsilon$)\\ \obar{V} [kick] ]
			[DP [it, roof] ]
		]
		[$\lambda P \lambda \epsilon P(\epsilon)$ $\wedge$ $\exists e' {[}e'<e$ $\wedge$ $P(\epsilon')${]}\\ AdvP [again, roof] ]
	]]
	\end{forest}
\end{exe}
\end{fullwidth}
The denotation of the largest \ibar{V} makes it a description of eventualities that map to True when they are kickings whose patient is whatever \emph{it} refers to, and there is another eventuality that held previously which also maps to True when it is a kicking whose patient is whatever \emph{it} refers to. This matches the paraphrases given above, then.

Although this denotation captures the rough meaning of \emph{again}, it misfires once we consider how it is put into the rest of (\ref{ex5:againbase1}b). Because \emph{nt} corresponds to the logical operator $\neg$, the denotation we'll get from (\ref{ex5:base3}) for (\ref{ex5:againbase1}b)---ignoring the contribution made by tense---is roughly~(\ref{ex5:base4}).
\begin{exe}
	\raggedright
	\ex \label{ex5:base4} 
	$\neg \exists \epsilon$ \parbox[t]{3in}{\textsc{agent}(Della)($\epsilon$) $\wedge$ \textsc{kick}($\epsilon$) $\wedge$ \textsc{patient}(it)($\epsilon$) $\wedge$\\
	$\exists \epsilon '$ [$\epsilon ' < \epsilon$ $\wedge$ \textsc{kick}($\epsilon '$) $\wedge$ \textsc{patient}(it)($\epsilon '$)]}
\end{exe}
A paraphrase of (\ref{ex5:base4}) is~(\ref{ex5:base4p}).
\begin{exe}
	\raggedright
	\ex \label{ex5:base4p} 
	There is no eventuality that has Della as its agent and is a kicking of it and has a prior eventuality that is a kicking of it.
\end{exe}
This wrongly puts inside the scope of $\neg$ that there is a prior eventuality of kicking it. The meaning in (\ref{ex5:base4}) would describe a situation in which Della \underline{did} kick it just so long as she didn't kick it previously.

What's needed is for the information that \emph{again} supplies about a prior eventuality to not be part of the main assertion of the clause. Instead, this information must be ``added'' to the assertion of the rest of the sentence. This is the same kind of thing, perhaps, that we see in cases like~(\ref{ex5:appositive}).%
%
\marginnote[-20pt]{An important difference between an appositive and \emph{again} is that material in an appositive seems to be genuinely ``outside'' the sentence semantically. This can be seen by noticing that pronouns inside cannot be understood as variables bound to DPs outside the appositive. In \emph{Every woman met Marcin, who she'd heard about for years}, \emph{she} cannot be a variable bound to \emph{every woman}. That is possible for information that \emph{again} conveys. In \emph{Every woman kicked her fridge again in front of me}, the previous kickings involve every woman's fridge. (See the discussion on the next page.)}%
%
\begin{exe}
	\raggedright
	\ex \label{ex5:appositive} 
	Della didn't kick Marcin, who teaches now in Canada.
\end{exe}
The so-called appositive relative clause, \emph{who teaches now in Canada}, isn't part of what \emph{nt} negates. Like the information supplied by \emph{again}, the information in the appositive of (\ref{ex5:appositive}) adds information that is separate from the main assertion of the clause. 

There are a variety of technical ways of expressing how the meaning of \emph{again} is separated from the meaning of the rest of the sentence it is in. \citet{vonStechow:1996} suggests making the contribution \emph{again} makes to the sentence a ``presupposition.'' It won't matter for our purposes here how this is done. Let us simply assume that the information about the existence of a previous eventuality that \emph{again} conveys is a separate message that accompanies the sentence that it is contained in.
\begin{exe}
	\raggedright
	\ex \label{ex5:againdenf} 
	\denotes{again} = $\lambda P \lambda \epsilon\ P(\epsilon) = T$.\\ Add to conversation: $\exists \epsilon '$ [$\epsilon ' < \epsilon \wedge\ P(\epsilon ')]$\\
	where $\lambda \epsilon P(\epsilon)=T$ is the denotation of the phrase\\ that \emph{again} combines with.
	\marginnote[-35pt]{The denotation in (\ref{ex5:againdenf}) isn't quite sufficient to capture all of what \emph{again} conveys. It isn't sufficient to simply require, as (\ref{ex5:againdenf}) does, that the previous eventuality merely exist. There are examples which indicate that \emph{again} makes reference to a particular, salient, eventuality that has occurred previously. See \citet{Beck:2006b}, who cites \citet{Heim:1990a}. It also needs to be added that the previous eventuality and the one that is described by the constituent that \emph{again} modifies, must lie between an eventuality in which the relevant property does not hold. In (\ref{ex5:base3}), for instance, the present kicking eventuality and the previous one that is being repeated must be separated by an eventuality which isn't a kicking. Without this additional element of meaning, a sentence like \emph{Paula is again crying} would have a meaning that overlaps with that of \emph{Paula is still crying}. See, for instance, \citet{Bale:2007}.}
\end{exe} 
With this denotation, (\ref{ex5:base4}) becomes~(\ref{ex5:base4new}).
\begin{exe}
	\raggedright
	\ex \label{ex5:base4new} 
	\begin{xlist}
		\ex $\neg \exists \epsilon$ \textsc{agent}(Della)($\epsilon$) $\wedge$ \textsc{kick}($\epsilon$) $\wedge$ \textsc{patient}(it)($\epsilon$)
		\sn $\approx$ there isn't an eventuality that Della is the agent of and is a kicking of it.
		\ex $\exists \epsilon '$ [$\epsilon ' < \epsilon$ $\wedge$ \textsc{kick}($\epsilon '$) $\wedge$ \textsc{patient}(it)($\epsilon '$)]
		\sn $\approx$ there was a previous eventuality that was a kicking of it.
	\end{xlist}
\end{exe}

For this to provide the right meaning for the sentence in (\ref{ex5:againbase1a}), repeated here, we'll have to understand the TP that \emph{again} combines with to be a description of eventualities. 
\begin{exe}
	\raggedright
	\exr{ex5:againbase1a} Again Della didn't kick it.
\end{exe}
We don't have enough of a semantics built up to include how tense makes its contribution, and so I won't attempt to work out how this can be achieved. What needs to emerge is for (\ref{ex5:againbase1a}) to get a meaning like that in~(\ref{ex5:againbase1ad}).
\begin{exe}
	\raggedright
	\ex \label{ex5:againbase1ad} 
	\begin{xlist}
		\ex $\neg \exists \epsilon$ \textsc{agent}(Della)($\epsilon$) $\wedge$ \textsc{kick}($\epsilon$) $\wedge$ \textsc{patient}(it)($\epsilon$)
		\sn $\approx$ there isn't an eventuality that Della is the agent of and is a kicking of it.
		\ex $\neg \exists \epsilon '$ [$\epsilon ' < \epsilon$ $\wedge$ \textsc{kick}($\epsilon '$) $\wedge$ \textsc{patient}(it)($\epsilon '$)]
		\sn $\approx$ there wasn't a previous eventuality that was a kicking of it by Della.
	\end{xlist}
\end{exe}
What's important here is that it's the phrase that \emph{again} is the sister to that determines the information that \emph{again} conveys. We can use this fact about \emph{again} to see what the meaning of its sister is.

Let's consider one more contrast that supports this conclusion. Consider the examples in~(\ref{ex5:bale1}), based on ones in a paper by Alan Bale.\cite{Bale:2007}
\begin{exe}
	\raggedright
	\sn
	\emph{Della kicked her broken refrigerator out of frustration while I was away. After I got home, \ldots}
	\ex \label{ex5:bale1}%
	%
	\marginnote[20pt]{The pronunciation of these sentences must not put pitch accent on \emph{again}. When \emph{again} is focussed, which is what pitch accent signals in English, the meaning of \emph{again} gets tangled up with the meaning of focus. See \citet{Beck:2006b}.}%
	%
	\begin{xlist}
		\ex[]{she kicked it again in front of me.}
		\ex[\#]{she kicked it in front of me again.}
	\end{xlist}
	\hspace*{\fill} \citep[inspired by (27):~462]{Bale:2007} \hspace*{.5in}
\end{exe}
I've set up the context for these two sentences so that it's clear what the relevant preceding eventuality is. Potential parses for the VPs in the two sentences are given in~(\ref{ex5:bale1p}).
\begin{fullwidth}
	\medskip
\begin{exe}
	\raggedright
	\ex \label{ex5:bale1p} 
	\small
	\begin{forest}
	[vP
		[DP [she, roof]]
		[\ibar{v} 
			[\obar{v} [agent] ]
			[VP [\ibar{V}
				[\ibar{V}
					[\ibar{V}
						[\obar{V} [kick] ]
						[DP [it, roof] ]
					 ]
					[AdvP [again, roof] ]
				 ]
				[PP [in front of me, roof] ]
			] ]
		]
	]
	\end{forest}
	\hspace*{10pt}
	\begin{forest}
	[vP
		[DP [she, roof]]
		[\ibar{v} 
			[\obar{v} [agent] ]
			[VP [\ibar{V}
				[\ibar{V}
					[\ibar{V}
						[\obar{V} [kick] ]
						[DP [it, roof] ]
					 ]
					[PP [in front of me, roof] ]
				 ]
				[AdvP [again, roof]]
			] ]
		]
	]
	\end{forest}
\end{exe}
\medskip
\end{fullwidth}
Imagine that \emph{in front of me} is a modifier that locates where the kicking eventuality takes place. That is:
\begin{exe}
	\raggedright
	\ex \label{ex5:denon} 
	\denotes{in the front door} is a function that maps eventualities to True when those eventualities are in front of me.
\end{exe}
The denotation of the \ibar{V} that \emph{again} combines with in (\ref{ex5:bale1}a) is (\ref{ex5:bale2}a), and the denotation of the \ibar{V} that \emph{again} combines with in (\ref{ex5:bale1}b) is~(\ref{ex5:bale2}b).
\begin{exe}
	\raggedright
	\ex \label{ex5:bale2} 
	\begin{xlist}
		\ex $\lambda \epsilon$ \textsc{kick}($\epsilon$) $\wedge$ \textsc{patient}(it)($\epsilon$)
		\ex $\lambda \epsilon$ \textsc{kick}($\epsilon$) $\wedge$ \textsc{patient}(it)($\epsilon$) $\wedge$ in\_front\_of\_me($\epsilon$)
	\end{xlist}
\end{exe}
When \emph{again} combines with these denotations, it introduces the meanings in~(\ref{ex5:bale3}).
\begin{exe}
	\raggedright
	\ex \label{ex5:bale3} 
	\begin{xlist}
		\ex $\exists$ $\epsilon '$ $\epsilon ' < \epsilon $ $\wedge$ \textsc{kick}($\epsilon '$) $\wedge$ \textsc{patient}(it)($\epsilon '$)
		\sn $\approx$ there was a previous kicking of it
		\ex $\exists$ $\epsilon '$ $\epsilon' < \epsilon$ $\wedge$ \textsc{kick}($\epsilon '$) $\wedge$ \textsc{patient}(it)($\epsilon '$) $\wedge$ in\_front\_of\_me($\epsilon '$)
		\sn $\approx$ there was a previous kicking of it in front of me.
	\end{xlist}
\end{exe}
In the contexts provided, only (\ref{ex5:bale3}a) is true, and so (\ref{ex5:bale1}b) doesn't fit the context (which is indicated by marking it with ``\#''). What we see here, then, is the same thing that was illustrated by (\ref{ex5:againbase1}). We can use \emph{again} to see what the meaning of \emph{again}'s sister is.

Consider now the meaning assigned to (\ref{ex5:vPa}).
\begin{exe}
	\raggedright
	\ex \label{ex5:vPa} Della again kicked it.
\end{exe}
The information that \emph{again} conveys in (\ref{ex5:vPa}) is that Della kicked it previously. That is, (\ref{ex5:vPa}) is synonymous with~(\ref{ex5:vPas}).
\begin{exe}
	\raggedright
	\ex \label{ex5:vPas} 
	Again Della kicked it.
\end{exe}
This is explained by the Derived Subjects Hypothesis, that is the idea that external arguments occupy a lower position at D-structure. If we set aside where \emph{again} is, the D-structure representation for (\ref{ex5:vPa}) is~(\ref{ex5:vPad}).
\begin{fullwidth}
\medskip
\begin{exe}
	\raggedright
	\ex \label{ex5:vPad} 
	\small
	\begin{forest}
	[TP[\ibar{T}
		[\obar{T} [past] ]
		[$\lambda \epsilon$ \textsc{agent}(Della)($\epsilon$) $\wedge$ \textsc{patient}(it)($\epsilon$)\\ vP
			[DP [Della, roof] ]
			[$\lambda x \lambda \epsilon$ \textsc{agent}($x$)($\epsilon$) $\wedge$ \textsc{patient}(it)($\epsilon$)\\ \ibar{v}
				[\obar{v}
					[\obar{V} [$\surd$kick, name=binder] ]
					[\obar{v} [agent] ]
				 ]
				[VP [\ibar{V}
					[{\ }, name=trace ]
					[DP [it, roof] ]
				] ]
			 ]
		 ]
	]]
	\draw[->,>=stealth', ultra thin, gray] (trace) to[out=south west,in=-80,looseness=1] node[pos=0.75,xshift=-10mm]{} (binder);
	\end{forest}
\end{exe}
\end{fullwidth}
\medskip
The constituent that \emph{again} combines with must be a description of eventualities. That is, it must be a function from eventualities to True or False. The denotation of \ibar{v} in (\ref{ex5:vPad}) doesn't have this meaning, and we should expect that \emph{again} cannot combine with it. The constituent that has the right meaning to be the sister to \emph{again} is the vP. This means that there must be a position to the left of the vP that \emph{again} resides in. Our present restriction on phrases provides only one way that can happen, and it's indicated in~(\ref{ex5:noagain}).
\begin{exe}
	\raggedright
	\ex \label{ex5:noagain} 
	\small
	\begin{forest}
	[\ibar{T}
		[\obar{T} [past]]
		[AdvP [\ibar{Adv}
			[\obar{Adv} [again] ]
			[vP
				[DP [Della, roof] ]
				[\ibar{v} [kick it, roof] ]
			 ]
		] ]
	]
	\end{forest}
\end{exe}
This gives to \emph{again} the same position that we've given to \obar{Neg}. It separates the vP from \obar{T} by an embedding head, the one occupied by \emph{again}. We might expect that this geometry would therefore have the same effects that the one in which \obar{Neg} occurs has. In chapter \ref{ch:features} we saw that when vP is embedded by an \obar{X} that separates it from \obar{T}, \textsc{agree} between \obar{T} and \obar{v} is blocked, and this results in ungrammaticality if there are features that need to get into an \textsc{agree}ment relation:
\begin{exe}
	\raggedright
	\ex[*]{\label{ex5:noagain2} Della not kicked it.}
\end{exe}
If this is correct, it's not clear why the similar representation in (\ref{ex5:noagain}) shouldn't have the same consequence.
\pagebreak

For this reason, let's allow phrases to have the shape in~(\ref{ex5:xpadj}).%
%
\marginnote[0pt]{Recall that ``\{$\alpha$,$\beta$\}'' means that $\alpha$ and $\beta$ are sisters, with no linear order imposed on $\alpha$ and $\beta$.}%
%
\begin{exe}
	\raggedright
	\ex \label{ex5:xpadj} 
	XP $\rightarrow$ \{YP, XP\}
\end{exe}
This isn't inconsistent with anything we've seen about English syntax so far, and it allows (\ref{ex5:vPa}) to get the parse in~(\ref{ex5:vPagt}).
\begin{fullwidth}
\medskip
\begin{exe}
	\raggedright
	\ex \label{ex5:vPagt} 
	\small
	\begin{forest}
	[\ibar{T}
		[\obar{T} [past] ]
		[vP
			[AdvP [\ibar{Adv} [\obar{Adv} [again]]] ]
			[$\lambda \epsilon$ \textsc{agent}(Della)($\epsilon$) $\wedge$ \textsc{kick}($\epsilon$) $\wedge$ \textsc{patient}(it)($e$)\\ vP
				[DP [Della, roof] ]
				[$\lambda x \lambda \epsilon$ \textsc{agent}($x$)($\epsilon$) $\wedge$ \textsc{kick}($\epsilon$) $\wedge$ \textsc{patient}(it)($\epsilon$)\\ \ibar{v}
					[\obar{v}
						[\obar{V} [$\surd$kick, name=binder] ]
						[\obar{v} [agent] ]
					 ]
					[VP [\ibar{V} 
						[{\ }, name=trace ]
						[DP [it, roof] ]
					] ]
				 ]
			 ]
		 ]
	]
	\draw[->,>=stealth', ultra thin, gray] (trace) to[out=south west,in=-80,looseness=1] node[pos=0.75,xshift=-10mm]{} (binder);
	\end{forest}
\end{exe}
\end{fullwidth}
\medskip
This D-structure produces the correct interpretation. The information that \emph{again} provides is~(\ref{ex5:vPam}).
\begin{exe}
	\raggedright
	\ex \label{ex5:vPam} $\exists$ $\epsilon '$ $\epsilon ' < \epsilon$ $\wedge$ \textsc{agent}(Della)($\epsilon '$) $\wedge$ \textsc{kick}($\epsilon '$) $\wedge$ \textsc{patient}(it)($\epsilon '$)
	\sn $\approx$ there was a previous kicking of it by Della
\end{exe}

As expected on this account, (\ref{ex5:vPa2}) can get the same interpretation that (\ref{ex5:vPa}) does.
\begin{exe}
	\raggedright
	\sn \emph{Della, frustrated by her noisy refrigerator, kicked it. Her anger unquenched, \ldots}
	\ex \label{ex5:vPa2} 
	Della kicked it again.
\end{exe}
This interpretation emerges from the D-structure representation in~(\ref{ex5:vPa2p}).
\begin{fullwidth}
\medskip
\begin{exe}
	\raggedright
	\ex \label{ex5:vPa2p} 
	\small
	\begin{forest}
	[\ibar{T}
		[\obar{T} [past] ]
		[vP
			[$\lambda \epsilon$ \textsc{agent}(Della)($\epsilon$) $\wedge$ \textsc{kick}($\epsilon$) $\wedge$ \textsc{patient}(it)($e$)\\ vP 
				[DP [Della, roof]]
				[$\lambda x \lambda \epsilon$ \textsc{agent}($x$)($\epsilon$) $\wedge$ \textsc{kick}($\epsilon$) $\wedge$ \textsc{patient}(it)($e$)\\ \ibar{v}
					[\obar{v}
						[\obar{V} [$\surd$kick, name=binder] ]
						[\obar{v} [agent] ]
					 ]
					[VP [\ibar{V} 
						[{\ }, name=trace ]
						[DP [it, roof] ]
					] ]
				 ]
			]
			[AdvP [\ibar{Adv} [\obar{Adv} [again]]] ]
		 ]
	]
	\draw[->,>=stealth', ultra thin, gray] (trace) to[out=south west,in=-80,looseness=1] node[pos=0.75,xshift=-10mm]{} (binder);
	\end{forest}
\end{exe}
\end{fullwidth}
\medskip
Because \emph{again} is combining with the same vP in (\ref{ex5:vPagt}) and (\ref{ex5:vPa2p}), the contribution it makes to the meaning of the sentence, i.e., (\ref{ex5:vPam}), is the same.

But (\ref{ex5:vPa2}) can have a meaning that (\ref{ex5:vPa}) cannot have. To see this, consider the situation portrayed in~(\ref{ex5:nosubj}).\marginnote[0pt]{This observation originates with \citet{McCawley:1973a}. The conclusions from this difference that I describe here, come from \citet{Bale:2007}, where it is also discovered that not all transitive verbs display this contrast. Bale uses this fact to argue that not all transitive verbs have a vP, and he reaches a similar conclusion about intransitive verbs. Indeed, his results are not entirely consistent with the portrayal of argument structure that has been developed in this book.}
\begin{exe}
	\raggedright
	\sn \emph{Shai, frustrated with his noisy refrigerator, kicked it and left for the relative quiet of the neighborhood coffee shop. Awoken by the refrigerator's racket, \ldots}
	\ex \label{ex5:nosubj} 
	\begin{xlist}
		\ex[\#]{Della again kicked it.}
		\ex[]{Della kicked it again.}
	\end{xlist}
\end{exe}
In this situation, the previous kicking was done by Shai, not Della. This is a context, then, which isn't consistent with the meaning in (\ref{ex5:vPam}), and as expected the representation for (\ref{ex5:nosubj}a) given in (\ref{ex5:vPagt}) produces an anomalous result. That the word order in (\ref{ex5:nosubj}b) does not have a similar effect suggests that there is a representation for this sentence that is different from (\ref{ex5:vPa2p}), which would produce the same anomalous meaning. Our parses for these sentences provide that representation; it is~(\ref{ex5:nosubjp}).
\begin{exe}
	\raggedright
	\ex \label{ex5:nosubjp}
	\small
	\begin{forest}
	[\ibar{T}
		[\obar{T} [past] ]
		[vP
			[DP [Della, roof] ]
			[\ibar{v}
				[\obar{v} [agent] ]
				[VP [\ibar{V}
					[$\lambda \epsilon$ \textsc{kick}($\epsilon$) $\wedge$ \textsc{patient}(it)($\epsilon$)\\\ibar{v} [kick it, roof] ]
					[AdvP [\ibar{Adv} [\obar{Adv} [again]]] ]
				] ]
			 ]
		 ]
	]
	\end{forest}
\marginnote[0pt]{Our new syntax would also allow for \emph{again} to combine with the VP headed by \emph{kick}, rather than its \ibar{V}.}
\end{exe}
This representation causes the meaning that \emph{again} contributes to the sentence to be (\ref{ex5:nosubjm}), which is consistent with the context in~(\ref{ex5:nosubj}).
\begin{exe}
	\raggedright
	\ex \label{ex5:nosubjm} 
	$\exists \epsilon '$ $\epsilon ' < \epsilon$ $\wedge$ \textsc{kick}($\epsilon '$) $\wedge$ \textsc{patient}(it)($\epsilon '$)
	\sn $\approx$ there was a previous kicking of it
\end{exe}
This supports the hypothesis adopted in chapter \ref{ch:positargs} that the position the external $\theta$-role occupies is different than the position the verb occupies. For these facts to emerge, it requires that there be a phrase that contains the verb that is distinct from the phrase that contains the external $\theta$-role.

With this in mind, let's reconsider the alternation illustrated by~(\ref{ex5:melt1}).
\begin{exe}
	\raggedright
	\ex \label{ex5:melt1} 
	\begin{xlist}
		\ex Ivy closed the boxes.
		\ex The boxes closed.
	\end{xlist}
\end{exe}
Our model of this is that the vocabulary item \emph{close} spells out the root $\surd$close with or without an accompanying agent $\theta$-role. Because the agent $\theta$-role is responsible for assigning accusative Case, its presence is also responsible for licensing the internal argument \emph{the boxes} in a lower position. When agent isn't present, the internal argument gets nominative Case, and this leads to its movement into clause initial position. The derivations sketched in (\ref{ex5:meltder}) illustrate.
\begin{fullwidth}
	\medskip
\begin{exe}
	\medskip
	\raggedright
	\ex \label{ex5:meltder} 
	\small
	\begin{forest}
	[TP
		[DP [Ivy, name=binders, roof]]
		[\ibar{T} 
			[\obar{T} [past] ]
			[vP
				[{\ }, name=traces ]
				[\ibar{v}
					[\obar{v}
						[\obar{V}, name=binderv
							[\obar{V} [$\surd$close] ]
							[\obar{$\mu$} [case: \emph{acc}] ]
						 ]
						[\obar{v} [agent] ]
					 ]
					[ $\mu$P	
						[DP [the boxes, name=bindero, roof] ]
						[\ibar{$\mu$}
							[{\ }, name=tracev ]
							[VP [\ibar{V}
								[{\ }, name=trace ]
								[{\ }, name=traceo ]
							] ]
						 ]
					]
				 ]
			 ]
		]
	]
	\draw[->,>=stealth', ultra thin, gray] (traceo) to[out=south west,in=south,looseness=1] node[pos=0.75,xshift=-10mm]{} (bindero);
	\draw[->,>=stealth', ultra thin, gray] (traces) to[out=-110,in=-80,looseness=1.3] node[pos=0.75,xshift=-10mm]{} (binders);
	\draw[->,>=stealth', ultra thin, gray] (tracev) to[out=-110,in=south,looseness=1] node[pos=0.75,xshift=-10mm]{} (binderv);
	\draw[->,>=stealth', ultra thin, gray] (trace) to[out=west,in=-80,looseness=1] node[pos=0.75,xshift=-10mm]{} (tracev);
	\end{forest}
\end{exe}
\end{fullwidth}
As expected, the transitive version of \emph{close} supports the ambiguity that arises when \emph{again} is placed at the end of the VP. There is a meaning in which what \emph{again} conveys is that there was a previous closing of the boxes by Ivy, as in (\ref{ex5:repivy1}), and there is a meaning which conveys merely that there was a previous closing of the boxes, as in~(\ref{ex5:repivy2}).\enlargethispage*{20pt}
\begin{exe}
	\raggedright
	\sn \emph{Ivy is responsible for boxing up the diodes for shipment at the end of the line. She had just finished closing all the boxes, when she realized she'd forgotten to include the inventory sheets that are supposed to accompany each shipment. After opening them all up and including the sheets,\ldots}
	\ex \label{ex5:repivy1} Ivy again closed the boxes.\\ Ivy closed the boxes again.
	\sn \emph{Mako is responsible for boxing up the diodes for shipment at the end of the line. She had just finished closing all the boxes, when she realized she'd forgotten to include the inventory sheets that are supposed to accompany each shipment. She opened them all up and included the sheets, but had to leave before finishing the job. When the truck arrived to take them off,\ldots}
	\ex \label{ex5:repivy2} \#Ivy again closed the boxes.\\ \hspace*{5.5pt}Ivy closed the boxes again.
\end{exe}
These meanings arise because the D-structure for the transitive sentence has both a vP and a VP for \emph{again} to combine with. 

There is no vP in the intransitive use of \emph{close}, however, and as expected therefore, \emph{again} cannot produce a reading that would arise by combining with vP. Moreover, the meaning \emph{again} conveys must include the object, as the contrast below indicates.
\pagebreak
\begin{exe}
	\sn \emph{The company in an effort to increase efficiency had started using self-closing boxes. A timer was set at the beginning of each work day that would cause the boxes to automatically close and seal at the end of the day. Ivy, used to years of closing the boxes, frequently forgot about the new self-closing boxes, and would close them herself as they came into her station. Today was one of the days that Ivy mistakenly closed the boxes, but she caught her mistake before the inspector arrived and re-opened them all so that\ldots}
	\raggedright
	\ex[]{\label{ex5:meltrep1} the boxes closed again at the end of the day as planned.}
	\sn \emph{The last of the self-closing boxes for the first run of the morning had just closed when the next group of boxes, on a delayed timer, arrived at Ivy's station. After being filled,\ldots}
	\ex[\#]{\label{ex5:meltrep2}the new boxes closed again.}
\end{exe}
A quick look at the D-structure for the intransitive case reveals why.
\begin{exe}
	\raggedright
	\ex \label{ex5:intmeltd} 
	\small
	\begin{forest}
	[VP [$\lambda \epsilon$ \textsc{close}($\epsilon$) $\wedge$ \textsc{patient}(the\_boxes)($\epsilon$)\\ \ibar{V}
		[$\lambda x \lambda \epsilon$ \textsc{close}($\epsilon$) $\wedge$ \textsc{patient}($x$)($\epsilon$)\\\obar{V} [$\surd$close] ]
		[DP [the boxes, roof] ]
	]
	]
	\end{forest}
\end{exe}
The only projection of \obar{V} that has the right meaning for \emph{again} to combine with is \ibar{V} and this constituent necessarily contains the object. To get the meaning necessary for the context in (\ref{ex5:meltrep2}) would require \emph{again} to combine with the verbal root \emph{$\surd$close} directly, and this isn't semantically possible. 

Interestingly, VPs headed by \emph{close} support another meaning when \emph{again} is present. To see this, consider the situation in~(\ref{ex5:rest1}).
\begin{exe}
	\raggedright
	\sn \emph{The boxes arrived in the Receiving department as they usually do closed up ready for use. Robert opened them up in preparation for filling, and sent them to Ivy's station. After filling them,\ldots}
	\ex \label{ex5:rest1} Ivy closed the boxes again.
\end{exe}
In this scenario, there is no previous event of closing the boxes. The boxes start this story in a closed state, and it is to this state that (\ref{ex5:rest1}) says Ivy returned them. This meaning exists also for the intransitive use of \emph{close}. 
\begin{exe}
	\raggedright
	\sn \emph{I was delighted when the tulips I'd planted the previous Fall finally bloomed. Their petals opened one warm afternoon to reveal a subtle multicolored pattern inside. But unaware of Nyctinasty, I was puzzled when\ldots}
	\ex \label{ex5:restint1} the tulips closed again that night.
\end{exe}
In this story, the tulips have never closed previously. They start in a closed state, and it is to that state that (\ref{ex5:restint1}) says they have returned. This meaning has a name. It is called the ``restitutive'' reading of \emph{again}, and it is to be contrasted with the two meanings we have seen previously, which are known as ``repetitive'' readings.%
%
\marginnote[0pt]{\citet{Dowty:1979} contains perhaps the first extended discussion of this ambiguity. The account of the restitutive reading adopted here is the one provided in \citet[]{vonStechow:1995,vonStechow:1996}.}

We can deduce from the presence of the restitutive reading that there is a constituent within the VPs of these sentences whose meaning describes the state of being closed. It is to that constituent that \emph{again} must be combining to give rise to these meanings. That constituent has a meaning rather like that of the adjective \emph{closed}, as we can see by considering~(\ref{ex5:aprep}).
\begin{exe}
	\raggedright
	\ex \label{ex5:aprep} 
	The tulips are again closed.
\end{exe}
If we set aside the meaning of \emph{closed} as a passive participle, and focus on the stative meaning it has when it is an adjective, then the message \emph{again} conveys in (\ref{ex5:aprep}) is the same that is conveyed in (\ref{ex5:restint1}).%
%
\marginnote[0pt]{For the difference between passive participles and adjectives, see chapter~\ref{sec:argmove}.} Like (\ref{ex5:restint1}), (\ref{ex5:aprep}) says that the tulips are restored to a closed state. Note that the verb \emph{closed} does not support that meaning when \emph{again} precedes the verb, as in~(\ref{ex5:leftrep}).%
%
\begin{exe}
	\raggedright
		\sn \emph{I was delighted when the tulips I'd planted the previous Fall finally bloomed. Their petals opened one warm afternoon to reveal a subtle multicolored pattern inside. But unaware of Nyctinasty, I was puzzled when\ldots}
	\ex[\#]{\label{ex5:leftrep}The tulips again closed that night.}	
\end{exe}
The constituent that describes a state of being closed is not syntactically available to \emph{again} in the position it occupies in~(\ref{ex5:leftrep}). 

All of this suggests that the VP associated with the verb \emph{close} contains more syntactic structure than we have given to it. In addition to the optional presence of the external $\theta$-role, there must be two phrases: one that denotes the state of being closed and another that describes the process of coming into that state. A parse that achieves that is~(\ref{ex5:restintrop}).
\begin{exe}
	\raggedright
	\ex \label{ex5:restintrop} 
	\small
	\begin{forest}
	[\ibar{V}
		[\obar{V} [become] ]
		[VP [\ibar{V}
			[\obar{V} [$\surd$close] ]
			[DP [the tulips, roof] ]
		]]
	]
	\end{forest}
\end{exe}
The intransitive \emph{close}, on this view, is made of two syntactically distinct predicates. One, \emph{$\surd$close}, describes states of being closed and takes, in this instance, \emph{the tulips} as its argument. The other, here indicated by \emph{become}, describes eventualities which lead to states. Roughly:
\begin{exe}
	\raggedright
	\ex \label{ex5:becomed} 
	\denotes{become} = $\lambda P$$\lambda \epsilon$ $\epsilon$ results in $\exists \epsilon '$ P($\epsilon$)=T
\end{exe}
When the VP in (\ref{ex5:restintrop}) is embedded under the vP headed by the agent $\theta$-role, we get the transitive version of \emph{close}. Because the vocabulary item \emph{close} can be matched to either the combination involving \obar{v} or not, we will want all of the morphemes that make up the meaning of \emph{close} to be in an \textsc{agree} relation. If those morphemes come equipped with a $\bullet$, it will require head movement to gather all of them into one \obar{X} position, as in (\ref{ex5:bigclosem}) for instance.
\begin{exe}
	\raggedright
	\ex \label{ex5:bigclosem} 
	\small
	\begin{forest}
	[vP
		[DP [Ivy, roof]]
		[\ibar{v} 
			[\obar{v}
				[\obar{V}, name=binder
					[\obar{V} [$\surd$close] ]
					[\obar{V} [become] ]
				 ]
				[\obar{v} [agent] ]
			 ]
			[VP [\ibar{V}
				[{\ }, name=tracev]
				[VP [\ibar{V}
					[{\ }, name=trace ]
					[DP [the boxes, roof] ]
				] ]
			] ]
		]
	]
	\draw[->,>=stealth', ultra thin, gray] (trace) to[out=south west,in=south,looseness=1] node[pos=0.75,xshift=-10mm]{} (tracev);
	\draw[->,>=stealth', ultra thin, gray] (tracev) to[out=south west,in=south,looseness=1] node[pos=0.75,xshift=-10mm]{} (binderv);
	\end{forest}
\end{exe}
The restitutive meaning for \emph{again} arises when it is combined with a projection of \emph{$\surd$close}.

Confirmation for this picture comes from considering how \emph{again} interacts with modifiers that distinguish between the stative meaning of \emph{$\surd$close} and the non-stative meaning of \emph{become}. If we take the adjective \emph{close} to represent the meaning of \emph{$\surd$close}, then (\ref{ex5:advclose1}) shows us that \emph{tightly} can modify it but \emph{quickly} cannot.
\begin{exe}
	\raggedright
	\ex \label{ex5:advclose1} 
	\begin{xlist}
		\ex[]{The boxes remained closed tightly.}
		\ex[*]{The boxes remained closed quickly.}
	\end{xlist}
\end{exe}
The model in (\ref{ex5:bigclosem}) predicts correctly that (\ref{ex5:noresta}) should not be able to support the restitutive reading because its syntax must be as indicated in~(\ref{ex5:advclose2}).
\begin{exe}
	\raggedright
	\sn \emph{I was delighted when the tulips I'd planted the previous Fall finally bloomed. Their petals opened one warm afternoon to reveal a subtle multicolored pattern inside. But unaware of Nyctinasty, I was puzzled when\ldots}
	\ex[\#]{\label{ex5:noresta} The tulips closed quickly again that night.}
	\ex \label{ex5:advclose2}
	\small
	\begin{forest}
	[VP[\ibar{V}
		[\ibar{V} 
			[\ibar{V}
				[\obar{V} [become] ]
				[VP [\ibar{V}
					[\obar{V} [$\surd$close] ]
					[DP [the tulips, roof] ]
				] ]
			 ]
			[AdvP [quickly, roof] ]
		]
		[AdvP [again, roof]]
	]]
	\end{forest}
\end{exe}
By contrast, (\ref{ex5:bigclosem}) predicts that the restitutive reading should be found in (\ref{ex5:norestb}) because the representation in (\ref{ex5:advclose2}) is available to it.
\begin{exe}
	\raggedright
	\sn \emph{I was delighted when the tulips I'd planted the previous Fall finally bloomed. Their petals opened one warm afternoon to reveal a subtle multicolored pattern inside. But unaware of Nyctinasty, I was puzzled when\ldots}
	\ex[]{\label{ex5:norestb} The tulips closed tightly again that night.}
	\ex \label{ex5:advclose3}
	\small
	\begin{forest}
	[VP[\ibar{V}
		[\obar{V} [become] ]
		[VP [\ibar{V}
			[\ibar{V}
				[\ibar{V}
					[\obar{V} [$\surd$close] ]
					[DP [the tulips, roof] ]
				 ]
				[AdvP [tightly, roof] ]
			 ]
			[AdvP [again, roof] ]
		] ]
	]]
	\end{forest}
\end{exe}

The presence of the repetitive and restitutive reading with \emph{again} is a phenomenon widely found in the world's languages, but there is considerable variation in how the readings correlate with the syntax. For instance, Minjeong Son\cite{Son:2006} shows that \emph{tasi}, the Korean equivalent of \emph{again}, invokes the same restitutive/repetitive ambiguity that we have seen in English, as in~(\ref{ex5:Korean}).
\pagebreak
\begin{exe}
	\raggedright
	\ex \label{ex5:Korean} 
	\gll Yenghi-ka koki-lul tasi el-le-ess-ta. \\ Yenghi-\textsc{nom} meat-\textsc{acc} again freeze-\textsc{caus}-\textsc{pst}-\textsc{dec} \\ `Yenghi froze the meat again'
	\begin{xlist}
		\ex $\approx$ Yenghi froze the meat and it had been frozen previously, or
		\ex $\approx$ Yenghi froze the meat and someone had done that previously.
	\end{xlist}
	\hspace*{\fill} \citep[(49):~59)]{Son:2006} \hspace*{.5in} 
\end{exe}
Notice that in Korean the verbal morphology overtly signals that there is a stative part (\emph{el}) and a ``become'' part (\emph{le}, glossed as \textsc{caus}). A natural way of modeling the difference between English and Korean is letting the two \obar{V} that makeup the relevant vocabulary item map onto distinct exponents in Korean, which are brought together by Head Movement.
\begin{exe}
	\raggedright
	\ex \label{ex5:korill1} 
	Yenghi-ka [\VP [\VP koki-lul \linkfrom{\ } \under ] [_{\text{\obar{V}}}\to{el} le] ]
\end{exe}
Korean syntax allows object DPs to be moved to the left. This process allows Korean to have S-structures like that sketched in~(\ref{ex5:korill2}).
\begin{exe}
	\raggedright
	\ex \label{ex5:korill2} 
	Yenghi-ka [\XP \linkto{koki-lul} \under [\VP [\VP \from{\ }  ] [_{\text{\obar{V}}}el le] ]]
\end{exe}
The ambiguity of (\ref{ex5:Korean}), then, arises because (\ref{ex5:korill2}) can be related to either of the D-structures in~(\ref{ex5:korill3}).
\begin{exe}
	\raggedright
	\ex \label{ex5:korill3} 
	\begin{xlist}
		\ex Yenghi-ka [\VP tasi [\VP [\VP koki-lul [_{\text{\obar{V}}}el] ] [_{\text{\obar{V}}}le] ]]
		\ex Yenghi-ka [\VP [\VP tasi [\VP koki-lul [_{\text{\obar{V}}}el] ] [_{\text{\obar{V}}}le] ]]
	\end{xlist}
\end{exe}
In (\ref{ex5:korill3}a), \emph{tasi} combines with the VP headed by \emph{le}, giving rise to a repetitive reading, and in (\ref{ex5:korill3}b), \emph{tasi} combines with the VP headed by \emph{el}, producing the restitutive reading. While it is easier to see in the verbal morphology of Korean the two components of meaning that the repetitive and restitutive readings dependent on, it is harder to see from the word-order of Korean that the syntactic placement of \emph{tasi} (``again'') correlates with this ambiguity.%
%
\marginnote[-10pt]{A challenge in the cross-linguistic study of \emph{again} is distinguishing it from adverbs whose meaning is closer to \emph{back}. One can appreciate the closeness in the meanings of these two expressions by comparing \emph{Meg went back into the room} and \emph{Meg went into the room again}. See \citet{Beck:2005} and \citet{Beck:2008} for discussion.}%
%
 A rapidly evolving project presently is understanding the cross-linguistic typology of adverbs like \emph{again} and what they reveal about the hidden structure of~VPs.

\subsection{Double Objects} 
\label{sub:doubleobjects}

Our present system derives the syntactic position of accusative Case marked arguments by putting them into an \textsc{agree} relation with \obar{$\mu$} and forcing them to move to Specifier of $\mu$P. This predicts that there should never be two accusative Case-marked DPs dependent on the same \obar{$\mu$}, and one expectation from that prediction is that no verb should have two accusative Case-marked objects. That expectation is frustrated by examples such as~(\ref{ex5:do1}).
\pagebreak
\begin{exe}
	\raggedright
	\ex \label{ex5:do1} 
	\begin{xlist}
		\ex Veena gave Rejean the answer.
		\ex Veena kicked Rejean the ball.
		\ex Veena sent Rejean a letter.
		\ex Veena baked Rejean some cookies.
		\ex Veena taught Rejean a song.
		\ex Veena promised Rejean the answer.
		\ex \label{ex5:deny}Veena denied Rejean his paycheck.
		\ex \label{ex5:envy}Veena envied Rejean his quick footwork.
	\end{xlist}
\end{exe}
These examples are sometimes described as involving the ``double object'' construction. The existence of the second DP is unexpected given our Case theory.

In the best known dialects of English, there is evidence that the Case that is borne by the second DP is not a structural Case. Recall that structural Case is what the Passive operation redirects.%
%
\marginnote[0pt]{See chapter~\ref{sec:thetacase} for a discussion of structural Case and the passive.} The contrast between (\ref{ex5:passdo1}) and (\ref{ex5:passdo2}) indicates that only the first DP bears structural Case.%
%
\begin{exe}
	\raggedright
	\ex \label{ex5:passdo1} 
	\begin{xlist}
		\ex Rejean was given the answer.
		\ex Rejean was kicked the ball.
		\ex Rejean was sent a letter.
		\ex Rejean was baked some cookies.
		\ex Rejean was taught a song.
		\ex Rejean was promised the answer.
		\ex Rejean was denied his paycheck.
		\ex Rejean was envied his quick footwork.
	\end{xlist}
	\ex \label{ex5:passdo2}
	\begin{xlist}
		\ex[*]{The answer was given Rejean.}
		\ex[*]{The ball was kicked Rejean.}
		\ex[*]{A letter was sent Rejean.}
		\ex[*]{Some cookies were baked Rejean.}
		\ex[*]{A song was taught Rejean.}
		\ex[*]{The answer was promised Rejean.}
		\ex[*]{His paycheck was denied Rejean.}
		\ex[*]{His quick footwork was envied Rejean.}
	\end{xlist}
\end{exe}
Our syntax, then, should be able to recognize the different status of the two Cases involved. 

The double object construction is common cross-linguistically and this asymmetry in the status of the two Cases is not always found. For instance in the Bantu language Kichaga either of the two arguments of a double object construction can switch to Nominative in the passive.%
%
\marginnote[-50pt]{The ``Nominative'' Case relation is not signaled by morphology on the DP in Kichaga,  but is instead signaled, as it is in Bantu generally, by ``subject agreement'' morphology. On the correspondence between subject agreement and nominative Case, see chapter~\ref{ch:features}.}%
%
\begin{exe}
	\raggedright
	\ex \label{ex5:kichaga} 
	\begin{xlist}
		\ex \gll \`{M}-k\`{a} n-\textipa{\H{a}}-\textipa{\H{\i}}-ly\`{i}-\'i-\`{o} k-\`{e}ly\^{a}. \\ 1-wife \textsc{foc}-1S-\textsc{pr}-eat-\textsc{ap}-\textsc{pas} 7-food  \\ `the wife is being affected by the food being eaten.'
		\ex \gll K-\`{e}ly\'{a} k-\textipa{\H{\i}}-ly\`{i}-\'{i}-\`{o} \`{m}-k\`{a}. \\ 7-food 7S-\textsc{pr}-eat-\textsc{ap}-\textsc{pas} 1-wife  \\ `the bood is being eaten for/on the wife'
	\end{xlist}
	\hspace*{\fill} \citep[(5):p.~150]{Bresnan:1990} \hspace*{.5in}
\end{exe}
The same is found in Norwegian.\marginnote[0pt]{These data originally from Tor {\AA}farli.}
\begin{exe}
	\raggedright
	\ex \label{ex5:norwegian} 
	\begin{xlist}
		\ex \gll Jon vart gitt ei klokke. \\ John was given a watch \\ `John was given a watch.'
		\ex \gll Ei klokke vart gitt Jon. \\ a watch was given Jon. \\ 
	\end{xlist}
	\hspace*{\fill} \citep[(2):p.~682]{Woolford:1993} \hspace*{.5in}
\end{exe}
There are also dialects of English that permit the second object to convert to Nominative in the passive. In Manchester English, for instance, \citet{Haddican:2010} reports that when the object is a pronoun it may convert to Nominatives; examples like \ref{ex5:manch} are grammatical.
\begin{exe}
	\raggedright
	\ex \label{ex5:manch} 
	\begin{xlist}
		\ex It was sent him.
		\ex It was given her.
	\end{xlist}
\end{exe}
And \citet{Biggs:2016} reports that this is possible for definites in Liverpool English; examples like (\ref{ex5:liver}) are grammatical.
\begin{exe}
	\raggedright
	\ex \label{ex5:liver} 
	\begin{xlist}
		\ex The book was given the teacher.
		\ex The package was sent her nan's.
	\end{xlist}
\end{exe}
These double objects are known as ``symmetric,'' and there are a variety of properties that correlate with the distinction between a symmetric and asymmetric double object construction. A complete account of the double-object construction would explain these correlations. This remains an area in which no consensus has emerged, and I will set this goal aside.%
%
\marginnote[0pt]{See \citet{Gary:1977}, \citet{Baker:1988c}, \citet{Bresnan:1990}, \citet{Woolford:1993}, \citet{Anagnostopoulou:2003}, \citet{Haddican:2019}, and references cited therein, for some discussion.}%
%

Some of the verbs that allow for the double object construction permit another expression of their arguments, as (\ref{ex5:do2}) shows.
\enlargethispage*{15pt}
\begin{exe}
	\raggedright
	\ex \label{ex5:do2} 
	\begin{xlist}
		\ex[]{Veena gave the answer to Rejean.}
		\ex[]{Veena kicked the ball to Rejean.}
		\ex[]{Veena sent the letter to Rejean.}
		\ex[]{Veena baked some cookies for Rejean.}
		\ex[]{Veena taught a song to Rejean.}
		\ex[*]{Veena promised the answer to/for Rejean.}
		\ex[*]{Veena denied his paycheck to/for Rejean.}
		\ex[*]{Veena envied his quick footwork to/for Rejean.}
	\end{xlist}
\end{exe}
There is reason to believe that the relationship between these two ways of expressing the arguments of the relevant verbs is rule governed, that is, it's not merely that speakers of English have learned for each verb that they may express their internal arguments in these two frames. One reason for thinking this is that when new words that have the right properties are coined, speakers associate that word with both argument structure frames even if they've only been presented with one of them. For instance, Alec Marantz notes that the verb \emph{shin}, can be introduced to a speaker with (\ref{ex5:prod1a}), and the speaker will readily use~(\ref{ex5:prod1b}). 
\begin{exe}
	\raggedright
	\ex \label{ex5:prod1} 
	\begin{xlist}
		\ex \label{ex5:prod1a} Brian shinned the ball to Jason.
		\ex \label{ex5:prod1b} Jason shinned Aynat the ball. 
	\end{xlist}
\end{exe}
% //todo Put in reference to Marantz for shin.

There is experimental evidence that supports this conclusion as well. \citet{Gropen:1989}, which is a study of the acquisition of the double object construction by children, demonstrated for children that they knew that a double object frame was possible for a verb if they were taught something about that verb's argument structure. They presented children with scenarios where a character transports another character in a gondola car to some goal, with sentences like ``The bear is pilking the pig to the giraffe.'' They were then presented with similar scenarios and asked to describe those scenarios. They found that children in the age range of 6 to 8 years old would respond with double object constructions (for example: \emph{X is pilking Y Z}) 44\% of the time. Knowing something about the meaning of a verb---perhaps its argument structure---confers knowledge about whether that verb can support the double object construction.

One way of capturing the relationship between the two argument structures is Transformational. We could imagine that there is a rule, or a series of rules, that convert a common D-structure representation into one or the other of these surface representations. This is the view taken by \citet{Larson:1988}, who builds on \citet{Chomsky:1975a}. Larson proposes a rule that takes a D-structure like (\ref{ex5:larson1}) to a surface structure like~(\ref{ex5:larson2}).%
%
\marginnote[0pt]{I am converting Larson's proposal into the framework developed here. Larson's work predated the adoption of vP, though one could argue that the vP hypothesis was nascent in his own syntax. He proposed that there are two VPs in the underlying structure of clauses, and that it is the verb heading the higher of these VPs that assigns structural (i.e., accusative) Case. This is exactly what the \obar{v} of vP is.}%
%
\begin{exe}
	\raggedright
	\ex \label{ex5:larson1} 
	\small
	\begin{forest}
	[vP
		[DP [Ania, roof]]
		[\ibar{v} 
			[\obar{v} ]
			[VP
				[DP [the candy, roof] ]
				[\ibar{V}
					[\obar{V} [$\surd$give] ]
					[PP [to Stephen, roof] ]
				 ]
			 ]
		]
	]
	\end{forest}
	\normalsize
	\ex \label{ex5:larson2}
	\small
	\begin{forest}
	[vP
		[DP [Ania, roof]]
		[\ibar{v} 
			[\obar{v} ]
			[VP
				[DP [Stephen, name=binder, roof] ]
				[\ibar{V} 
					[\ibar{V}
						[\obar{V} [$\surd$give] ]
						[{\ }, name=trace ]
					 ]
					[DP [the candy, roof] ]
				]
			 ]
		]
	]
	\draw[->,>=stealth', ultra thin, gray] (trace) to[out=south west,in=south,looseness=1] node[pos=0.75,xshift=-10mm]{} (binder);
	\end{forest}
\end{exe}
Note that he places the \textsc{theme} argument in Specifier of VP at D-structure.

His proposed Transformation consists of two rules that work in concert. One deletes the preposition, here \emph{to}, thereby depriving its argument of Case. The second moves the DP in Specifier of VP to a position that makes it sister to \ibar{V}. The Case filter is satisfied by using Argument Movement to bring the DP that was dependent on the preposition for Case into Specifier of VP. There it can get Case from \obar{v}. The DP that is now sister to \ibar{V} gets Case from a special rule that interprets the \ibar{V} it is a sister to as a Case assigner. 

As in the theory developed here, the D-structure representation on Larson's view is responsible for the semantic interpretation of the sentence. A way of understanding the alternation, then, is to posit that knowing that a verb takes the two arguments shown in (\ref{ex5:larson1}) licenses this D-structure and the consequent transformations that give rise to the double object frame. When an English speaker learns that the meaning of \emph{give} determines the D-structure in (\ref{ex5:larson1}), the suite of transformations that the speaker has learned constitute part of the grammar of English are available to produce (\ref{ex5:larson2}). Thus, learning the argument structure of \emph{give} in English confers the knowledge that it will allow both frames, thus explaining facts like Marantz's. 

For this to be workable, discrepancies in the meanings of some of the alternants will have to be accounted for. Consider, for instance, the pair in~(\ref{ex5:bake}).
\begin{exe}
	\raggedright
	\ex \label{ex5:bake} 
	\begin{xlist}
		\ex Johannes baked cookies for Christine.
		\ex Johannes baked Christine cookies.
	\end{xlist}
\end{exe}
(\ref{ex5:bake}a) describes eventualities in which Johannes baked cookies on behalf of Christine, but (\ref{ex5:bake}b) does not. (\ref{ex5:bake}b) requires that Johannes baked cookies for a particular purpose: for Christine to have them. To capture this, Larson suggests that there is more than one \emph{for}, and that they are distinguished by their meaning. The transformation that deletes \emph{for}, thereby causing its argument to undergo A Movement, targets the particular \emph{for} that denotes the ``intends to have'' type relation. For those verbs that do not have the frame for expressing their arguments that involve a preposition, Larson suggests that the process of preposition deletion is obligatory, perhaps because English doesn't have a vocabulary item that corresponds to that preposition. In this way, then, \emph{promise}, \emph{deny}, and verbs like them obligatorily undergo the transformations that give rise to the double object frame. The transformations that are responsible for the alternation, then, will have to map one particular argument structure onto a D-structure that gives rise to either S-structure.

But the argument structure of a verb doesn't seem sufficient to predict whether a verb will be able to participate in the alternation. For instance, \emph{donate} and \emph{give} seem to involve the same $\theta$-roles, and yet only \emph{give} allows for the double object frame.
\begin{exe}
	\raggedright
	\ex \label{ex5:donate} 
	\begin{xlist}
		\ex[]{Ania donated candy to Stephen.}
		\ex[*]{Ania donated Stephen candy.}
	\end{xlist}
\end{exe}
If the productivity of the relationship is merely tied to argument structure, then the difference between \emph{give} and \emph{donate} would have to arise because they do not have the same $\theta$-roles, or are in some other way different in their argument structure. But this seems difficult to maintain. Note that \emph{X donates Y to Z} entails \emph{X gives Y to Z} whereas \emph{X lends Y to Z}, does not. And yet \emph{lend} unlike \emph{donate}, allows for the alternation. \citet{Pinker:1989} suggests that what distinguishes \emph{donate} from \emph{give} is that \emph{donate} derives historically from Latin, while \emph{give} has its roots in Germanic. One way of conceiving of this is that the process of learning involves the ability to subdivide verbs into arbitrary classes, and this process is sensitive enough to phonological patterns to be able to distinguish those verbs which inherit characteristics of Latin morpho-phonology from those which don't.%
%
\marginnote[0pt]{See \citet{Yang:2016} for a good discussion of this issue.}%
%
One of the characteristics is length: English verbs of Latinate origin have, on average, more syllables than those of Germanic origin. What an English speaker might have acquired, then, is a subcategorization of verbs that includes criteria like their argument structure and phonological pattern, in addition to the transformational processes that deliver the two frames of the alternation from a common D-structure to one of those subcategorizations. These transformational processes would have to be keyed to the relevant subcategory of verb. 

Or consider the problem posed by the contrast in (\ref{ex5:pull1}), also discussed in \citet[pp. 110-1]{Pinker:1989}.
\begin{exe}
	\raggedright
	\ex \label{ex5:pull1} 
	\begin{xlist}
		\ex[]{Ania tossed Stephen a ball.}
		\ex[]{Ania threw Stephen a ball.}
		\ex[]{Ania flipped Stephen a ball.}
		\ex[]{Ania kicked Stephen a ball.}
		\ex[*]{Ania pulled Stephen a ball.}
		\ex[*]{Ania schlepped Stephen the balls.}
		\ex[*]{Ania lifted Stephen the balls.}
		\ex[*]{Ania hauled Stephen the balls.}
	\end{xlist}
\end{exe}
Both \emph{toss} and \emph{pull} describe eventualities in which something is moved and something at whose location the moved term ends up. Both sentences in (\ref{ex5:pull2}) entail that a ball moved and that its trajectory aims at Stephen.
\begin{exe}
	\raggedright
	\ex \label{ex5:pull2} 
	\begin{xlist}
		\ex Ania tossed a ball to Stephen.
		\ex Ania pulled a ball to Stephen.
	\end{xlist}
\end{exe}
Our way of modeling this is to give \emph{toss} and \emph{pull} the same $\theta$-roles. And yet this match in argument structure is not sufficient to allow \emph{pull} to express its arguments in the double object frame. Both of these verbs are Germanic in origin, and they do not differ in the morpho-phonological ways that mark the distinction between Latinate and Germanic origin. Pinker suggests that what might distinguish these cases is in the structure of the eventuality that these verbs describe. He notes that \emph{toss} describes a punctual eventuality: what makes an eventuality a tossing is the kind of action that puts an item into motion. By contrast, \emph{pull} describes an eventuality that includes the trajectory of motion and cannot therefore be punctual. Pinker notes that many verbs which contrast along this dimension also differ in the way that (\ref{ex5:pull1}) illustrates.%
%
\marginnote[0pt]{See \citet{Krifka:1999} and \citet{Krifka:2003a} for attempts to define this difference, and suggestions about how to use it to explain accessibility to the double object frame.}%
% //todo Put Pinker's ballistic motion verb lists here.
To the extent that this observation holds, it suggests that there is more to the meaning of the verbs that participate in the alternation than just their argument structure. To account for facts like these, Larson's set of rules would have to be made sensitive to information that goes beyond the $\theta$-roles associated with a verb.

Richard Oehrle\cite{Oehrle:1976} argues against an account like Larson's that relies on putting the alternation in the set of rules that make derivations. Instead he suggests that the process belongs to rules that determine how the pieces that make up the denotation of a verb are mapped to vocabulary items. His evidence consists of observations about how the two frames differ semantically in ways that suggest that they do \uline{not} have the same argument structure. To the extent that the argument structure of a verb is given by the shape of the VP hosting that verb, then these facts speak against giving both frames the same D-structure. He notes, for instance, that the $\theta$-role assigned to the subject of \emph{give} seems to differ depending on how its internal arguments are arranged. For instance, consider the pair of sentences in~(\ref{ex5:mailer}).
\begin{exe}
	\raggedright
	\ex \label{ex5:mailer} 
	\begin{xlist}
		\ex Nixon gave a book to Mailer.
		\ex Nixon gave Mailer a book.
	\end{xlist}
	\hspace*{\fill} \citep[section 1.1, p.~19]{Oehrle:1976} \hspace*{.5in}
\end{exe}
Oehrle notes that (\ref{ex5:mailer}b) allows for an interpretation in which Nixon's behavior is responsible for Mailer writing a book. (\ref{ex5:mailer}a), by contrast, cannot have that interpretation; instead, Nixon must be responsible for Mailer coming into possession of a book. This difference could be partly modeled by letting the external $\theta$-roles be different. The \emph{give} in (\ref{ex5:mailer}a) might assign an \textsc{agent} $\theta$-role, whereas the \emph{give} in (\ref{ex5:mailer}b) might assign a \textsc{cause} $\theta$-role. This shift in the external $\theta$-role doesn't accompany the alternation with every verb. The external $\theta$-role for \emph{bake}, \emph{send}, etc. remains \textsc{agent} in both frames. 

For some verbs, there is a shift in the internal $\theta$-roles assigned by the verb. Most famous, perhaps, are cases that involve certain kinds of motion verbs, including those in~(\ref{ex5:motion1}).
\begin{exe}
	\raggedright
	\ex \label{ex5:motion1} 
	\begin{xlist}
		\ex Jon Ander sent the report to Aleksei.
		\ex Jon Ander kicked the report to Aleksei.
		\ex Jon Ander threw the report to Aleksei.
		\ex Jon Ander mailed the report to Aleksei.
		\ex Jon Ander tossed the report to Aleksei.
	\end{xlist}
\end{exe}
The \emph{to}-phrase in these examples refers to a location, and so we can surmise that the $\theta$-role assigned to this phrase calls for something with the denotation of a location. Let's call this $\theta$-role \textsc{location}. We can see that the \emph{to}-phrases refer to a location because these verbs allow this PP to be expressed with the locative pro-forms \emph{there} and \emph{where}.
\begin{exe}
	\raggedright
	\ex \label{ex5:motion2} 
	\begin{xlist}
		\ex Jon Ander sent the report there.
		\ex Jon Ander kicked the report there.
		\ex Jon Ander threw the report there.
		\ex Jon Ander mailed the report there.
		\ex Jon Ander tossed the report there.
	\end{xlist}
	\ex \label{ex5:motion3}
	\begin{xlist}
		\ex Who sent the report where?
		\ex Who kicked the report where?		
		\ex Who threw the report where?
		\ex Who mailed the report where?
		\ex Who tossed the report where?
	\end{xlist}
\end{exe}
In the double object frame, however, the argument that corresponds to the \emph{to}-phrase is not a locative, as can be seen from the ungrammaticality of~(\ref{ex5:motion4}).
\begin{exe}
	\raggedright
	\ex \label{ex5:motion4} 
	\begin{xlist}
		\ex[*]{Jon Ander sent there the report.}
		\ex[*]{Who sent where the report?}
		\ex[*]{Jon Ander kicked there the report.}
		\ex[*]{Who kicked where the report?}
		\ex[*]{Jon Ander threw there the report.}
		\ex[*]{Who threw where the report?}
		\ex[*]{Jon Ander mailed there the report.}
		\ex[*]{Who mailed where the report?}
		\ex[*]{Jon Ander tossed there the report.}
		\ex[*]{Who tossed where the report?}
	\end{xlist}
\end{exe}
The examples in (\ref{ex5:motion4}) are grammatical if \emph{there} or \emph{where} is replaced by a pronoun that refers to an individual, like \emph{him}. One way of describing the shift in meaning illustrated here is that the non-theme internal argument is a possessor, or perhaps goal. In saying \emph{Jon Ander sent Aleksei the report}, one communicates that the sending eventuality was intended to result in Aleksei's having the report.

Indeed, as Georgia Green\cite[]{Green:1974}, and many following her including Oehrle, have noted, the double object constructions frequently lend themselves to this kind of paraphrase. This is, as noted above, a way of describing the difference between \emph{bake X for Y} and \emph{bake Y X} for example. In each of the cases in (\ref{ex5:have1}), a suitable paraphrase of the meanings is (\ref{ex5:have0}), where X is the subject, Y the second object, V the verb and Z the first object.%
%
\marginnote[0pt]{The idea that the meaning of double object constructions involves ``goals,'' in the way described by (\ref{ex5:have0}), comes from \citet{Koenig:2001}. See \citet{Beavers:2011a} for discussion.}%
%
\begin{exe}
	\raggedright
	\ex \label{ex5:have0} X V Y with the goal of Z having Y.
	\ex \label{ex5:have1}
	\begin{xlist}
		\ex Hotze sewed Jose a handkerchief.
		\ex Pius ironed Hotze a shirt.
		\ex Pat cooked Pius a biscuit.
		\ex Katy faxed Bernhard the money.
		\ex Bernhard painted Katy a picture.
		\ex Pius took Janina the gift.
		\ex Pat left Bernhard the report.
		\ex Jose brought Hotze a chalice.
		\ex Pius guaranteed Pat a position.
		\ex Jose granted Janina a pass.
		\ex Katy drew Jose a tesseract.
		\ex Hotze bought Jose a ring.
		\ex Janina advanced Jose a portion of his pay.
		\ex Jose promised Pat the answer.
	\end{xlist}
\end{exe}
The sense of \emph{have} in these paraphrases is difficult to nail down, as it is sometimes not ``possession'' in the normal use of that verb. This is clear in examples like~(\ref{ex5:have2}).
\begin{exe}
	\raggedright
	\ex \label{ex5:have2} 
	\begin{xlist}
		\ex Katy read Bernhard the report.
		\ex Janina sang Hotze a song.
		\ex Janina taught Jose a madrigal.
		\ex Pat assigned Bernhard the hardest problem.
		\ex Jose quoted Katy the price.
		\ex Janina showed Pius the answer.
	\end{xlist}
\end{exe}
What (\ref{ex5:have2}a) claims Bernhard possesses as a result of Katy's reading the report is not the report, but its contents. Similarly, what (\ref{ex5:have2}e) claims Katy has as a result of Jose's quote is not the price, but knowledge of the price. For discussion, see, in addition to \citet{Green:1974} and \citet[]{Oehrle:1976}: \citet{Kayne:1984d}, \citet{Harves:2012}, \citet{Manzini:2011}, \citet{Manzini:2016}, \citet{Harley:2004}, \citet{Harley:2015}, \citet{Beavers:2011a}, and especially \citet{Beavers:2009}. Because locations cannot ``possess'' things (\emph{*There has a ball}), the ungrammaticality of (\ref{ex5:motion4}) will follow if the argument structure of the double object construction requires the ``have'' predicate in~(\ref{ex5:have0}) to relate the two objects.

Let's express Green's thesis, then, as the claim that all double object constructions have an argument structure that can be expressed with the formula~(\ref{ex5:greenform}).\marginnote[0pt]{VPs that describe eventualities that have a goal are said to be ``telic,'' and telic events can be identified by the ways in which they interact with certain modifiers. See \citet{Vendler:1967}, and \citet{Dowty:1979} who is a good overall guide to the relevant phenomena. \citet{Jackendoff:1996} claims that double object constructions have the properties of telic events.}
\begin{exe}
	\raggedright
	\ex \label{ex5:greenform} 
	{}[Subject verb \nth{2} object ] $\mathcal{R}$ [\nth{1} object have \nth{2} object]
\end{exe}
For \emph{buy}, \emph{throw}, \emph{mail}, etc., $\mathcal{R}$ is something like ``with the goal of.'' The examples in (\ref{ex5:give1}) illustrate that $\mathcal{R}$ can also be something like ``cause.''
\begin{exe}
	\raggedright
	\ex \label{ex5:give1} 
	\begin{xlist}
		\ex Meredith lent Andrew the car.
		\ex Della fed Andrew dinner.
		\ex Andrew served Della lunch.
		\ex Meredith rented Della the car.
		\ex Andrew awarded Meredith the prize.
		\ex Della sold Meredith the car.
		\ex Andrew gave Della the picture.
	\end{xlist}
\end{exe}
What distinguishes these cases from (\ref{ex5:have1}) and (\ref{ex5:have2}) is that the possession relation between the two objects in (\ref{ex5:give1}) is entailed. This can be seen by comparing (\ref{ex5:entail1}) with (\ref{ex5:entail2}).%
%
\marginnote[0pt]{Which verbs support this entailment in the double object construction and which don't is not always clear. Green classed \emph{teach}, for instance, among those that do, but Oehrle argues otherwise.}%
% //todo Find reference to Oehrle for meaning of "teach"
\begin{exe}
	\raggedright
	\ex \label{ex5:entail1} 
	\begin{xlist}
		\ex Pius ironed Hotze a shirt, but it got destroyed before Hotze got home.
		\ex Katy drew Jose a tesseract, but he never saw it.
		\ex Pat left Bernhard the report, but he never found it.
		\ex Jon Ander tossed Della the ball, but it got intercepted.
	\end{xlist}
	\ex \label{ex5:entail2}\marginnote[0pt]{The ``\#'' sign indicates infelicity, or semantic ill-formedness.}
	\begin{xlist}
		\ex[\#]{Meredith lent Andrew the car, but he never got it.}
		\ex[\#]{Della fed Andrew dinner, but he didn't eat it.}
		\ex[\#]{Della sold Meredith the car, but she didn't own it.}
		\ex[\#]{Andrew gave Della the picture, but she never got it.}
	\end{xlist}
\end{exe}
\citet{Hovav:2008} point out that this entailment also holds when these verbs are used in the other frame.\enlargethispage*{20pt}
\begin{exe}
	\raggedright
	\ex \label{ex5:entail3} 
	\begin{xlist}
		\ex[\#]{Meredith lent the car to Andrew, but he never got it.}
		\ex[\#]{Della fed dinner to Andrew, but he didn't eat it.}
		\ex[\#]{Della sold the car to Meredith, but she didn't own it.}
		\ex[\#]{Andrew gave the picture to Della, but she never got it.}
	\end{xlist}	
\end{exe}
What the $\mathcal{R}$ is in the double object construction varies with the root, it appears.

Even (\ref{ex5:deny}) might be able to be fit to this model if one of the ways of combining the two pieces of (\ref{ex5:have0}) can include negation.
\begin{exe}
	\raggedright
	\exr{ex5:deny} 
	Veena denied Rejean his paycheck.
\end{exe}
An awkward, but not obviously incorrect, paraphrase of the meaning of (\ref{ex5:deny}) is~(\ref{ex5:denyp}).
\begin{exe}
	\raggedright
	\ex \label{ex5:denyp} 
	Veena denied his paycheck and this caused not Rejean have his paycheck
\end{exe} 
Note that \emph{deny} belongs to the class of verbs that support the entailment that we see in~(\ref{ex5:entail2}).
\begin{exe}
	\raggedright
	\ex[\#]{\label{ex5:denyentail} Veena denied Rejean his paycheck, but he got it anyway.}
\end{exe}
So we should build into its argument structure the \emph{cause} component that is responsible for that entailment. And finally, we might capture the meaning of (\ref{ex5:envy}) with an argument structure that corresponds to the paraphrase in~(\ref{ex5:envyp}).
\begin{exe}
	\raggedright
	\exr{ex5:envy} Veena envied Rejean his quick footwork.
	\ex \label{ex5:envyp} {} [Rejean has his quick footwork] causes Veena envy
\end{exe}
On this view, the possession of quick footwork by Rejean is the cause of Veena's state of envy. The causation relation reverses the two pieces of the meaning.

Green suggested that all versions of $\mathcal{R}$ included a \emph{cause} component. She didn't countenance the goal oriented meaning that (\ref{ex5:have0}) uses. Capturing the differing entailment facts of (\ref{ex5:entail1}) and (\ref{ex5:entail2}) could be done by letting the $\mathcal{R}$ for the verbs that don't support the entailment have a meaning paraphrased by~(\ref{ex5:prog}).%
%
\marginnote[0pt]{See \citet{Beck:2004} for this view.}%
% //todo Check if Dowty has inertia worlds view of double object
\begin{exe}
	\raggedright
	\ex \label{ex5:prog} 
	$\mathcal{R}$ = cause becoming
\end{exe}
On this view, the meaning of (\ref{ex5:throwp}) matches the (terrible) paraphrase in~(\ref{ex5:throwpd}).
\begin{exe}
	\raggedright
	\ex \label{ex5:throwp} 
	Ania threw Stephen the ball.
	\ex \label{ex5:throwpd}
	Ania threw the ball which caused [Stephen to have the ball] becoming.
\end{exe}
On this view, the eventuality of throwing the ball causes another eventuality which, if things progress normally, leads to Stephen having the ball. If $\mathcal{R}$ is restricted to flavors of \emph{cause}, then this might reflect the fact that the double object construction is built upon an argument structure that necessarily includes a ``cause X have Y'' ingredient.

Richard Kayne\cite{Kayne:1984d} has argued that the meaning component associated with the \emph{have} part of our paraphrases is found in the syntax of the double object construction. On his view the D-structure of a double object construction has the shape in~(\ref{ex5:kayne1}).
\begin{exe}
	\raggedright
	\ex \label{ex5:kayne1} 
	\small
	\begin{forest}
	[vP
		[DP [Ania, roof] ]
		[\ibar{v}
			[\obar{v} [agent] ]
			[VP [\ibar{V}
				[\obar{V} [$\surd$throw] ]
				[VP
					[DP [Stephen, roof]]
					[\ibar{V}
						[\obar{V} [\shade{have}] ]
						[DP [the ball, roof] ]
					 ]
				 ]
			] ]
		 ]
	]
	\end{forest}
	\marginnote[0pt]{I am fitting Kayne's proposal into a syntax which gives the external $\theta$-role the position~\obar{v}.}
\end{exe}
Kayne's argument involves comparing the argument structure of verbs with the argument structure of nouns morphologically related to those verbs. I postpone looking at his argument until the next section.%
%
\marginnote[0pt]{\citet{Beavers:2009} credits \citet{McCawley:1974b} with a similar proposal about making the \emph{have} component part of the syntax of double object constructions.}%
%
\marginnote[0pt]{\citet{Harley:1996} argues that only languages in which the possession relation is lexicalized into a verb like English's \emph{have} are capable of hosting the double object frame. In her (small) sample of languages, there was a correlation between the languages that permit the double object construction and those that use a verb to express possession. (A common alternative way of expressing possession is with the copula (e.g., English \emph{be}) and a goal preposition or case.) She suggests this supports Kayne's proposal to the extent that the structure in (\ref{ex5:kayne1}) requires an \obar{X} that can express possession.}%
%

If (\ref{ex5:kayne1}) is correct, there is a phrase within the double object construction that describes the possession relationship between the two objects. We should be able to detect the presence of that phrase using \emph{again}. Sigrid Beck\cite{Beck:2004} has argued that this expectation is fulfilled. To see this, consider the contrast between~(\ref{ex5:doagain1}).
\begin{exe}
	\raggedright
	\ex \emph{Cecilia lost her job and for the first time in her life, became despondent about her future. After he was introduced to Cecilia at the meeting, \ldots}
	\ex \label{ex5:doagain1} 
	\begin{xlist}
		\ex[]{Marcin gave Cecilia hope again during the discussion.}
		\ex[\#]{Marcin gave Cecilia hope during the discussion again.}
	\end{xlist}
\end{exe}
The word order in (\ref{ex5:doagain1}a) permits the parse in (\ref{ex5:doagain1pa}), but the word order in (\ref{ex5:doagain1}b) forces a parse like that in~(\ref{ex5:doagain1pb}).
\begin{exe}
	\raggedright
	\ex \label{ex5:doagain1pa} 
	\small
	\begin{forest}
	[\ibar{V}
		[\ibar{V} 
			[\obar{V} [$\surd$give] ]
			[VP
				[VP 
					[DP [Cecilia, roof] ]
					[\ibar{V}
						[\obar{V} [\shade{have}] ]
						[DP [hope, roof] ]
					 ]
				 ]
				[AdvP [again] ]
			 ]
		]
		[PP [during the discussion, roof] ]
	]
	\end{forest}
	\ex \label{ex5:doagain1pb}
	\begin{forest}
	[\ibar{V}
		[\ibar{V} 
			[\obar{V} 
				[\obar{V} [$\surd$give] ]
				[VP
					[DP [Cecilia, roof] ]
					[\ibar{V}
						[\obar{V} [\shade{have}] ]
						[DP [hope, roof] ]
					 ]
				 ]
			 ]
			[PP [during the discussion, roof] ]
		]
		[AdvP [again, roof] ]
	]
	\end{forest}
\end{exe}
In the context provided, there is no prior eventuality of giving Ceclia hope, and that is the meaning that would be required of (\ref{ex5:doagain1pb}). Instead, the context supports the inference that Cecilia was in a prior state of hopefulness, and that is the meaning our semantics for \emph{again} would deliver to~(\ref{ex5:doagain1pa}). This is just what is expected if there is the \emph{have}-VP within the double-object construction that Kayne proposes.

We can see the same confirmation by considering how certain duration modifiers work. For instance, the meaning of (\ref{ex5:lentf}) has a meaning in which \emph{for the week} specifies how long the possession state lasts.%
%
\marginnote[0pt]{See \citet{Beavers:2009}, and references cited therein, for this argument.}%
%
\begin{exe}
	\raggedright
	\ex \label{ex5:lentf} 
	She lent me the lawnmower for the week this morning.
\end{exe}
Note that the two temporal modifiers do not lead to a conflict in (\ref{ex5:lentf}), like they do in~(\ref{ex5:nlentf}).
\begin{exe}
	\raggedright
	\ex[\#]{\label{ex5:nlentf} She hugged me for the week this morning.}
\end{exe}
In (\ref{ex5:nlentf}), both temporal expressions must describe the time at which the one eventuality of \emph{hugging} transpires, leading to the conflict. But in (\ref{ex5:lentf}), by contrast, there are two descriptions of eventualities that the temporal modifiers can combine with, as shown in~(\ref{ex5:lentfp}).
\begin{exe}
	\raggedright
	\ex \label{ex5:lentfp} 
	\small
\begin{forest}
[\ibar{V}
	[\ibar{V} 
		[\obar{V} [$\surd$lend] ]
		[VP 
			[VP 
				[DP [me, roof] ]
				[\ibar{V}
					[\obar{V} [\shade{have}] ]
					[DP [the lawnmower, roof] ]
				 ]
			 ]
			[PP [for the week, roof] ]
		 ]
	]
	[DP [this morning, roof] ]
]
\end{forest}
\end{exe}
Unlike \emph{again}, however, the word order doesn't correlate with the availability of this reading.
\begin{exe}
	\raggedright
	\ex \label{ex5:lentf2} 
	She lent me the lawnmower this morning for the week.
\end{exe}
Perhaps this is because, unlike \emph{again}, \emph{for the week} can undergo Complex NP Shift.
\begin{exe}
	\raggedright
	\ex \label{ex5:hnpslend} 
	\small
	\begin{forest}
	[TP
		[TP 
			[DP [she, name=binders, roof]]
			[\ibar{T}
				[\obar{T} [past] ]
				[vP
					[{\ }, name=traces]
					[\ibar{v}
						[\obar{v} [agent] ]
						[VP [\ibar{V}
							[\obar{V}
								[\obar{V} [$\surd$lend] ]
								[VP
									[VP
										[DP [me, roof]me ]
										[\ibar{V} 
											[\obar{V} [\shade{have}] ]
											[DP [the lawnmower, roof] ]
										 ]
									 ]
									[{\ }, name=trace ]
								 ]
							 ]
							[DP [this morning, roof] ]
						] ]
					 ]
				 ]
			 ]
		]
		[PP [for the week, name=binder, roof] ]
	]
	\draw[->,>=stealth', ultra thin, gray] (trace) to[out=east,in=south east,looseness=2] node[pos=0.75,xshift=10mm]{} (binder);
	\draw[->,>=stealth', ultra thin, gray] (traces) to[out=south west,in=south,looseness=1.3] node[pos=0.75,xshift=-10mm]{} (binders);
	\end{forest}
\end{exe}

Kayne's syntax does not combine with our semantics to deliver the right meaning. If the meaning of the root is a simple description of eventualities, then our present semantics doesn't include a way of putting the pieces together.
\medskip
\begin{fullwidth}
\begin{exe}
	\raggedright
	\ex \label{ex5:nosemill} 
	\small
	\begin{forest}
	[??\\ \ibar{V}
		[$\lambda \epsilon$ \textsc{lend}($\epsilon$)\\ \obar{V} [$\surd$lend]]
		[$\lambda \epsilon$ \textsc{have}($\epsilon$) $\wedge$ \textsc{theme}(the\_lawmower)($\epsilon$) $\wedge$ \textsc{goal}(me)($\epsilon$)\\ VP
			[DP [me, roof] ]
			[$\lambda x \lambda \epsilon$ \textsc{have}($\epsilon$) $\wedge$ \textsc{theme}(the\_lawnmower)($\epsilon$) $\wedge$ \textsc{goal}($x$)($\epsilon$)\\\ibar{V} 
				[$\lambda y \lambda x \lambda \epsilon$ \textsc{have}($\epsilon$) $\wedge$ \textsc{theme}($y$)($\epsilon$) $\wedge$ \textsc{goal}($x$)($\epsilon$)\\\obar{V} [\shade{have}] ]
				[DP [the lawnmower, roof] ]
			]
		 ]
	]
	\end{forest}
\end{exe}
\end{fullwidth}
\medskip
What is missing here, of course, is the $\mathcal{R}$ that connects the lending eventuality with the having eventuality. As we've seen, the content of $\mathcal{R}$ varies with the lexical item chosen. How these two pieces are semantically composed should therefore be part of the information that comes with the lexical item. One way of doing that would be to build into the denotation of the root this meaning. See \citet{Beavers:2017,Beavers:2020} for that direction. Alternatively, there could be another syntactic piece to the VP which, when assembled with the others, is matched to the relevant lexical item. For \emph{lend}, that could be achieved with the D-structure in~(\ref{ex5:causelend}).
\medskip
\begin{fullwidth}
\begin{exe}
	\raggedright
	\ex \label{ex5:causelend} 
	\small
	\begin{forest}
	[$\lambda \epsilon \exists \epsilon'$ \textsc{lend}($\epsilon$) $\wedge$ {[}\textsc{have}($\epsilon'$) $\wedge$\textsc{theme}(the\_mower)($\epsilon'$) $\wedge$ \textsc{goas}(me)($\epsilon'$) $\wedge$ \textsc{cause}($\epsilon'$)($\epsilon$){]}\\ \ibar{V}
		[$\lambda \epsilon$ \textsc{lend}($\epsilon$)\\ \obar{V} [$\surd$lend] ]
		[$\lambda Q$ $\lambda \epsilon$ $\exists \epsilon'$ $Q$($\epsilon$) $\wedge$ {[}\textsc{have}($\epsilon'$) $\wedge$ \textsc{theme}(the\_mower)($\epsilon'$) $\wedge$ \textsc{goal}(me)($\epsilon'$){]}\\ $\wedge$ \textsc{cause}($\epsilon'$)($\epsilon$)\\ VP [\ibar{V}
			[$\lambda P$ $\lambda Q$ $\lambda \epsilon$ $\exists \epsilon'$ $Q$($\epsilon$) $\wedge$ $P$($\epsilon'$) $\wedge$ \textsc{cause}($\epsilon'$)($\epsilon$)\\ \obar{V} [\shade{cause}]]
			[$\lambda \epsilon$ \textsc{have}($\epsilon$) $\wedge$ \textsc{theme}(the\_lawnmower)($\epsilon$) $\wedge$ \textsc{goal}(me)($\epsilon$)\\ VP 
				[DP [me, roof] ]
				[\ibar{V}
					[\obar{V} [\shade{have}] ]
					[DP [the lawnmower, roof] ]
				 ]
			]
		] ]
	]
	\end{forest}
\end{exe}
\end{fullwidth}
\medskip
Other verbs---\emph{deny}, \emph{toss}---could express the alternative $\mathcal{R}$s put into the \obar{V} position that \emph{cause} occupies in~(\ref{ex5:causelend}). 

What these implementations of Green's thesis miss is how to capture the fact that the second object of the double object construction is an object of the root as well. That seems semantically to be the case, as we can see from~(\ref{ex5:objsem}).
\begin{exe}
	\ex \label{ex5:objsem} 
\begin{xlist}
	\ex[\#]{Laura melted Amalia a puddle.}
	\ex[]{Laura melted Amalia the ice.}
\end{xlist}
\end{exe}
The anomaly of (\ref{ex5:objsem}a) is the same as the anomaly of \emph{Laura melted a puddle}: puddles cannot be the patient of a melting eventuality.

Note also that obligatorily transitive verbs can fit into the double object construction. 
\enlargethispage*{\baselineskip}
\begin{exe}
	\ex \label{ex5:transdo} 
	\begin{xlist}
		\ex[]{Tohru sent (Suzanne) the letter yesterday.}
		\sn[*]{Tohru sent yesterday.}
		\ex[]{Tohru rolled (me) the marbles yesterday.}
		\sn[*]{Tohru rolled yesterday.}
		\ex[]{Tohru brought (me) the marbles yesterday.}
		\ex[*]{Tohru brought yesterday.}
		\ex[]{Tohru mixed (me) a drink yesterday.}
		\ex[*]{Tohru mixed yesterday.}
		\ex[]{Tohru bought (me) a new map yesterday.}
		\ex[*]{Tohru bought yesterday.}
	\end{xlist}
\end{exe}
And no obligatorily intransitive ones can, even though intransitives can sometimes figure in resultative constructions which have many of the same features of the syntax and semantics of~(\ref{ex5:causelend}), as (\ref{ex5:intransitiveres}) illustrates.%
%
\marginnote[0pt]{See \citet{Hoekstra:1988} and \citet{Kratzer:2005} for a syntax and semantics of this kind for resultatives.}%
%
\begin{exe}
	\ex \label{ex5:intransitiveres} Suzanne danced herself silly.
	\sn
	\small
	\begin{forest}
	[$\lambda \epsilon$ $\exists \epsilon'$ \textsc{dance}($\epsilon$) $\wedge$ \textsc{silly}($\epsilon'$) $\wedge$ \textsc{holder}(herself)($\epsilon'$) $\wedge$ \textsc{cause}($\epsilon'$)($\epsilon$) \\\ibar{V}
		[$\lambda \epsilon$ \textsc{dance}($\epsilon$)\\ \obar{V} [$\surd$dance]]
		[$\lambda Q$ $\lambda \epsilon$ $\exists \epsilon'$ $Q$($\epsilon$) $\wedge$ \textsc{silly}($\epsilon'$) $\wedge$ \textsc{holder}(herself)($\epsilon'$) $\wedge$ \textsc{cause}($\epsilon'$)($\epsilon$)\\ VP [\ibar{V}
			[$\lambda P$ $\lambda Q$ $\lambda \epsilon$ $\exists \epsilon'$ $Q$($\epsilon$) $\wedge$ $P$($\epsilon'$)$\wedge$ \textsc{cause}($\epsilon'$)($\epsilon$) \\ \obar{V} [\shade{cause}]]
			[$\lambda \epsilon$ \textsc{silly}($\epsilon$) $\wedge$ \textsc{holder}(herself)($\epsilon$)\\ AP
				[DP [herself, roof] ]
				[\ibar{A} [silly, roof] ]
			 ]
		]]
	]
	\end{forest}
\end{exe}
(\ref{ex5:causelend}) would lead us to expect that something similar would be possible with the double object construction. But all obligatorily intransitive verbs are, like~(\ref{ex5:intransitivedo}), ungrammatical in the double object construction.
\begin{exe}
	\raggedright
	\ex[*]{\label{ex5:intransitivedo} Suzanne danced herself two blisters. }
\end{exe}
This all makes sense only if the second object is a genuine object of the verbal root.

This has led \citet{Bruening:2010a,Bruening:2010}, among others, to a syntax like~(\ref{ex5:toolow}).\marginnote[0pt]{\citet{Marantz:1993} might be the first to propose this structure. See also \citet{Pylkkanen:2008}.}
\begin{exe}
	\ex \label{ex5:toolow} 
	\small
	\begin{forest}
	[\ibar{v}
		[\obar{v} 
			[\obar{V} [throw, name=binder] ]
			[\obar{v} [agent] ]
		 ]
		[VP
			[DP [Stephen, roof] ]
			[\ibar{V}
				[\obar{V} [<appl>] ]
				[VP [\ibar{V}
					[{\ }, name=trace ]
					[DP [the ball, roof] ]
				] ]
			] 
		] 
	]
	\draw[->,>=stealth', ultra thin, gray] (trace) to[out=south west,in=south,looseness=1] node[pos=0.75,xshift=-10mm]{} (binder);
	\end{forest}
\end{exe}
Here the <appl> head must somehow encode what I've expressed as \textsc{have}. But this syntax does not allow for there to be a constituent that is made up of the two objects and \textsc{have} (or its equivalent) that does not include the root. That is what we seem to require to get the facts about temporal modifiers and \emph{again}. What seems to be required is for the second object to be both an argument of \textsc{have} and, independently, an argument of the root. Our syntax does not presently provide a means for doing that.

If we set this problem aside for a moment, we can now ask how to capture the relationship between the two ways of expressing arguments of the verbs that participate in the alternation. If the argument structures of the two variants are not (always) identical, an approach that credits the regular relationship between them to rules that make a derivation seems wrong. As Oehrle concluded, what we look for is a process that relates different D-structures. We need a process that says that if a vocabulary item expresses such-and-such an argument structure, then it will express such-and-such other argument structure too. \citet{Harley:2002a} suggests that the two D-structure representations we want to relate to a motion verb, like \emph{throw}, are~(\ref{ex5:harleya}).%
%
\marginnote[-20pt]{See \citet{Harley:2012} and \citet{Harley:2015} for some modifications. I change the particulars of her proposal to give a place for $\mathcal{R}$, and fit our constraints on phrase markers.}%
%
\begin{fullwidth}
\medskip
\begin{exe}
	\raggedright
	\ex \label{ex5:harleya} 
	\small
	\begin{forest}
	[vP
		[DP [Ania, roof]]
		[\ibar{v} 
			[\obar{v} [agent] ]
			[VP [\ibar{V}
				[\obar{V} [$\surd$throw] ]
				[VP [\ibar{V}
					[\obar{V} [$\mathcal{R}$] ]
					[VP
						[DP [Stephen, roof] ]
						[\ibar{V}
							[\obar{V} [\shade{have}] ]
							[DP [the ball, roof] ]
						 ]
					 ]
				] ]
			] ]
		]
	]
	\end{forest}
	\hspace*{-22pt}
	\begin{forest}
	[vP
		[DP [Ania, roof]]
		[\ibar{v} 
			[\obar{v} [agent] ]
			[VP [\ibar{V}
			[\obar{V} [$\surd$throw]]
			[VP [\ibar{V}
				[\obar{V} [$\mathcal{R}$] ]
				[PP
					[DP [the ball, roof] ]
					[\ibar{P}
						[\obar{P} [to] ]
						[DP [Stephen, roof] ]
					 ]
				 ]
			] ]
		]
	]]]
	\end{forest}
\end{exe}
\medskip
\end{fullwidth}
Recall that the $\mathcal{R}$ which accompanies \emph{throw} is the ``cause becoming'' or ``with the goal of'' sort; it does not entail that the state denoted by its complement comes to exist. And recall that the \emph{to} which shows up with \emph{throw} is a locative; its meaning can be paraphrased by ``is at the location of.'' On Harley's view, then, both instantiations of \emph{throw} involve a root that describes a throwing eventuality and a state that is the target result of that eventuality. In the case of the double object construction, that state is possession (i.e., \emph{have}), and in the PP frame it is a thing residing in a location (i.e. \emph{to}). That there is a phrase that denotes a locative state in the PP frame is confirmed by \emph{again}.
\begin{exe}
	\raggedright
	\sn \emph{The coach put the balls in the endzone before the game. When one of the balls rolled into the field,\ldots}
	\ex \label{ex5:againto} 
	The waterboy threw [[it to the endzone] again].
\end{exe}
In the context provided for (\ref{ex5:againto}), there is no previous throwing eventuality, and so \emph{again} cannot be combining with the VP headed by \emph{throw}. Instead, \emph{again} attaches to the PP headed by \emph{to}, conveying that the state of the ball being in the endzone is being restored.
\enlargethispage*{\baselineskip}

If these were the correct D-structures for the two frames, then the relationship between them could be expressed with a rule like~(\ref{ex5:lexrule}).
\begin{exe}
	\raggedright
	\ex \label{ex5:lexrule} 
	If $\alpha$ is a vocabulary item that expones $\surd \alpha$+$\mathcal{R}+\emph{have}$,\\ then $\alpha$ also expones $\surd \alpha$+$\mathcal{R}$+\emph{to}
\end{exe}
What an English speaker has done is learn which verbs can be $\alpha$.

The picture of the alternation between the PP and double object frames given here harbors several problems. In addition to the problem of understanding how the second object of the double object frame is an argument of the root,%
%
\marginnote[-20pt]{Note that (\ref{ex5:harleya}) also (wrongly) makes the first object of the PP frame not an argument of the root.}%
%
there is also the problem of understanding how rules like (\ref{ex5:lexrule}) can be made part of a natural class of processes, rather than simple \emph{ad hoc} descriptions of what is needed. Nonetheless, if something like Kayne's proposal for the double object construction can be maintained, it does provide an account of how two structural cases arise in these constructions. Consider how our system of accusative Case assignment works in a Kaynian model of the double object construction.
\begin{exe}
	\raggedright
	\ex \label{ex5:casekayne} 
	\small
	\begin{forest}
	[vP
		[DP [Veena, roof]]
		[\ibar{v} 
			[\obar{v} [agent] ]
			[$\mu$P [\ibar{$\mu$}
				[\obar{$\mu$} [case: \emph{acc}] ]
				[VP [\ibar{V}
					[\obar{V} [$\surd$give] ]
					[VP [\ibar{V}
						[\obar{V} [\shade{cause}] ]
						[VP
							[DP [Rejean, roof] ]
							[\ibar{V}
								[\obar{V} [\shade{have}] ]
								[DP [the answer, roof] ]
							 ]
						 ]
					] ]
				] ]
			] ]
		]
	]
	\end{forest}
\end{exe}
Suppose that \emph{have} comes with an accusative Case feature. Perhaps this is telling us that the category of \emph{have} is actually \obar{P},%
%
\marginnote[0pt]{See \citet{Kayne:1993} for the idea that \emph{have} is a preposition in the double object construction. Occasionally, the idea that what I've represented with ``\emph{have}'' would be lexicalized as the preposition \emph{with} is suggested. Expressions like \emph{the woman with a hat/all the answers} involve a meaning for \emph{with} that is close to what's needed.}% 
%
a category that English otherwise imbues with an accusative Case feature. The second DP is in a position that allows its unvalued Case feature to be valued by the accusative Case feature on \emph{have}. But it is too distant to have its Case feature valued by $\mu$, since there is a another DP---the first object---between them. The first object's unvalued Case feature can get valued by the accusative feature in $\mu$, however. Since it is the Case feature that is found in $\mu$ that is redirected in the Passive, we have an account for why in (some dialects of) English, the first DP, but not the second, converts to Nominative when the double object verb is passivized.

Kayne's thesis that the double object construction houses a small clause aids in capturing the facts about \emph{again} and it sheds light on the availability of two accusative Cases in these constructions. His argument for the small clause, however, involves a deduction he makes from the fact that the double object construction does not show up in the NPs headed by nouns derived from verbs. Let's now examine this argument.

\section{Argument Structure of Nominalizations} 
\label{sec:nominalizations}

To appreciate Kayne's argument, we'll first have to learn a few things about the argument structure of nouns. We'll focus our discussion on nouns that are morphologically related to verbs. If we imagine that the nouns and verbs which are morphologically related share a root---that is are built from the same underlying morphological piece---then what we can learn by studying their different argument structures is something about how nouns and verbs differ. This is what is relevant to Kayne's argument.

\subsection{How the Theta Criterion applies to nominals} 
\label{sub:nominaltheta}

When nouns are related to verbs, the arguments they can combine with are often parallel to those found with the verbs. Rather than nominative Case, the external arguments of nouns appear in the genitive Case, and in the Specifier of DP position. Some examples are in~(\ref{ex5:nomin1}).
\begin{exe}
	\raggedright
	\ex \label{ex5:nomin1} 
	\begin{xlist}
		\ex Presley's fear of marigolds
		\ex Meg's dislike of MaxElide
		\ex Noah's run to the store
		\ex Kathryn's tumble down the stairs
		\ex Michael's discussion of squirrels
		\ex Maribel's observation of the star
		\ex Aynat's advocacy of your proposal
		\ex Martin's toss of the baseball to you
		\ex Luis's exchange of coins with you
		\ex Jan's placement of the book on the table
		\ex Wendell's gift of the book to you
		\ex Cherlon's proposal that we should leave
		\ex Anna's suggestion that natto kills
		\ex Karen's attempt to fold the origami
		\ex Ilaria's promise to stay
	\end{xlist}
\end{exe}
The $\theta$-roles and c-selection requirements of these nouns match those of their cognate verbs, with the exception that there are no DP internal arguments. This flows from the observation that accusative Case is not licensed in nominals. Instead, the argument that would appear as a structurally case-marked object of a verb appears with the preposition \emph{of} in the noun version of that verb. If we import the model we've developed for verbs to nominals, we'll have a syntax like that in~(\ref{ex5:noun1}).
\begin{exe}
	\raggedright
	\ex \label{ex5:noun1} 
	\small
	\begin{forest}
	[DP
		[DP's\\ case: \emph{gen} [Presely, name=binder, roof] ]
		[\ibar{D}
			[\obar{D} [case: \emph{gen}] ]
			[nP
				[{\ }, name=trace ]
				[\ibar{n}
					[\obar{n} [$\theta$-role: \emph{experiencer}	]]
					[NP [\ibar{N}
						[\obar{N} [$\surd$fear\\ $\theta$-role: \emph{cause}\\ $\theta$-role: \emph{experiencer}] ]
						[PP [of marigolds, roof] ]
					] ]
				 ]
			 ]
		 ]
	]
	\draw[->,>=stealth', ultra thin, gray] (trace) to[out=south west,in=south,looseness=1] node[pos=0.75,xshift=-10mm]{} (binder);
	\end{forest}
\end{exe}
The external $\theta$-role resides outside the projection of the noun, just as it does with verbs. I've called the position that the external $\theta$-role resides in here \obar{n}, to designate that this is to be understood as part of a nominal. As with verbs, the external $\theta$-role that resides outside of the projection of the word that expresses that $\theta$-role is in an \textsc{agree} relation with that word. The external argument resides in the Specifier of nP at D-structure, and when that argument is a DP it comes with an unvalued Case feature. That feature is valued by the genitive Case that comes with the determiner heading the DP. This Case must have a $\bullet$ associated with it, as movement of the Case-marked DP into the Specifier of DP is obligatory. 

Interestingly, the external $\theta$-role in nominals is not able to assign accusative Case; if it were, of course, we'd expect to see accusative Case-marked internal arguments of nouns. We could assume that the structural Case which nominalized external $\theta$-roles assign is the one expressed by \emph{of}. But if that were correct, it would require that the \emph{of} which shows up with nouns that do not support an external $\theta$-role, as in (\ref{ex5:noexof}), is different.
\begin{exe}
	\raggedright
	\ex \label{ex5:noexof} 
	\begin{xlist}
		\ex the sister of Amy Rose
		\ex the piece of paper
		\ex the cup of Harissa 
	\end{xlist}
\end{exe}
Indeed, even with the nominals in (\ref{ex5:nomin1}), the external $\theta$-role bearer can be absent and the \emph{of} remains.
\begin{exe}
	\raggedright
	\ex \label{ex5:nomin2} 
	\begin{xlist}
		\ex the fear of marigolds
		\ex a dislike of MaxElide
		\ex a run to the store
		\ex the tumble down the stairs
		\ex a discussion of squirrels
		\ex an observation of the star
		\ex the advocacy of your proposal
		\ex the toss of the baseball to you
		\ex the exchange of coins with you
		\ex the placement of the book on the table
		\ex a gift of the book to you
		\ex the proposal that we should leave
		\ex the suggestion that natto kills
		\ex an attempt to fold the origami
		\ex every promise to stay
	\end{xlist}	
\end{exe}
On the other hand, \emph{of} does not seem like a semantically contentful preposition. The DPs with \emph{of} in (\ref{ex5:nomin2}) seem to combine with the nouns they are complements to in the same way those DPs would to the related verbs. It is the $\theta$-roles that these nouns and verbs assign that is responsible for their semantic composition. The idea that \emph{of} is a kind of Case morpheme, and moreover a semantically contentless one---a ``structural'' Case---seems appropriate. If all of this is correct, then the examples in (\ref{ex5:nomin2}) are situations where an internal argument is getting a structural Case from a term that has no external $\theta$-role. These are violations of Burzio's Generalization. An explanation of Burzio's Generalization should include why it doesn't hold in (English) NPs. For concreteness, I will assume that the term that assigns the structural Case \emph{of} is the noun itself. 

(\ref{ex5:nomin2}) indicates that the arguments of nouns are not as obligatory as they are for verbs. This is true of internal arguments as well, as (\ref{ex5:nomin3}) shows.
\begin{exe}
	\raggedright
	\ex \label{ex5:nomin3} 

	\begin{tabular}{ll}
		a. the fear & i. the exchange\\
		b. a dislike & j. the placement\\
		c. a run & k. a gift of the book to you\\
		d. the tumble & l. the proposal that we should leave\\
		e. a discussion & m. the suggestion that natto kills\\
		f. an observation & n. an attempt to fold the origani\\
		g. the advocacy & o. every promise to stay
	\end{tabular}
\end{exe}
If the examples in (\ref{ex5:nomin1}) suggest that the argument structure of nouns is the same as the argument structure of verbs then there must be something that allows nouns to avoid the effects of Theta Criterion. That is, if the denotations of nouns like \emph{exchange} is roughly the same as that of the verb \emph{exchange}, then we need an explanation for why the NP headed by \emph{exchange} is grammatical without its internal argument, but a VP headed by \emph{exchange} is not.
\begin{exe}
	\raggedright
	\ex \label{ex5:exchange1} 
	\begin{xlist}
		\ex[]{Luis's exchange}
		\ex[*]{Luis exchanged.}
	\end{xlist}
\end{exe}
And we also need an explanation for why the external $\theta$-role for the verb is obligatory, but not for the noun.
\begin{exe}
	\raggedright
	\ex \label{ex5:exchange2} 
	\begin{xlist}
		\ex[]{the exchange (of coins)}
		\ex[*]{It (\emph{expletive}) exchanged.}
	\end{xlist}
\end{exe}

Our account for the Theta Criterion,%
%
\marginnote[0pt]{See the discussion in chapter \ref{sec:reconsidered}.}%
%
relies on the stipulation that vP must be of a particular semantic type, and a combination of denotations assigned to the verb and its $\theta$-roles along with a theory that limits the ways in which terms can be semantically combined. These pieces of the explanation are assembled in (\ref{ex5:tcexpl}), and their application to a vP headed by \emph{exchange} is shown in~(\ref{ex5:exchange3}).
\begin{fullwidth}
\medskip
\begin{exe}
	\raggedright
	\ex \label{ex5:tcexpl} 
	\begin{xlist}
		\ex \denotes{vP} must be a function whose domain is eventualities and range truth values.
		\ex {}[_{\text{\obar{V}}} exchange] $\Rightarrow$ \textsc{agent}+$\surd$exchange
		\ex \denotes{$\surd$exchange} = $\lambda x\ \lambda \epsilon$ \textsc{exchange}($\epsilon$) $\wedge$ \textsc{theme}($x$)($\epsilon$)=T
		\ex \denotes{\textsc{agent}} = $\lambda P\ \lambda x\ \lambda \epsilon$ $P(\epsilon)\ \wedge$ \textsc{agent}($x$)($\epsilon$)=T
	\end{xlist}
	\ex \label{ex5:exchange3}
	\small
	\begin{forest}
	[$\lambda \epsilon$ \textsc{exchange}($\epsilon$) $\wedge$ \textsc{theme}(coins)($\epsilon$) $\wedge$ \textsc{agent}(Luis)($\epsilon$)\\ vP
		[DP [Luis, roof] ]
		[$\lambda x\ \lambda \epsilon$ \textsc{exchange}($\epsilon$) $\wedge$ \textsc{theme}(coins)($\epsilon$) $\wedge$ \textsc{agent}($x$)($\epsilon$)\\ \ibar{v}
			[$\lambda P\ \lambda x\ \lambda \epsilon$ $P(\epsilon)\ \wedge$ \textsc{agent}($x$)($\epsilon$)\\ \obar{v} [agent] ]
			[$\lambda \epsilon$ \textsc{exchange}($\epsilon$) $\wedge$ \textsc{theme}(coins)($\epsilon$)\\ VP [$\lambda \epsilon$ \textsc{exchange}($\epsilon$) $\wedge$ \textsc{theme}(coins)($\epsilon$)=T\\ \ibar{V}
				[$\lambda x\ \lambda \epsilon$ \textsc{exchange}($\epsilon$) $\wedge$ \textsc{theme}($x$)($\epsilon$)\\ \obar{V} [$\surd$exchange] ]
				[DP [coins, roof] ]
			] ]
		 ]
	]
	\end{forest}
\end{exe}
\end{fullwidth}
\medskip
If the DPs \emph{Luis} and \emph{coins} had not been present, it would not have been possible ot use function application---our only semantic rule of composition---to produce from the denotations that \emph{exchange} introduces a vP with the requisite meaning. Something in nominals must interrupt this effect.

Let's focus first on the absence of the external argument in nominals.%
%
\marginnote[0pt]{My discussion here is led by \citet{Moulton:2014}, with some tweaks to fit to our syntax. Moulton builds on \citet{Grimshaw:1990b} and \citet{Borer:2005a}. All of these approaches to the semantics of nominalizations have their roots in \citet{Higginbotham:1983}.}%
%
We have a mechanism that is designed for removing the Theta Criterion's effects by existentially closing the function that otherwise requires an argument. This process is designed so that it only applies to internal arguments, capturing the fact that external arguments of verbs are required whenever the external $\theta$-role associated with that verb is present. However, we also saw that this process is capable of applying to an external $\theta$-role's argument in the Passive. In the Passive, the external argument appears in a PP headed by \emph{by}, or it is existentially closed:
\begin{exe}
	\raggedright
	\ex \label{ex5:passivee} 
	\begin{xlist}
		\ex Coins were exchanged by Luis.
		\ex Coins were exchanged.
	\end{xlist}
\end{exe}

Perhaps the process of existential closure can apply to external $\theta$-roles within nominals more liberally. One way of detecting the presence of an external $\theta$-role in Passives is with adverbs that require them. For instance, the adverb \emph{intentionally} can describe eventualities only if those eventualities have a suitable external $\theta$-role.%
%
\marginnote[0pt]{See \citet{Jackendoff:1972}.}%
%
\begin{exe}
	\raggedright
	\ex \label{ex5:agento} 
	\begin{xlist}
		\ex[]{Luis intentionally exchanged the coins for buttons.}
		\ex[\#]{Luis intentionally seemed happy.}
		\ex[\#]{Luis intentionally existed.}
		\ex[\#]{Coins intentionally equal buttons in value.}
		\ex[\#]{The election intentionally affected us.}
		\ex[\#]{We were intentionally affected by the election.}
	\end{xlist}
\end{exe}
The grammaticality of (\ref{ex5:passiveto}) suggests that there is an external $\theta$-role in it.
\begin{exe}
	\raggedright
	\ex \label{ex5:passiveto} 
	Coins were intentionally exchanged for buttons.
\end{exe}
The same reasoning suggests that an external $\theta$-role remains in (\ref{ex5:exnome}) as well.
\begin{exe}
	\raggedright
	\ex \label{ex5:exnome} 
	The intentional exchange of coins for buttons (was shocking).
\end{exe}
A working hypothesis, then, is that nouns allow for a wider application of existential closure than verbs. This would allow for DPs like~(\ref{ex5:eventn1}).
\begin{fullwidth}
\medskip
\begin{exe}
	\raggedright
	\ex \label{ex5:eventn1} 
	\small
	\begin{forest}
	[DP
		[\obar{D} [the] ]
		[$\lambda \epsilon\ \exists x\ \exists y$ \textsc{exchange}($\epsilon$) $\wedge$ \textsc{theme}($y$)($\epsilon$) $\wedge$ \textsc{agent}($x$)($\epsilon$)\\ nP [\ibar{n}
			[$\lambda P\ \lambda \epsilon\ \exists x$  $P(\epsilon)$ $\wedge$ \textsc{agent}($x$)($\epsilon$)\\ \obar{n} [agent] ]
			[$\lambda \epsilon\ \exists y$ \textsc{exchange}($\epsilon$) $\wedge$ \textsc{theme}($y$)($\epsilon$)\\ NP [\ibar{N} [\obar{N} [exchange]]] ]
		] ]
	]
	\end{forest}
\end{exe}
\end{fullwidth}
\medskip
From this meaning for nP must be derived a denotation that allows the DP to refer to something of the appropriate sort. This is what is achieved by the meaning of the determiner \emph{the}. The denotation of \emph{the} is complex and must allow for various outcomes, it appears. We needn't engage the semantics of \emph{the} fully here. We can instead work with an informal sketch. Roughly, what the denotation of \emph{the} does is use the meaning of its sister to identify an object to which the DP it heads refers. Because the meaning of its sister is a function whose range is True or False, we can use that function to identify the relevant objects by paying attention to those that the function maps to True. Our definition of \emph{the} will be~(\ref{ex5:denthep}).
\begin{exe}
	\raggedright
	\ex \label{ex5:denthep} 
	\denotes{the} = $\lambda P$ refer to $\epsilon$ if $P(\epsilon)$=T
\end{exe}
There are other aspects to the meaning of \emph{the} which are responsible for making sure that the object referred to is salient and ``definite.'' These are the parts to the meaning of \emph{the} that make it complex. I will try to evoke these parts by describing the denotation with English words that have a similar flavor. So, in place of (\ref{ex5:denthep}), I'll use~(\ref{ex5:denthe}).
\begin{exe}
	\raggedright
	\ex \label{ex5:denthe} 
	\denotes{the} = $\lambda P$ that $\epsilon$ such that $P(\epsilon)$=T
\end{exe}
This gives (\ref{ex5:eventn1c}) from (\ref{ex5:eventn1}).
\begin{fullwidth}
\medskip
\begin{exe}
	\raggedright
	\ex \label{ex5:eventn1c} 
	\small
	\begin{forest}
	[that $\epsilon$ such that $\exists x$ $\exists y$ \textsc{exchange}($\epsilon$) $\wedge$ \textsc{theme}($y$)($\epsilon$) $\wedge$ \textsc{agent}($x$)($\epsilon$) \\ DP [that $\epsilon$ such that $\exists x$ $\exists y$ \textsc{exchange}($\epsilon$) $\wedge$ \textsc{theme}($y$)($\epsilon$) $\wedge$ \textsc{agent}($x$)($\epsilon$)\\ \ibar{D}
		[$\lambda P$ that $x$ such that $P(x)$\\ \obar{D} [the] ]
		[$\lambda \epsilon\ \exists x\ \exists y$ \textsc{exchange}($\epsilon$) $\wedge$ \textsc{theme}($y$)($\epsilon$) $\wedge$ \textsc{agent}($x$)($\epsilon$)\\ nP
			[$\lambda P\ \lambda \epsilon\ \exists x$  $P(\epsilon)$ $\wedge$ \textsc{agent}($x$)($\epsilon$)\\ \obar{n} [agent]]
			[$\lambda \epsilon\ \exists y$ \textsc{exchange}($\epsilon$) $\wedge$ \textsc{theme}($y$)($\epsilon$)\\ NP [\ibar{N} [\obar{N} [$\surd$exchange]]] ]
		 ]
	]]
	\end{forest}
\end{exe}
\end{fullwidth}
\medskip
These are sometimes called ``event'' nominalizations.

Notice that in this example, the DP refers to an eventuality. Eventualities have durations (sometimes), and it is for this reason that sentences like~(\ref{ex5:exchangedur}) are sensible.
\begin{exe}
	\raggedright
	\ex \label{ex5:exchangedur} 
	\begin{xlist}
		\ex The exchange took over an hour.
		\ex The exchange began in the morning.
		\ex The exchange got interrupted.
	\end{xlist}
\end{exe}
Not all DPs whose nouns have verbal correlates refer to eventualities, however. For instance, \emph{gift} refers to things, and not events, and for this reason sentences like~(\ref{ex5:gift1}) are odd.
\pagebreak
\begin{exe}
	\raggedright
	\ex \label{ex5:gift1} 
	\begin{xlist}
		\ex The gift took an hour.
		\ex The gift began in the morning.
		\ex The gift got interrupted.
	\end{xlist}
\end{exe}
What \emph{the gift} refers to is an object that was given, not the event of giving. Moreover, it may well be that there is no external $\theta$-role present in NPs headed by \emph{gift}. At least, I find awkward the use of modifiers that depend on the presence of an external $\theta$-role.
\begin{exe}
	\raggedright
	\ex[\#]{\label{ex5:gift2} The intentional gift was in the mailbox.}
\end{exe}
What this suggests is that \emph{gift} occurs without the external $\theta$-role that is required of the verb \emph{give}, and that only the meaning associated with \emph{$\surd$give} attaches to the noun \emph{gift}. To get the fact that \emph{gift} refers to the thing given, we can let Existential Closure bind off the event variable and leave the Theme variable open. The result is portrayed in~(\ref{ex5:gift2}).
\begin{exe}
	\raggedright
	\ex \label{ex5:gift2} 
	\small
	\begin{forest}
	[that $x$ such that $\exists \epsilon$ \textsc{give}($\epsilon$) $\wedge$ \textsc{theme}($x$)($\epsilon$) \\ DP [that $x$ such that $\exists \epsilon$ \textsc{give}($\epsilon$) $\wedge$ \textsc{theme}($x$)($\epsilon$) \\ \ibar{D}
		[$\lambda P$ that $x$ such that $P(x)$\\ \obar{D}[the] ]
		[$\lambda x$ $\exists \epsilon$ \textsc{give}($\epsilon$) $\wedge$ \textsc{theme}($x$)($\epsilon$)=T\\ NP [\ibar{N} [\obar{N} [$\surd$gift]]] ]
	]]
	\end{forest}
\end{exe}
These are sometimes called ``object'' nominals, as the thing they refer to is the object of the noun.

This account predicts that these nouns cannot combine with an \emph{of}-phrase that gets the internal $\theta$-role.%
%
\marginnote[0pt]{See \citet{Zucchi:1989} for this observation. There are apparent counter-examples, as for example: \emph{the gift of grace} and \emph{the gift of apples}. These arise, so far as I can tell, when the object of the noun is indefinite or abstract. There are modes of semantic composition that are reserved for objects of this semantic type, and it seems to me likely that this is what's involved in these apparent counterexamples. This mode of semantic composition is consistent with the account sketched here. If this is how these cases arise, they are not counterexamples then. See \citet{Chung:2004a} and \citet{Dayal:2011} for some discussion.}%
%
That seems to be a correct prediction, as can be seen from~(\ref{ex5:noobj}).
\begin{exe}
	\raggedright
	\ex[]{\label{ex5:noobj} the gift of this basket}
\end{exe}
To the extent that (\ref{ex5:noobj}) has a grammatical interpretation, it refers to a gift that the basket, somehow, is responsible for. This is just what the analysis illustrated in(\ref{ex5:gift2}) predicts.

If this is correct, it requires that \emph{the} be flexible enough for it to combine with NPs or nPs. Further, the denotation of \emph{the} must be flexible enough for it to combine with a phrase whose meaning is a function with a domain of things (as in (\ref{ex5:gift2})) as well as with a phrase whose meaning is a function with a domain of eventualities (as in~(\ref{ex5:eventn1c})). In recognition of this, let's refine the definition of \denotes{the}.
\pagebreak
\begin{exe}
	\raggedright
	\ex \label{ex5:denthef} 
	\denotes{the} = $\lambda P$ that $a$ such that $P(a)$=T, where
	\begin{xlist}
		\ex $a$ is a thing or an eventuality, and
		\ex $P$ is a function with domain $a$.
	\end{xlist}
	%
%
\marginnote[-30pt]{Recall that our convention is to use $\ldots,x,y,z$ to represent objects drawn from the domain of things, $\epsilon$, $\epsilon '$, $\epsilon ''$,\ldots to represent objects drawn from the domain of eventualities, and $P,Q,R$ to represent functions whose domains are either things or eventualities. We now add the convention that $a,b,c,\ldots$ will be used to represent objects from the domains of things or eventualities.}
\end{exe}

Nouns related to verbs can also refer to things that are reifications of the events they describe. For instance, the noun \emph{report} can be used to make a DP that refers to an eventuality, as in~(\ref{ex5:report1}).
\begin{exe}
	\raggedright
	\ex \label{ex5:report1} 
	The report took an hour.
\end{exe}
But it can also be used to build a DP that refers to a thing, and not an eventuality, as in~(\ref{ex5:report2}).
\begin{exe}
	\raggedright
	\ex \label{ex5:report2} 
	The report is several pages long.
\end{exe}
In neither of these cases is \emph{report} making a DP that refers to a thing that would be the object of the verb \emph{report}, as can be seen from~(\ref{ex5:report3}).
\begin{exe}
	\raggedright
	\ex \label{ex5:report3} 
	I reported the violation to the authorities.
	\begin{xlist}
		\ex The report took an hour.
		\ex The report is several pages long.
	\end{xlist}
\end{exe}
(\ref{ex5:report3}a) does not claim that the violation took an hour, and (\ref{ex5:report3}b) does not report that the violation is several pages long. The thing that \emph{the report} refers to in (\ref{ex5:report2}) corresponds in some fashion to an eventuality of reporting. The thing it refers to carries the same information that would be needed to individuate a reporting eventuality. We can imagine that nouns with this meaning are built from those that describe eventualities from a process that converts those descriptions of eventualities to descriptions of things with the same individuating content. However, unlike nouns that describe eventualities, the ones that describe things that correspond to eventualities do not seem to have an external $\theta$-role associated with them. They do not support the use of modifiers that depend on the presence of an external $\theta$-role:
\begin{exe}
	\raggedright
	\ex[\#]{\label{ex5:report4} The intentional report (of the infraction) is several pages long.}
\end{exe}
Here again, the external $\theta$-role does not seem to accompany the root \emph{$\surd$report} in its nominal form, at least when it refers to things. The process that converts a description of eventualities to a description of things must correlate with this suppression of the external $\theta$-role. Nouns that describe things that are related to eventualities are sometimes called ``concrete'' nouns. They are also sometimes called ``result nominals.'' 

We've seen now two ways in which nouns that are related to verbs can build DPs that refer to things, and one way in which nouns refer to eventualities. For object nominals like \emph{gift}, the method involves manipulating the meaning of the noun so that it ends up describing things that are given. For event nominals like \emph{report}, there is a process that takes the event-referring meaning of \emph{report} and makes a noun that refers to eventualities of that kind. And, finally, for concrete nominals like \emph{report}, there is a process that takes the event nominal \emph{report} and converts it to a thing-referring meaning. As these examples indicate, sometimes nouns are ambiguous, as for example \emph{report}; they are able to undergo more than one of these three processes. But some nouns are not, as for example \emph{gift}, which can only undergo one of these processes.

We'll look at one last way in which nouns with verbal meanings can be manipulated to refer to things. In this last example, illustrated by (\ref{ex5:baker1}), the noun refers to a thing that would be the external $\theta$-role bearer of the related verb.\marginnote[0pt]{See \citet{Rappaport:1992} for a discussion of these nominalizations. There are apparent counter-examples to the generalization that nominals built with the suffix \emph{er} refer to the ``subject'' argument of the verb---``bestseller'' is one. See \citet{Alexiadou:2010} for an analysis.}
\begin{exe}
	\raggedright
	\ex \label{ex5:baker1} 
	\begin{xlist}
		\ex The baker arrived late.
		\ex The worriers fretted all night.
		\ex The rollers flattened the dough.
	\end{xlist}
\end{exe}
This meaning correlates with the presence of the nominalizer, \emph{er}. Note that nouns built with \emph{er} do not refer to the things that can more generally be understood to bear the relation that genitives have to the noun. For instance, whereas the genitive DP in (\ref{ex5:painte}a) refers to the owner of the painting, the referent of (\ref{ex5:painte}b) does not.
\begin{exe}
	\raggedright
	\ex \label{ex5:painte} 
	\begin{xlist}
		\ex The Vatican's painting of God touching Adam has thousands of visitors a day.
		\ex The painter of God touching Adam has thousands of visitors a day.
	\end{xlist}
\end{exe}
The referent of \emph{painter of X} is to the person who painted X. In general, it is the $\theta$-role of the root that determines what a noun with \emph{er} refers to.

To capture this, we can let the presence of \emph{er} require that the external $\theta$-role function not be existentially closed, but allow for all other functions to be existentially closed. This produces representations like~(\ref{ex5:baker2}).
\begin{fullwidth}
\medskip
\begin{exe}
	\raggedright
	\ex \label{ex5:baker2} 
	\small
	\begin{forest}
	[that $a$ such that $\exists \epsilon$ $\exists y$ \textsc{bake}($\epsilon$) $\wedge$ \textsc{patient}($y$)($\epsilon$) $\wedge$ \textsc{agent}($a$)($\epsilon$)\\ DP [that $a$ such that $\exists \epsilon$ $\exists y$ \textsc{bake}($\epsilon$) $\wedge$ \textsc{patient}($y$)($\epsilon$) $\wedge$ \textsc{agent}($x$)($\epsilon$)\\ \ibar{D}
		[$\lambda P$ that $a$ such that $P(a)$=T\\\obar{D} [the] ]
		[$\lambda x$ $\exists \epsilon$ $\exists y$ \textsc{bake}($\epsilon$) $\wedge$ \textsc{patient}($y$)($\epsilon$) $\wedge$ \textsc{agent}($x$)($\epsilon$)\\ nP [$\lambda x$ $\exists \epsilon$ $\exists y$ \textsc{bake}($\epsilon$) $\wedge$ \textsc{patient}($y$)($\epsilon$) $\wedge$ \textsc{agent}($x$)($\epsilon$)\\\ibar{n}
			[$\lambda P$ $\lambda x$ $\exists \epsilon$ $P(a)$ $\wedge$ \textsc{agent}($x$)($\epsilon$)=T\\ \obar{n} [agent] ]
			[$\lambda \epsilon$ $\exists y$ \textsc{bake}($\epsilon$) $\wedge$ \textsc{patient}($y$)($\epsilon$)=T\\ NP [$\surd$baker, roof] ]
		] ]
	]]
	\end{forest}
\end{exe}
\end{fullwidth}
\medskip
I'll call these ``subject nominalizations.''

Subject nominalizations allow us to see that the meaning associated with the genitive Case-marked DP can have more than one source, a conclusion that was first, perhaps, defended by Mona Anderson.\cite{Anderson:1983a} In (\ref{ex5:baker3}), the genitive Case-marked DPs cannot be understood as the external $\theta$-role bearers of the nouns.
\begin{exe}
	\raggedright
	\ex \label{ex5:baker3} 
	\begin{xlist}
		\ex Emily's baker (of pies) was late.
		\ex Jason's worriers (about the election) fretted all night.
		\ex The cylinder's rollers (of the dough) did their job.
	\end{xlist}
\end{exe}
This is what is expected if to get the meaning associated with the external $\theta$-role of the verb related to the nouns in (\ref{ex5:baker3}) requires that it be assigned to the DP in genitive Case. Because nouns built with the \emph{er} morpheme require that $\theta$-role to be bound by the determiner \emph{the}, it cannot be assigned to a DP. Instead, the genitive DPs in these examples must be arguments of another predicate, one that leads us to understand that the baker in (\ref{ex5:baker3}a), for instance, ``belongs'' to Emily. If we adopt the view that the determiner responsible for assigning genitive Case has the denotation of \emph{the},%
%
\marginnote[0pt]{See chapter \ref{ch:phrasestructure}.}%
%
and we represent the term that gives the genitive DPs in (\ref{ex5:baker3}) their meaning with $\mathcal{P}$, then (\ref{ex5:baker3}a) could get the representation in~(\ref{ex5:baker3t}) on the next page.

We've seen, then, a variety of ways in which the Theta Criterion is satisfied by projections of nouns. First, unlike verbs, it seems that the nominal versions of a verb's root can expone the root without its accompanying external $\theta$-role. Second, it looks like Existential Closure can apply in a less constrained way to nouns and their $\theta$-roles, including the external $\theta$-role. And finally, it appears that the projection of a noun and its external $\theta$-role need not be a function whose domain is eventualities. It is possible instead for the projection of a noun, or its external $\theta$-role, to be a function whose domain is things. This last option allows for result nominals, and also for DPs that refer to a thing that is the bearer of a $\theta$-role that comes associated with a noun: object and subject nominalizations. This model of how the nominal and verbal versions of a shared root can differ is summarized in~(\ref{ex5:diff}).
\begin{fullwidth}
\medskip
\begin{exe}
	\raggedright
	\ex \label{ex5:baker3t} 
	\small
	\begin{forest}
	[that $a$ such that $\exists \epsilon$  \textsc{bake}($\epsilon$) $\wedge$ \textsc{patient}(pies)($\epsilon$)\\ $\wedge$ \textsc{agent}($a$)($\epsilon$) $\wedge$ \textsc{belongs\_to}(Emily)($a$)  \\ DP
		[DP's\\ case:\emph{gen} [Emily, name=binder, roof] ]
		[that $a$ such that $\exists \epsilon$  \textsc{bake}($\epsilon$) $\wedge$ \textsc{patient}(pies)($\epsilon$)\\ $\wedge$ \textsc{agent}($a$)($\epsilon$) $\wedge$ \textsc{belongs\_to}(Emily)($a$)  \\ \ibar{D}
			[$\lambda P$ that $a$ such that $P(a)$\\\obar{D} 
				[\denotes{the}]
				[$\bullet$case: \emph{gen} ]
			 ]
			[$\lambda w$ $\exists \epsilon$  \textsc{bake}($\epsilon$) $\wedge$ \textsc{patient}(pies)($\epsilon$)\\ $\wedge$ \textsc{agent}($w$)($\epsilon$) $\wedge$ \textsc{belongs\_to}(Emily)($w$) \\ XP
				[{\ }, name=trace ]
				[$\lambda z$ $\lambda w$ $\exists \epsilon$  \textsc{bake}($\epsilon$) $\wedge$ \textsc{patient}(pies)($\epsilon$)\\ $\wedge$ \textsc{agent}($w$)($\epsilon$) $\wedge$ \textsc{belongs\_to}($z$)($w$)\\ \ibar{X} 
					[$\lambda P$ $\lambda z$ $\lambda w$ $P(w)$ $\wedge$ \textsc{belongs\_to}($z$)($w$)\\ \obar{X} ]
					[$\lambda x$ $\exists \epsilon$ bake($\epsilon$) $\wedge$\\ \textsc{patient}(pies)($\epsilon$) $\wedge$ \textsc{agent}($x$)($\epsilon$)\\ nP [$\lambda x$ $\exists \epsilon$ bake($\epsilon$) $\wedge$ \textsc{patient}(pies)($\epsilon$) $\wedge$ \textsc{agent}($x$)($\epsilon$)\\ \ibar{n}
						[$\lambda P$ $\lambda x$ $\exists \epsilon$ $P(a)$ $\wedge$ \textsc{agent}($x$)($\epsilon$)\\ \obar{n} [agent] ]
						[$\lambda \epsilon$ \textsc{bake}($\epsilon$) $\wedge$ \textsc{patient}(pies)($\epsilon$)\\ NP [$\lambda x$ $\exists \epsilon$ bake($\epsilon$) $\wedge$ \textsc{patient}(pies)($\epsilon$) $\wedge$ \textsc{agent}($x$)($\epsilon$)\\ \ibar{N}
							[$\lambda y$ $\lambda \epsilon$ \textsc{bake}($\epsilon$) $\wedge$ \textsc{patient}($y$)($\epsilon$)\\ \obar{N} [baker] ]
							[PP [of pies, roof] ]
						] ]
					 ]
				]
			 ]
		 ]]
	]\draw[->,>=stealth', ultra thin, gray] (trace) to[out=south west,in=south,looseness=1] node[pos=0.75,xshift=-10mm]{} (binder);
	\end{forest}
\end{exe}
\end{fullwidth}
\begin{exe}
	\raggedright
	\ex \label{ex5:diff} 
	\begin{xlist}
		\ex A root need not be accompanied by its External $\theta$-role in nominal projections even when it must in verbal projections.
		\ex Existential Closure is more widely available in NP/nP than it is in VP/vP.
		\ex VP/vP must be a function with a domain of eventualities, but NP/nP may also be a function with a domain of things.
	\end{xlist}
\end{exe}

There are event nominals which behave more like verbal projections with regard to how the Theta Criterion applies to them. Jane Grimshaw\cite{Grimshaw:1990b} discovered one that involves putting into the NP a modifier that depends on the so-called aktionsart of the root. Aktionsart refers to ways of characterizing certain logical properties of the eventualities that verbs, and nouns, are descriptions of. When we first adopted a semantics of verbs built around eventualities, I noted that there are two ``kinds'' of eventualities: states and events. These are two Aktionsarten. Eventualities that are states have a kind of homogeneity relative to time that other eventualities don't have. Like all other things that natural language describes, eventualities are situated in a temporal matrix. Let's call the duration of an eventuality, its ``run time.'' The eventuality referred to in \emph{Satoshi walked} occurs at a particular (previous) slice of the time line---that slice is its run time.%
%
\marginnote[-25pt]{Sometimes the distinctions I'm describing here with the term ``aktionsarten'' are instead said to involve ``aspect.'' In general terms, ``aspect'' refers to ways in which language relates an eventuality to a time. We can distinguish two ways in which languages do this. One is built into the denotation of the verbs, adjectives, and nouns that describe eventualities. The other is determined by morphological means that are somewhat independent of the lexical item that describes the eventuality. The difference between the perfect (e.g., \emph{Satoshi had walked}) and the progressive (\emph{Satoshi is walking}) is an example of the latter type. I will reserve the term ``aspect'' for this type, and use aktionsart for the kind that is built into the denotation of a root. Sometimes aktionsart is called ``lexical aspect'' or ``inner aspect.''}%
%
 If we examine the eventuality of Satoshi's walking moment by moment, we can observe that there will be slices of time within the run time that do not have the required attributes to be a walking eventuality. There will be moments, for instance, in which Satoshi is in a two-feet-on-the-ground stance that we might also find in portions of the run time of the eventuality described by \emph{Satoshi stood} or \emph{Satoshi fell}. There is a minimum portion of the run time of the eventuality described by \emph{walk} below which we no longer have an eventuality that can be individuated as a walking. Eventualities that are states differ in this regard. The eventuality referred to in \emph{Satoshi loved nattoo}, for example, has a run time all of whose moments are lovings of nattoo.%
 %
 \marginnote[0pt]{This way of characterizing the difference between states and events comes from \citet{Dowty:1979}. There are other properties that distinguish events and states. A better definition would explain why those properties correlate.}%
 %
 Eventualities that are like those described by \emph{love} are states. Other state descriptions are \emph{know, fear, be, sad, blue}. Eventualities that are like those described by \emph{walk} are events. Other verbs that describe events are \emph{sneeze, arrive, bake}.

Events can be divided further into two types of aktionsarten. There are events that have an endpoint built into their meaning. For instance, the eventuality referred to in \emph{Satoshi baked something} necessarily contains a moment at which something is baked. To see this, compare the eventuality referred to in \emph{Satoshi baked something} to the eventuality referred to in \emph{Satoshi heated something}. What distinguishes these eventualities is their results. The baking eventuality includes an eventuality that could be described by \emph{heat}. A baking is a heating that culminates with something being baked. Eventualities that are individuated by virtue of their culminations---the properties of the final moment in their run time---are known as ``telic.''%
%
\marginnote[-10pt]{The term ``telic'' comes from \citet{Garey:1957}, though the distinction is found in Aristotle's \emph{Metaphysics}.}%
%
Other telic predicates are \emph{die, make, build}, and \emph{appear}. Events which aren't telic are called ``atelic.''

It may seem counter-intuitive to think of telic events as requiring endpoints because it is possible to use telic predicates to describe eventualities which don't have endpoints. This is possible, for instance, in the so-called ``progressive'' construction in English. As (\ref{ex5:baking}) indicates, there need be no resultant baked thing when using the verb \emph{bake}.
\begin{exe}
	\raggedright
	\ex \label{ex5:baking} 
	Satoshi was baking a pie when the house caught fire and burnt to the ground.
\end{exe}
What the progressive does in (\ref{ex5:baking}) is refer to a portion of the baking eventuality whose run time doesn't include the endpoint. What matters for the purpose of determining the aktionsart of an eventuality is how it is individuated. That is, what portions of its run time are necessary in distinguishing one kind of eventuality from another kind of eventuality. Telic eventualities require that their run time have a last moment by virtue of which they are individuated, and atelic ones don't. Once individuated, language has ways of referring to portions of an eventuality. The progressive is one of those ways.

Note too that the syntax can affect the aktionsart of an event description. For instance, though the verb \emph{walk} describes an atelic event, the VP \emph{walk a kilometer} describes a telic event. That is, whether \emph{walk a kilometer} is true of an event depends on properties of the last moment in its run time.%
%
\marginnote[-20pt]{\citet{Vendler:1957} divides telic events further into ones that are punctual (known as ``achievements'') and ones that are gradual (known as ``accomplishments''). The event described by \emph{die} is an achievement; the one described by \emph{bake} is an accomplishment. \citet{Comrie:1976} similarly divides atelic events into ones that are punctual (known as ``semelfactives'') and ones that are aren't (known as ``activities''). The event described by \emph{cough} is semelfactive, and the event described by \emph{walk} is an activity.}%
%

Many grammatical phenomena are sensitive to these distinctions in aktionsarten. One of interest to the task at hand is that certain temporal expressions which measure the duration of an event are compatible with atelic events but not telic ones. A classic example of this is \emph{for X}, where ``X'' is an expression that refers to an interval of time. We can see this in the contrasts in~(\ref{ex5:for1}).%
%
\marginnote[20pt]{Note that the use of \emph{for}-phrases with telic predicates is ``wrong'' if they are meant to tell the duration of the event described by that predicate. They are grammatical under other interpretations. For instance (\ref{ex5:for1}g) can have the meaning that Jesse was dead for a full second; that is, it isn't ill-formed if Jesse isn't dead, but was previously.}%
%
\begin{exe}
	\raggedright
	\ex \label{ex5:for1} 
	\begin{xlist}
		\sn \emph{atelic:}
		\ex[]{Eva ran for an hour.}
		\ex[]{Jesse danced for just three minutes.}
		\ex[]{Eva sneezed for a full second.}
		\ex[]{Jesse observed the experiment for an hour.}
		\sn \emph{telic:}
		\ex[\#]{Jesse built it for an hour.}
		\ex[\#]{Eva arrived for just three minutes.}
		\ex[\#]{Jesse died for a full second.}
		\ex[\#]{Eva destroyed the experiment for an hour.}
	\end{xlist}
\end{exe}
There is something semantically anomalous about using these temporal modifiers with predicates that describe telic events. By contrast, the modifier \emph{in X}, where ``X'' is an expression that refers to an interval of time, is compatible with predicates that describe telic events, but anomalous when used with predicates that describe atelic ones:
\begin{exe}
	\raggedright
	\ex \label{ex5:in1} 
	\begin{xlist}
		\sn \emph{atelic:}
		\ex[\#]{Eva ran in an hour.}
		\ex[\#]{Jesse danced in just three seconds.}
		\ex[\#]{Eva sneezed in a full second.}
		\ex[\#]{Jesse observed the experiment in an hour.}
		\sn \emph{telic:}
		\ex[]{Jesse built it in an hour.}
		\ex[]{Eva arrived in just three seconds.}
		\ex[]{Jesse died in a full second.}
		\ex[]{Eva destroyed the experiment in an hour.}
	\end{xlist}
\end{exe}
There is something about the meanings of these modifiers that interacts with the properties of eventualities to produce these results. 

What's of interest to us is that nouns that share a root with a verb can also be divided with respect to this classification of eventuality types. We can see this from the contrasts in~(\ref{ex5:nounaks1}).
\begin{exe}
	\raggedright
	\ex \label{ex5:nounaks1} 
	\begin{xlist}
		\sn \emph{atelic:}
		\ex[\#]{Only observation of the patient in several weeks (can determine the most likely course of action).}
		\ex[]{Only observation of the patient for several weeks (can determine the most likely course of action).}
		\sn \emph{telic:}
		\ex[]{The total destruction of the city in two days (appalled everyone).}
		\ex[\#]{The total destruction of the city for two days (appalled everyone).}
	\end{xlist}
	\hspace*{\fill} \citep[(28-9):~58]{Grimshaw:1990b} \hspace*{.5in}
\end{exe}
Grimshaw discovered that nominals with modifiers that require a certain aktionsart must contain their object arguments. This is why there are the contrasts in~(\ref{ex5:objakt}).
\begin{exe}
	\raggedright
	\ex \label{ex5:objakt} 
	\begin{xlist}
		\ex[]{Only observation (can determine the most likely course of action).}
		\ex[*]{Only observation for several weeks (can determine the most likely course of action).}
		\ex[]{The total destruction (appalled everyone).}
		\ex[*]{The total destruction in two days (appalled everyone).}
	\end{xlist}
\end{exe}
Note that the subject argument isn't required; it's hard to tell whether that is because the external $\theta$-role need not accompany the noun or that is because existential closure is available for the external $\theta$-role. What we learn from these examples is that the aktionsart of the event described by some term is tied to the presence of that term's object. As we've seen, it's possible for a noun to describe an eventuality without its object arguments being expressed overtly. But if those eventualities are to have the properties of aktionsarten that support modification, they must come with an overt expression of the object. This suggests that the denotation of the object contributes to determining the aktionsart of an event description.%
%
\marginnote[0pt]{See \citet{Borer:2005} for an account that would have this consequence.}%
%

One scenario, then, in which nouns behave more like verbs with respect to the Theta Criterion is when they match the aktionsart of the verb they are related to. Another situation that has this effect on nouns arises, Grimshaw argues, when nouns bring with them their external $\theta$-role. We've seen that the ``subject'' arguments of nouns can get a $\theta$-role that is not the one characteristically associated with the root. In those scenarios, I suggested that there was hidden inside the nominal another term that assigned that $\theta$-role (see~(\ref{ex5:baker3})). In order to see Grimshaw's effect, it is therefore necessary to create examples in which the subject argument is getting the $\theta$-role associated with the root, and not this other one. If we pick roots that come with an Agent $\theta$-role, we can do this by using the modifiers that are licensed by the presence of an Agent $\theta$-role. Grimshaw uses the modifier \emph{deliberately} to do this.
\begin{exe}
	\raggedright
	\ex \label{ex5:grims} 
	\begin{xlist}
		\ex[]{The instructor's deliberate examination of the patient (took a long time).}
		\ex[*]{The instructor's deliberate examination (took a long time).}
	\end{xlist}
	\hspace*{\fill} \citep[(11):~51]{Grimshaw:1990b} \hspace*{.5in}
\end{exe}
The same fact emerges if we use \emph{intentional}. Grimshaw didn't determine whether it was the presence of the external $\theta$-role, or the presence of the external argument, that has this effect. If it is the presence of the external argument that forces the expression of the object argument, then the explanation for the effect must come from how the Theta Criterion is satisfied. In our system, we force the presence of arguments through our semantics by constraining Existential Closure. To force the presence of an object argument in these situations, then, will require that Existential Closure be denied of the internal $\theta$-role when the external $\theta$-role is present and Existential Closure does not apply to it. If, by contrast, the effect illustrated by (\ref{ex5:grims}) holds when the external $\theta$-role is present, but not expressed with an argument, then Existential Closure of the internal $\theta$-role is blocked when there is an external $\theta$-role. Whether that external $\theta$-role is existentially closed or not doesn't matter on this second scenario. We can tell which is the correct generalization by looking at examples where there is no external argument, but we can detect the presence of the external $\theta$-role in another way. Our method for detecting the presence of an external $\theta$-role is to use the modifiers that are sensitive to their presence. So, what's at stake is whether the contrast in (\ref{ex5:grims}) shows up without the external argument, but with the agent oriented modifier. Are the contrasts in (\ref{ex5:compgrims1}) mirrored in~(\ref{ex5:compgrims2})?
\begin{exe}
	\raggedright
	\ex \label{ex5:compgrims1} 
	\begin{xlist}
		\ex[]{her deliberate examination of the problem}
		\ex[*]{her deliberate examination}
		\ex[]{his deliberate construction of the building}
		\ex[*]{his deliberate construction}
		\hspace*{\fill} \citep[(8):~4]{Moulton:2014} \hspace*{.5in}
		\ex[]{your intentional observation of the conflagration}
		\ex[*]{your intentional observation}
		\ex[]{my intentional destruction of these theories}
		\ex[*]{my intentional destruction}
	\end{xlist}
	\pagebreak
	\ex \label{ex5:compgrims2} 
	\begin{xlist}
		\ex[]{the deliberate examination of the problem}
		\ex[?]{the deliberate examination}
		\ex[]{the deliberate construction of the building}
		\ex[?]{the deliberate construction}
		\ex[]{the intentional observation of the conflagration}
		\ex[?]{the intentional observation}
		\ex[]{the intentional destruction of these theories}
		\ex[?]{the intentional destruction}
	\end{xlist}
\end{exe}
I am frankly unable to tell.

The last way to make nominals behave more like the verbs they are related to with respect to the Theta Criterion is to use the so-called ``gerund.'' This term is used to refer to verbal forms that appear with the suffix \emph{ing}, but are not progressive. We can distinguish three cases, illustrated by the bracketed portions of~(\ref{ex5:gerundsint}).%
%
\marginnote[0pt]{\citet{Abney:1987} has a good introduction to English gerunds, and see also \citet{Milsark:1988}.}%
%
\begin{exe}
	\raggedright
	\ex \label{ex5:gerundsint} 
	\begin{xlist}
		\ex Della remembered [him making cookies].
		\ex Della talked about [his making cookies].
		\ex Della talked about [his making of cookies].
	\end{xlist}
\end{exe}
The form in (\ref{ex5:gerundsint}a) is a clause, whose accusative Case-marked subject comes from the $\mu$P embedding \emph{remember}, as we can see from the effect that the Passive has on this kind of sentence: \emph{He was remembered making cookies}. (\ref{ex5:gerundp1}) sketches its syntax.
\begin{exe}
	\raggedright
	\ex \label{ex5:gerundp1} 
	\small
	\begin{forest}
	[\ibar{V}
		[\obar{V} [$\surd$remember]]
		[TP 
			[DP [him, roof] ]
			[\ibar{T} [making cookies, roof] ]
		]
	]
	\end{forest}
\end{exe}

The examples in (\ref{ex5:gerundsint}b) and (\ref{ex5:gerundsint}c) are DPs, and of more interest to us. The nominal in (\ref{ex5:gerundsint}c) is eventive, and has the properties we have already seen of eventive nouns.%
%
\marginnote[-80pt]{\citet{Chomsky:1970} suggested that gerunds are related to what we would now call vPs by a syntactic process that is part of the derivation. Gerunds, in other words, are built syntactically from projections of verbs, unlike other nouns which are constructed morphologically from a root that is shared by the nominal and verbal subclasses. \citet{Fraser:1970ab} has the first defense of this view of gerunds.}%
%
Existential Closure can apply to make either of its external or internal arguments non-obligatory.
\begin{exe}
	\raggedright
	\ex \label{ex5:cookies2} 
	\marginnote[20pt]{The only $\theta$-role available to the genitive DP in (\ref{ex5:cookies2}c) is the external $\theta$-role that comes with the root \emph{$\surd$make}. For some reason, $\mathcal{P}$, the term that provides the general relation that genitives can have to the NP that follows them, is absent in these gerunds. For this reason, the internal argument is obligatory in this configuration.}
	\begin{xlist}
		\ex[]{the making of cookies}
		\ex[?]{the making}
		\ex[*]{his making}
	\end{xlist}
\end{exe}
In (\ref{ex5:gerundsint}b), however, we have a DP with what appears to be a VP inside. The object of \emph{make} here bears accusative Case---the structural Case associated with $\mu$. A preliminary sketch of its syntax is~(\ref{ex5:accingp}).
\begin{exe}
	\raggedright
	\ex \label{ex5:accingp} 
	\small
	\begin{forest}
	[\ibar{D}
		[\obar{D} [case: \emph{gen}]]
		[IP 
			[\obar{I} [ing] ]
			[vP 
				[DP [his, roof] ]
				[\ibar{v} 
					[\obar{v} [agent] ]
					[$\mu$P [\ibar{$\mu$}
						[\obar{$\mu$} [case: \emph{acc}] ]
						[VP [\ibar{V}
							[\obar{V} [$\surd$make] ]
							[DP [cookies, roof] ]
						] ]
					]]
				]
			 ]
		]
	]
	\end{forest}
\end{exe}
As expected, because the accusative Case assigner comes into a D-structure with the external $\theta$-role---Burzio's Generalization correlates them---the \obar{v} that carries the external $\theta$-role is found in these nominals. Like external $\theta$-roles outside nominals, Existential Closure cannot apply, and the external argument is obligatory, as (\ref{ex5:accings}) indicates.
\begin{exe}
	\raggedright
	\ex[*]{\label{ex5:accings} Della talked about [the making cookies].}
\end{exe}
These kinds of gerunds might have no noun in them at all. As expected, the Theta Criterion applies to these cases in the same way that it applies to~VPs.

\subsection{A gap in clausal complementation} 
\label{sub:clausalcomplementation}

As we've seen, when nouns that are related to verbs take arguments, they match the $\theta$-role and c-selection of the verbs. We've explained this by letting these nouns be built from the same root that the related verbs are built from, and therefore the $\theta$-roles and c-selection requirements that travel with these roots are present in the nominals as well. But there are some systematic ways in which the argument structure of nominals is not like that of the verbs that they are related to. It is this fact that \citet{Kayne:1984d} uses to argue for the structure of double object constructions we were examining at the end of the previous section. We're now ready to examine this argument.

One of these differences concerns the kinds of clauses that nouns can take as objects. As expected, nouns that are related to verbs can select clausal complements if the verbs do. Thus, we have pairs like~(\ref{ex5:clausencom1}).
\pagebreak
\begin{exe}
	\raggedright
	\ex \label{ex5:clausencom1} 
	\begin{xlist}
		\ex 
		\begin{xlist}
			\ex Amalia's saying that nouns are interesting (puzzled us).
			\ex Amalia said that nouns are interesting.
		\end{xlist}
		\ex
		\begin{xlist}
			\ex Amalia's declaration that nouns are interesting (puzzled us).
			\ex Amalia declared that nouns are interesting.
		\end{xlist}
		\ex
		\begin{xlist}
			\ex Christine's attempt to understand nouns (amazed us).
			\ex Christine attempted to understand nouns.
		\end{xlist}
		\ex 
		\begin{xlist}
			\ex Christine's desire to understand nouns (puzzled us).
			\ex Christine desired to understand nouns.
		\end{xlist}
		\ex 
		\begin{xlist}
			\ex Susanne's trying to untangle nominalizations (took all day).
			\ex Susanne tried to untangle nominalizations.
		\end{xlist}
		\ex
		\begin{xlist}
			\ex Amalia's question whether to let roots be shared (wasn't seriously entertained).
			\ex Amalia questioned whether to let roots be shared.
		\end{xlist}
	\end{xlist}
\end{exe}
But interestingly, verbs and adjectives that select raising or exceptional Case marking infinitives are not matched by nouns that can. This is what we see in the contrasts in~(\ref{ex5:clausencom2}).
\begin{exe}
	\raggedright
	\ex \label{ex5:clausencom2} 
	\begin{xlist}
		\ex 
		\begin{xlist}
			\ex[]{Wendell considers Tanja to be strong.}
			\ex[*]{Wendell's consideration of Tanja to be strong}
		\end{xlist}
		\ex
		\begin{xlist}
			\ex[]{Wendell appears to be clever.}
			\ex[*]{Wendell's appearance to be clever}
		\end{xlist}
		\ex 
		\begin{xlist}
			\ex[]{This proposal is likely to be untenable.}
			\ex[*]{This proposal's likelihood to be untenable}
		\end{xlist}
	\end{xlist}
\end{exe}
Not all speakers seem to find fully acceptable the gerunds that include VPs within them, instead of NPs, like that illustrated in~(\ref{ex5:accingp}). But for those that do, we have minimal pairs of the kind in~(\ref{ex5:calusencom3}).
\begin{exe}
	\raggedright
	\ex \label{ex5:calusencom3} 
	\begin{xlist}
		\ex 
		\begin{xlist}
			\ex[?]{Wendell's considering Tanja to be strong (amuses him).}
			\ex[*]{Wendell's considering of Tanja to be strong (amuses him).}
		\end{xlist}
		\ex
		\begin{xlist}
			\ex[?]{Wendell's believing Tanja to be clever (amuses him).}
			\ex[*]{Wendell's believing of Tanja to be clever (amuses him).} 
		\end{xlist}
	\end{xlist}
\end{exe}
Infinitival clauses to nouns are possible if they are not of the raising or exceptional Case marking type. In particular, infinitival clauses that have \textsc{pro} subjects are fine.
\pagebreak
\begin{exe}
	\raggedright
	\ex \label{ex5:controlncom} 
	\begin{xlist}
		\ex Nick's eagerness to please
		\ex Helen's promise to be on time
		\ex Noah's attempt to leave
		\ex Suzi's refusal to give up
		\ex Claire's desire to stay
		\ex Aynat's offer to continue
		\ex Robert's demand to run
	\end{xlist}
	\marginnote[-80pt]{If the \textsc{pro} in the infinitival is ``controlled'' by an object, however, these infinitives seem to be blocked in nominals. 
\begin{exe}
	\raggedright
	\exi{i.} 
	\begin{xlist}
		\ex[*]{Mary's compulsion of her husband to quit his job}
		\ex[*]{her encouragement of John to take up linguistics}
		\ex[*]{their obligation of John to join the army}
		\ex[*]{its impulsion of Mary to sue for divorce}
		\ex[*]{her permission of John to buy himself a new car}
		\ex[*]{their allowance of John to smoke}
	\end{xlist}
	\hspace*{\fill} \citep[(84):~170]{Kayne:1981b} \hspace*{.5in}
\end{exe}}
\end{exe}
There is something about the properties of raising and exceptional Case marking infinitives that makes them incapable of being complements to nouns. Kayne discovered this pattern, though the cases that support the generalization can be found in \citet{Chomsky:1970}.

Kayne's suggestion about what that might be involves the fact that in complement clauses of the sort that are blocked in nominals, there is a relationship between an argument in that complement clause and something outside of the complement clause. To see this, consider the difference between infinitival complements of verbs that house \textsc{pro}, for instance, and infinitival complements that house an accusative Case-marked~DP.
\begin{fullwidth}
\begin{exe}
	\medskip
	\raggedright
	\ex \label{ex5:contrn1} 
	\small
	\begin{forest}
	[\ibar{$\mu$}
		[\obar{$\mu$} [case: \emph{acc}] ]
		[VP [\ibar{V}
			[\obar{V} [$\surd$question] ]
			[CP [\ibar{C}
				[\obar{C} [whether] ]
				[TP
					[DP [\textsc{pro}, roof] ]
					[\ibar{T}
						[\obar{T} [to] ]
						[vP [share roots, roof] ]
					 ]
				 ]
			] ]
		] ]
	]
	\end{forest}
	\hspace*{-20pt}
	\begin{forest}
	[\ibar{$\mu$}
		[\obar{$\mu$} [case: \emph{acc}] ]
		[VP [\ibar{V}
			[\obar{V} [$\surd$consider] ]
			[TP
				[DP\\ case: \emph{acc} [Tanja, roof] ]
				[\ibar{T}
					[\obar{T} [to] ]
					[vP [be strong, roof] ]
				 ]
			 ]
		] ]
	]
	\end{forest}
	\medskip
\end{exe}
\end{fullwidth}
In the case of \emph{consider}, the subject argument, \emph{Tanja}, receives its Case from the $\mu$ that belongs with the VP headed by \emph{consider}. Nothing like that is happening with the complement to \emph{wonder}, however. In raising, there is a similar relationship between an argument in the complement clause and a Case marking term outside the complement, as can be seen in~(\ref{ex5:raisingn1}).
\enlargethispage*{60pt}
\begin{exe}
	\raggedright
	\ex \label{ex5:raisingn1} 
	\small
	\begin{forest}
	[\ibar{T}
		[\obar{T} [$\bullet$case: \emph{nom}] ]
		[VP [\ibar{V}
			[\obar{V} [$\surd$appear] ]
			[TP [\ibar{T}
				[\obar{T} [to] ]
				[VP [\ibar{V}
					[\obar{V} [be] ]
					[AP
						[DP\\ case: \emph{nom} [Wendell, roof] ]
						[\ibar{A} [happy, roof] ]
					 ]
				] ]
			] ]
		] ]
	]
	\end{forest}
\end{exe}
In (\ref{ex5:raisingn1}), \emph{Wendell} gets its Case feature valued by the nominative in~\obar{T}. 

Kayne suggests that the clausal complement to nouns is opaque for the relations that need to occur across infinitives of these kinds.%
%
\marginnote[0pt]{Kayne's own proposal was expressed in terms quite different from those being used here. He suggested that government across the infinitive is blocked in nominals, where ``government'' is a relation akin, and ancestral, to our present conditions on Case assignment.}%
%
\ Consider the form that nominal versions of these configurations will have.
\begin{fullwidth}
\begin{exe}
	\medskip
	\raggedright
	\ex \label{ex5:kclaus1} 
	\small
	\begin{forest}
	[\ibar{D}
		[\obar{D} ]
		[NP [\ibar{N}
			[\obar{N} [$\surd$appear\\ $\surd$consider] ]
			[TP [\ibar{T}
				[\obar{T} [to] ]
				[VP [\ibar{V} 
					[\obar{V} [be] ]
					[AP 
						[DP\\ case: [Wendell, roof] ]
						[AP [\ibar{A} [\obar{A} [happy]]] ]
					 ]
				] ]
			] ]
		] ]
	]
	\end{forest}
	\bigskip
\end{exe}
\end{fullwidth}
The DP, \emph{Wendell}, must get its Case feature valued. If \obar{D} in (\ref{ex5:kclaus1}) is equipped with genitive Case, then that is one source for its value. Alternatively, the structural Case expressed with \emph{of}, which I have assumed is assigned by a noun, could also be the source for its value. If TP, however, is opaque for \textsc{agree}, then neither of these will be able to value the Case feature on \emph{Wendell}. Kayne's hypothesis, then, amounts to the claim that cases such as these are blocked by the Case Filter. 

This account requires that clauses always have a DP within them that is subject to the Case filter. For only then can the ungrammaticality of this complementation structure be credited to a failure to satisfy the Case filter; there must always be a DP within these infinitives that requires Case form outside the infinitive if all instances of these infinitives are to be blocked. When the verb of the embedded clause c-selects one or more DP arguments, this will occur as long as one of those DPs relies on Case from a term outside the clause. Because these clauses are infinitives, the nominative Case that is typically capable of satisfying the Case needs of one DP will be absent. If the verb of the embedded clause c-selects a DP external argument, this DP will have to get its Case from a term outside the clause, and the failure of Case assignment that Kayne's proposal entails will lead to a violation of the Case filter. If the verb of the embedded clause doesn't c-select an external argument, but does c-select a DP internal argument, then that internal argument will not get structural Case from the verb. Burzio's generalization requires of a verb that comes equipped with structural Case that it also assign an external $\theta$-role, so verbs that don't c-select an external argument won't deliver Case to their object DPs. In this case too, then, Kayne's proposal would prevent these clauses from being embedded within a TP. 

But what happens when the embedded clause contains a verb that does not c-select a DP? If the embedded TP is subject to the EPP, and the EPP is satisfied when a DP occupies the Specifier of the relevant TP, then there will have to be some DP that occupies its Specifier in one of the parses in its derivation. If this DP is subject to the Case filter, then it will be that DP which violates the Case filter. The syntax that this view requires is shown in~(\ref{ex5:raiseexpl}).
\begin{exe}
	\raggedright
	\ex \label{ex5:raiseexpl} 
	\small
	\begin{forest}
	[TP
		[DP\\ case: \emph{nom} [it, name=binder, roof] ]
		[\ibar{T}
			[\obar{T} [will\\ $\bullet$case: \emph{nom}] ]
			[VP [\ibar{V}
				[\obar{V} [appear] ]
				[TP
					[{\ }, name=trace ]
					[\ibar{T}
						[\obar{T} [to] ]
						[VP [seem that pigs fly, roof] ]
					 ]
				 ]
			] ]
		 ]
	]
	\draw[->,>=stealth', ultra thin, gray] (trace) to[out=south west,in=south,looseness=1] node[pos=0.75,xshift=-10mm]{} (binder);
	\end{forest}
\end{exe}
\pagebreak
The expletive \emph{it} must occupy the Specifier of the embedded infinitival TP, and in that position \textsc{agree} with the nominative Case assigner, \emph{will}, in the higher clause. Because of the $\bullet$, this \textsc{agree} relation triggers the movement that is shown. If that syntax could be enforced, then Kayne's proposal will extend to explain why these infinitives with this argument structure are not found in nominals. Their syntax would require \textsc{agree} to hold across the TP as indicated in~(\ref{ex5:raiseexplnom}).
\begin{exe}
	\raggedright
	\ex \label{ex5:raiseexplnom} 
	\small
	\begin{forest}
	[DP 
		[\obar{D}\\ case \emph{gen} ]
		[NP [\ibar{N}
			[\obar{N} [$\surd$appear] ]
			[TP
				[DP\\ case: [it, roof] ]
				[\ibar{T}
					[\obar{T} [to] ]
					[VP [seem that pigs fly, roof] ]
				 ]
			 ]
		] ]
	]
	\end{forest}
\end{exe}
If the TP complements to nouns block \textsc{agree}, then this requirement will be frustrated in~(\ref{ex5:raiseexplnom}).

Finally, let's consider scenarios in which the infinitive contains no DP arguments that require Case, and the EPP is satisfied without the use of an expletive. The EPP can be satisfied by CPs, as (\ref{ex5:cpepps}) shows, and so we might expect on Kayne's proposal that infinitives that contain no DPs whatsoever can show up in nominals.
\begin{exe}
	\raggedright
	\ex \label{ex5:cpepps} 
	That pigs fly seems to be unlikely.
\end{exe}
These infinitives would contain CPs in their Specifiers. Unexpectedly, however, CPs cannot surface in the Specifier of infinitival TPs, as the ungrammaticality of (\ref{ex5:nocpepp}) shows.
\begin{exe}
	\raggedright
	\ex \label{ex5:nocpepp} 
	\begin{xlist}
		\ex[*]{She believes [\TP [\CP that pigs fly] to be unlikely].}
		\ex[*]{She considers [\TP [\CP that pigs fly] to be unlikely].}
		\ex[*]{She finds [\TP [\CP that pigs fly] to be unlikely].}
	\end{xlist}
\end{exe}
Instead, the word-order that is preferred employs an expletive:
\begin{exe}
	\raggedright
	\ex \label{ex5:nocpepp2} 
	\begin{xlist}
		\ex She believes it to be unlikely [\CP that pigs fly].
		\ex She considers it to be unlikely [\CP that pigs fly].
		\ex She finds it to be unlikely [\CP that pigs fly].
	\end{xlist}
\end{exe}
In some of these cases, it is also possible to put the CP in the final position without an accompanying expletive.
\pagebreak
\begin{exe}
	\raggedright
	\ex \label{ex5:nocpepp3} 
	\begin{xlist}
		\ex[?]{She believes to be unlikely [\CP that pigs fly].}
		\ex[?*]{She considers to be unlikely [\CP that pigs fly].}
		\ex[*]{She finds to be unlikely [\CP that pigs fly].}
	\end{xlist}
\end{exe}
This word order might arise by extraposition of the CP. If so, we might expect these cases to arise in nominalizations as there is no obvious term within the infinitive that requires an \textsc{agree} relation to a term outside of the infinitive. This expectation is frustrated, however:\marginnote[0pt]{The particulars of Kayne's account did extend to account for the ungrammaticality of (\ref{ex5:nocpepp4}). This is because his syntax for movement required that they left silent pronouns, and these could be viewed as subject to his condition. We will encounter this account of movement in a later chapter.}
\begin{exe}
	\raggedright
	\ex[*]{\label{ex5:nocpepp4} Her belief to be unlikely that pigs fly.}
\end{exe}
If the reason infinitives of this kind are not possible complements to nouns is that the complements to nouns are opaque to \textsc{agree}, then the ungrammaticality of (\ref{ex5:nocpepp4}) is unexpected.

There are other systematic absences in the complement types to nouns. Another that Kayne proposes is that ``small clauses'' cannot be found as complements to nouns. There are no instances where a verb that selects a small clause complement can continue to do so in its noun version. The pairs in~(\ref{ex5:7.14}) illustrate this generalization.
\begin{exe}
	\ex \label{ex5:7.14} 
	\begin{xlist}
	\ex[]{I believe this fact interesting.}
	\ex[*]{my belief of this fact interesting}
	\ex[]{I consider this sentence instructive.}
	\ex[*]{her consideration of this sentence instructive}
	\ex[]{Some find these facts remarkable.}
	\ex[*]{your finding of these facts remarkable}
	\end{xlist}
\end{exe}
In each case, a verb that selects a small clause is matched with a cognate noun, and the subject argument of that small clause is assigned Case from a term outside the small clause: either accusative Case in the situations with verbs, or the Case expressed with \emph{of} in the situations with nouns. The same kind of contrast is found when the Case that the subject of the small clause is assigned comes from a higher position, as (\ref{ex5:7.15}) shows. 
\begin{exe}
	\ex \label{ex5:7.15} 
	\begin{xlist}
	\ex[]{This fact is believed interesting.}
	\ex[*]{this fact's belief interesting}
	\ex[]{This sentence was considered instructive.}
	\ex[*]{this sentence's consideration instructive}
	\ex[]{These facts were found remarkable.}
	\ex[*]{these facts' finding remarkable}
	\end{xlist}
\end{exe}
When the small clause's argument is assigned nominative Case, and therefore moves, the result is grammatical. But when that argument is assigned genitive Case, again triggering movement, the result is ungrammatical. Like infinitival TPs, small clauses cannot be complements to a noun. If small clauses become opaque to \textsc{agree} when they are complements to nouns, the ungrammaticality of these examples is understood.

\citet{Kayne:1981b}'s argument for the existence of a small clause in the double object comes from the observation that double objects are also not possible in nominalizations. There are contrasts parallel to those in (\ref{ex5:7.14}) and~(\ref{ex5:7.15}) which involve double object constructions.
\begin{exe}
	\raggedright
	\ex \label{ex5:kaynedoa} 
	\begin{xlist}
		\ex[]{Her teacher gave Mary the letter.}
		\ex[*]{Her teacher's gift of Mary (of) the letter}
		\ex[]{Her classmate offered her a crayon.}
		\ex[*]{Her classmate's offer of her (of) a crayon}
		\ex[]{This salesman sold us that defective car.}
		\ex[*]{This saleman's sale of us (of) that defective car}
		\ex[]{A friend of his rented John office space.}
		\ex[*]{A friend's rental of John (of) office space}
	\end{xlist}
\hspace*{\fill} \citep[based on (61)\&(62): 161]{Kayne:1981b} \hspace*{.5in}
\end{exe}
If what makes small clause complements incompatible with nouns is that they produce a barrier for \textsc{agree}, and the source of the ungrammaticality of the double object frame in nominals has the same source, then this is reason to give them parallel structures. This is the argument \citet{Kayne:1984d} provides for giving the double object frame a structure like that in (\ref{ex5:kayne1}), repeated below. 
\begin{exe}
	\raggedright
	\exr{ex5:kayne1} 
	\small
	\begin{forest}
	[vP
		[DP [Ania, roof]]
		[\ibar{v} 
			[\obar{v} [agent] ]
			[VP [\ibar{V}
				[\obar{V} [$\surd$throw] ]
				[VP 
					[DP [Stephen, roof] ]
					[\ibar{V}
						[\obar{V} [\shade{have}] ]
						[DP [the ball, roof] ]
					 ]
				 ]
			] ]
		]
	]
	\end{forest}
\end{exe}
This will make it parallel to the structure of small clause complements---like that in~(\ref{ex5:kaynesmcl}).
\begin{exe}
	\raggedright
	\ex \label{ex5:kaynesmcl} 
	\small
	\begin{forest}
	[\ibar{V}
		[\obar{V} [$\surd$believe] ]
		[AP 
			[DP [this fact, roof] ]
			[\ibar{A} [\obar{A} [interesting]] ]
		]
	]
	\end{forest}
\end{exe}
In both structures, there is a clause-like phrase that separates the DP in the Specifier of that phrase and the term outside the phrase that is responsible for providing that DP with Case. This clause-like phrase is, by hypothesis, a barrier for Case assignment when it's the complement to a noun, and these structures are then all violations of the Case filter.

Interestingly the other complement structure that these verbs support is not similarly incompatible with nominals.
\begin{exe}
	\raggedright
	\ex \label{ex5:kaynedoc} 
	\begin{xlist}
		\ex John's presentation of a medal to Mary
		\ex his offer of a book to Mary
		\ex his recommendation of Mary to the Board
		\ex their surrender of the city to the enemy
		\ex her explanation of the problem to the students
	\end{xlist}
	\hspace*{\fill} \citep[(91):~171]{Kayne:1981b} \hspace*{.5in}
\end{exe}
If Kayne's argument is correct, this suggests that the DP+PP frame doesn't make use of a small clause structure. This is inconsistent with Harley's proposal about how to relate the two frames, since her account requires that the feature which determines whether a verb participates in the alternation is its ability to c-selecte PP small clauses in both frames (see~(\ref{ex5:harleya})). As far as I know, these conflicting pieces of evidence have not been reconciled.

There are other problems with Kayne's proposal. \citet{Borer:2013} discovered that there are certain kinds of small clauses that can be found in nominals, casting doubt on Kayne's method. There is evidence that so-called resultatives have a small clause structure, as we briefly touched upon when considering~(\ref{ex5:intransitiveres}).
\pagebreak
\begin{exe}
	\raggedright
	\exr{ex5:intransitiveres} Suzanne danced herself silly.
	\sn
	\small
	\begin{forest}
	[$\lambda \epsilon$ $\exists \epsilon'$ \textsc{dance}($\epsilon$) $\wedge$ \textsc{silly}($\epsilon'$) $\wedge$ \textsc{holder}(herself)($\epsilon'$) $\wedge$ \textsc{cause}($\epsilon'$)($\epsilon$) \\\ibar{V}
		[$\lambda \epsilon$ \textsc{dance}($\epsilon$)\\ \obar{V} [$\surd$dance] ]
		[$\lambda Q$ $\lambda \epsilon$ $\exists \epsilon'$ $Q$($\epsilon$) $\wedge$ \textsc{silly}($\epsilon'$) $\wedge$ \textsc{holder}(herself)($\epsilon'$) $\wedge$ \textsc{cause}($\epsilon'$)($\epsilon$)\\ VP [\ibar{V}
			[$\lambda P$ $\lambda Q$ $\lambda \epsilon$ $\exists \epsilon'$ $Q$($\epsilon$) $\wedge$ $P$($\epsilon'$)\\ $\wedge$ \textsc{cause}($\epsilon'$)($\epsilon$) \\ \obar{V} [\shade{case}] ]
			[$\lambda \epsilon$ \textsc{silly}($\epsilon$) $\wedge$ \textsc{holder}(herself)($\epsilon$)\\ AP
				[DP [herself, roof] ]
				[\ibar{A} [silly, roof] ]
			 ]
		] ]
	]
	\end{forest}
\end{exe}
But Borer points out that these constructions can survive in nominalizations.
\begin{exe}
	\raggedright
	\ex \label{ex5:borerres} 
	\begin{xlist}
		\ex Kim's hammering of the metal flat
		\ex the hammering of the metal flat (by Kim)
		\ex your burning of the kitchen walls black
		\ex the burning of the kitchen walls black (by you)
	\end{xlist}
	\hspace*{\fill} \citep[(66):~108]{Borer:2013} \hspace*{.5in}
\end{exe}
Borer's own account of the absence of the double object frame in nominals is that they do not have the resources for assigning two object Cases in English, and this is what distinguishes nominals from verbs. \citet{Pesetsky:1995} also has an account for the absence of the double object frame in nominals that divorces it from the fact that nominals don't permit many small clauses, raising and exceptional Case marking infinitives. His proposal targets the hidden morphemes that the verb expones in the double object frame. He suggests that these interact fatally with the morphology that creates nouns from verbs. If (\ref{ex5:intransitiveres}) is correct, then the grammaticality of (\ref{ex5:borerres}) is also problematic for Pesetsky's proposal. I do not know of an account for the full gamut of effects described here about how the argument structure of nouns differs from the argument structure of the verbs they are derived from.